\documentclass[11pt]{article}
\usepackage[utf8]{inputenc}
\usepackage[T1]{fontenc}
\usepackage[spanish,es-nodecimaldot,es-lcroman]{babel}
\usepackage{amsmath,amssymb,amsthm,mathtools}
\usepackage[margin=2.5cm]{geometry}
\usepackage{booktabs,tabularx,array}
\usepackage{enumitem}
\usepackage[hidelinks]{hyperref}
\setlength{\emergencystretch}{2em}
\newcolumntype{Y}{>{\raggedright\arraybackslash}X}

\newtheorem{teorema}{Teorema}[section]
\newtheorem{lema}[teorema]{Lema}
\newtheorem{proposicion}[teorema]{Proposición}
\newtheorem{corolario}[teorema]{Corolario}
\theoremstyle{definition}
\newtheorem{definicion}[teorema]{Definición}
\newtheorem{ejemplo}[teorema]{Ejemplo}
\newtheorem{ejercicio}[teorema]{Ejercicio}
\theoremstyle{remark}
\newtheorem{observacion}[teorema]{Observación}

\newcommand{\R}{\mathbb{R}}
\newcommand{\N}{\mathbb{N}}
\newcommand{\E}{\mathbb{E}}
\newcommand{\Prob}{\mathbb{P}}
\DeclareMathOperator{\Var}{Var}
\DeclareMathOperator{\Cov}{Cov}
\DeclareMathOperator{\Corr}{Corr}
\DeclareMathOperator*{\argmin}{arg\,min}
\DeclareMathOperator*{\argmax}{arg\,max}
\DeclareMathOperator{\med}{med}
\DeclareMathOperator{\IQR}{IQR}
\DeclareMathOperator{\sgn}{sgn}
\DeclareMathOperator{\rango}{rango}
\DeclareMathOperator{\idf}{idf}
\DeclareMathOperator{\tf}{tf}
\DeclareMathOperator{\PMI}{PMI}
\newcommand{\ind}[1]{\mathbf{1}\!\left[#1\right]}
\newcommand{\BC}{\mathrm{BC}}
\newcommand{\YJ}{\mathrm{YJ}}
\newcommand{\norm}[1]{\left\lVert #1\right\rVert}
\newcommand{\abs}[1]{\left\lvert #1\right\rvert}
\newcommand{\ip}[2]{\left\langle #1,#2\right\rangle}
\newcommand{\Dtr}{\mathcal{D}_{\mathrm{tr}}}
\newcommand{\Dte}{\mathcal{D}_{\mathrm{te}}}

\title{Ingeniería de características:\\ desarrollo matemático, contraste de técnicas\\ y servicios de AWS donde se aplican}
\author{Desarrollo del Capítulo 3 de la guía oficial MLA-C01}
\date{Septiembre de 2026}

\begin{document}
\maketitle
\tableofcontents

\section{Marco general}\label{sec:marco}

\noindent Este texto desarrolla con rigor las técnicas de transformación de datos e ingeniería de características que presenta el Capítulo~3 de la guía de estudio del examen AWS Certified Machine Learning Engineer -- Associate: limpieza (valores faltantes, atípicos, duplicados), selección y extracción de características, escalado, transformaciones no lineales, discretización, codificación de variables categóricas, características de series de tiempo, de imágenes y de texto, tratamiento del desbalance de clases y división de datos. Para cada técnica se da su formulación matemática, se demuestran las propiedades que la guía afirma (y algunas que da por sentadas), se señalan sus límites mediante contraejemplos y se contrasta con las técnicas que compiten con ella. La última sección enumera los servicios de AWS en los que se aplica cada técnica, con el mismo reparto que hace la guía.

\medskip

\noindent El nivel supuesto es el de un lector con formación en probabilidad y estadística matemática (esperanza, varianza, cuantiles, estimación por máxima verosimilitud, regresión lineal por mínimos cuadrados) y álgebra lineal (valores propios, rango, proyecciones). No se usa teoría de la medida.

\subsection{Datos, características y transformaciones ajustadas}

\noindent El problema que motiva toda la disciplina es el siguiente. Un algoritmo de aprendizaje supervisado recibe pares $(x_i,y_i)$ con $x_i\in\R^p$ y busca una función $f$ tal que $f(x_i)$ aproxime a $y_i$. Pero los datos crudos rara vez son vectores de $\R^p$ adecuados: contienen cadenas de texto, fechas, imágenes, huecos, escalas incomparables y errores. La ingeniería de características es la construcción de una aplicación que lleva cada registro crudo a un vector numérico con el que el algoritmo aprende mejor. Conviene fijar la notación con la que se hablará de ella.

\begin{definicion}[Registro, característica, matriz de diseño]\label{def:diseno}
\noindent Sea $\mathcal{Z}$ el conjunto de registros crudos posibles (filas de una tabla con columnas de tipos mixtos, imágenes, documentos). Un \emph{conjunto de datos} es una sucesión finita $\mathcal{D}=((z_1,y_1),\dots,(z_n,y_n))$ con $z_i\in\mathcal{Z}$ y $y_i$ en un conjunto de etiquetas $\mathcal{Y}$ (por ejemplo $\mathcal{Y}=\R$ en regresión o $\mathcal{Y}=\{0,1\}$ en clasificación binaria). Una \emph{característica} es una función $\phi_j:\mathcal{Z}\to\R$; un \emph{mapa de características} es $\phi=(\phi_1,\dots,\phi_p):\mathcal{Z}\to\R^p$. La \emph{matriz de diseño} es $X\in\R^{n\times p}$ con entradas $X_{ij}=\phi_j(z_i)$; su fila $i$ se denota $x_i^\top$ y su columna $j$, $x_{\cdot j}$.
\end{definicion}

\noindent En vocabulario estadístico, las características son las covariables o predictores y $y$ es la variable respuesta. La definición separa deliberadamente el registro $z$ de su representación $\phi(z)$: la ingeniería de características consiste precisamente en elegir $\phi$. Dos ejemplos del texto original muestran el alcance. Si $z$ es una fila con una columna de fecha, la característica ``día de la semana'' es $\phi(z)\in\{0,\dots,6\}$, que no está escrita en la tabla pero se calcula de ella. Si $z$ es una imagen, $\phi(z)$ puede ser la intensidad de cada píxel (con $p$ del orden de $10^5$) o un vector corto calculado por una red neuronal (Sección~\ref{sec:imagenes}).

\medskip

\noindent Casi todas las transformaciones de este texto dependen de parámetros que se estiman con los datos: la media y la desviación estándar de un escalador, los cuartiles del método IQR, el vocabulario de un codificador, la $\lambda$ de Box--Cox, los ejes de un PCA. Esta dependencia es la fuente del error más grave de la ingeniería de características, la fuga de datos, y por eso se formaliza desde el principio.

\begin{definicion}[Transformación ajustada]\label{def:ajustada}
\noindent Una \emph{transformación ajustada} es un par $(\hat\theta,T)$ donde $\hat\theta$ es una función que asigna a cada conjunto de datos $\mathcal{D}$ un parámetro $\hat\theta(\mathcal{D})$ en un espacio $\Theta$ (la etapa \texttt{fit}), y $T:\Theta\times\R^p\to\R^{q}$ aplica la transformación a un vector con un parámetro dado (la etapa \texttt{transform}). Escribimos $T_\theta(x)=T(\theta,x)$.
\end{definicion}

\noindent Por ejemplo, la estandarización corresponde a $\hat\theta(\mathcal{D})=(\bar x,s)$ (media y desviación estándar muestrales de la columna) y $T_{(\mu,\sigma)}(x)=(x-\mu)/\sigma$. La separación entre $\hat\theta$ y $T$ es exactamente la separación entre los métodos \texttt{fit} y \texttt{transform} de scikit-learn que describe la guía.

\medskip

\noindent El uso correcto de una transformación ajustada se deduce de cómo se usará el modelo. En producción, el modelo recibe un registro nuevo $z_{\mathrm{new}}$ que no existía cuando se entrenó; la transformación que se le aplica solo puede depender de los datos de entrenamiento. Si la evaluación ha de imitar la producción, lo mismo debe valer para los datos de prueba.

\begin{definicion}[Fuga de datos por preprocesamiento]\label{def:fuga}
\noindent Sea $\mathcal{D}=\Dtr\cup\Dte$ una partición en entrenamiento y prueba. Se dice que un procedimiento de evaluación tiene \emph{fuga de datos por preprocesamiento} si la transformación aplicada a los registros de $\Dte$ es $T_{\hat\theta(\mathcal{D}')}$ con $\mathcal{D}'\cap\Dte\neq\emptyset$, es decir, si el parámetro ajustado usó algún registro de prueba. El procedimiento correcto usa $T_{\hat\theta(\Dtr)}$ tanto para $\Dtr$ como para $\Dte$.
\end{definicion}

\noindent La condición es precisa: no basta con que el modelo no ``vea'' los datos de prueba, tampoco deben verlos los parámetros de la transformación. El ejemplo siguiente muestra que la fuga cambia efectivamente lo que se evalúa.

\begin{ejemplo}[La media ajustada contiene información de prueba]\label{ej:fuga-media}
\noindent Sea $\Dtr=\{x_1,\dots,x_m\}$ y $\Dte=\{x_{m+1},\dots,x_n\}$ en una sola columna, y sea $\bar x_{\mathrm{tr}}=\frac1m\sum_{i\le m}x_i$, $\bar x_{\mathrm{te}}=\frac1{n-m}\sum_{i>m}x_i$. La media de todos los datos es
\begin{equation}\label{eq:media-total}
\bar x=\frac{m}{n}\,\bar x_{\mathrm{tr}}+\frac{n-m}{n}\,\bar x_{\mathrm{te}},
\end{equation}
lo que se obtiene separando la suma $\sum_{i=1}^n x_i=m\bar x_{\mathrm{tr}}+(n-m)\bar x_{\mathrm{te}}$ y dividiendo entre $n$. Si se centra con $\bar x$, el valor transformado de un punto de entrenamiento es $x_i-\bar x$, que por \eqref{eq:media-total} depende de $\bar x_{\mathrm{te}}$: al variar un solo dato de prueba en $\delta$, cada dato de entrenamiento transformado cambia en $-\delta/n$. El modelo entrenado sobre esos valores es, por tanto, función de los datos de prueba, y la evaluación deja de medir el desempeño sobre datos no vistos. En este ejemplo el efecto es pequeño (orden $1/n$); en transformaciones con muchos parámetros, como el vocabulario de un codificador o el sobremuestreo sintético (Sección~\ref{sec:desbalance}), puede ser grande.
\end{ejemplo}

\noindent La regla práctica que la guía enuncia sin justificar, ``normaliza o estandariza siempre después de dividir'', es la Definición~\ref{def:fuga} aplicada a los escaladores. La misma regla rige para la media de imputación, los límites del IQR, el vocabulario de un codificador, las componentes de un PCA y el sobremuestreo. A lo largo del texto se señalará qué parámetros ajusta cada técnica.

\section{Limpieza de datos}\label{sec:limpieza}

\subsection{Valores faltantes}\label{sec:faltantes}

\noindent \emph{Motivación.} Casi todos los algoritmos clásicos (regresión lineal y logística, máquinas de vectores de soporte, $k$ vecinos) están definidos sobre vectores de $\R^p$ completos: la expresión $\beta^\top x$ no tiene sentido si falta una coordenada. Descartar las filas incompletas es la solución más simple, pero con $p$ columnas en las que cada entrada falta de manera independiente con probabilidad $r$, la probabilidad de que una fila esté completa es $(1-r)^p$; con $r=0.02$ y $p=50$ esto da $0.98^{50}\approx 0.36$, así que se perdería cerca de dos tercios del conjunto de datos por huecos que afectan al $2\%$ de las celdas. Esta pérdida motivó la imputación. El análisis de cuándo la imputación es inocua requiere precisar \emph{por qué} falta el dato, que es lo que formalizó Rubin (1976).

\begin{definicion}[Mecanismos de falta]\label{def:mcar}
\noindent Sea $X=(X_{\mathrm{obs}},X_{\mathrm{mis}})$ el vector completo de un registro, separado en la parte observada y la faltante, y sea $R$ el vector de indicadores de respuesta ($R_j=1$ si la coordenada $j$ se observó). El mecanismo de falta es:
\begin{enumerate}[label=(\alph*)]
\item \emph{completamente al azar} (MCAR) si $R$ es independiente de $X$;
\item \emph{al azar} (MAR) si la distribución de $R$ condicionada a $X$ depende solo de $X_{\mathrm{obs}}$;
\item \emph{no al azar} (MNAR) en cualquier otro caso, es decir, cuando la probabilidad de faltar depende del propio valor no observado.
\end{enumerate}
\end{definicion}

\noindent La guía describe el caso de ``pocos valores faltantes distribuidos aleatoriamente'' por una falla de ingesta; eso es MCAR. Un ejemplo de MNAR es el ingreso declarado, cuando los ingresos altos se declaran con menos frecuencia.

\begin{proposicion}[La media observada bajo MCAR y MNAR]\label{prop:mcar}
\noindent Sean $X_1,\dots,X_n$ i.i.d.\ con media $\mu$ finita y $R_1,\dots,R_n$ sus indicadores de respuesta, y sea $M=\sum_i R_i$.
\begin{enumerate}[label=(\roman*)]
\item Si el mecanismo es MCAR, con $(R_1,\dots,R_n)$ independiente de $(X_1,\dots,X_n)$, entonces $\E[\bar X_{\mathrm{obs}}\mid R_1,\dots,R_n]=\mu$ en el suceso $M\ge1$, donde $\bar X_{\mathrm{obs}}=M^{-1}\sum_i R_iX_i$.
\item Si $X_i\sim\mathrm{Exp}(1)$ y cada valor se observa si y solo si $X_i\le c$ (MNAR), entonces $\E[X_i\mid R_i=1]=1-\dfrac{c\,e^{-c}}{1-e^{-c}}<1=\mu$.
\end{enumerate}
\end{proposicion}
\begin{proof}
\noindent (i) Condicionando a $R=(R_1,\dots,R_n)$ con $M\ge1$, $M$ y los $R_i$ son constantes, así que por linealidad $\E[\bar X_{\mathrm{obs}}\mid R]=M^{-1}\sum_iR_i\,\E[X_i\mid R]$. Por la independencia entre $R$ y $X$, la distribución condicional de $X_i$ dado $R$ es la marginal, de modo que $\E[X_i\mid R]=\mu$. Por tanto $\E[\bar X_{\mathrm{obs}}\mid R]=M^{-1}\mu\sum_iR_i=\mu$.

\medskip

\noindent (ii) El suceso $R_i=1$ es $\{X_i\le c\}$, con probabilidad $\int_0^c e^{-x}\,dx=1-e^{-c}$. Integrando por partes con $u=x$, $dv=e^{-x}dx$,
\[
\int_0^c x e^{-x}\,dx=\big[-xe^{-x}\big]_0^c+\int_0^c e^{-x}\,dx=-ce^{-c}+1-e^{-c}.
\]
Dividiendo entre $\Prob(X_i\le c)$ se obtiene $\E[X_i\mid X_i\le c]=\dfrac{1-e^{-c}-ce^{-c}}{1-e^{-c}}=1-\dfrac{ce^{-c}}{1-e^{-c}}$, y el sustraendo es positivo para $c>0$.
\end{proof}

\noindent Con $c=2$ la media condicionada vale $1-2e^{-2}/(1-e^{-2})\approx 0.687$: imputar con la media observada subestima la media real en un $31\%$, por grande que sea la muestra. La imputación simple, por tanto, solo es defendible bajo MCAR (o bajo MAR si se imputa condicionando a lo observado). La guía también distingue entre imputar con la media, la mediana o la moda. Las tres son el mejor valor constante bajo pérdidas distintas.

\begin{proposicion}[Media, mediana y moda como constantes óptimas]\label{prop:constantes}
\noindent Sean $x_1,\dots,x_m$ los valores observados de una característica.
\begin{enumerate}[label=(\roman*)]
\item $\sum_i(x_i-c)^2$ se minimiza en $c=\bar x$, y solo ahí.
\item $\sum_i\abs{x_i-c}$ se minimiza en cualquier mediana muestral.
\item Si los valores pertenecen a un conjunto finito, $\sum_i\ind{x_i\neq c}$ se minimiza en cualquier moda.
\end{enumerate}
\end{proposicion}
\begin{proof}
\noindent (i) Sumando y restando $\bar x$, $\sum_i(x_i-c)^2=\sum_i(x_i-\bar x)^2+2(\bar x-c)\sum_i(x_i-\bar x)+m(\bar x-c)^2$. El término cruzado es nulo porque $\sum_i(x_i-\bar x)=\sum_ix_i-m\bar x=0$. Queda $\sum_i(x_i-\bar x)^2+m(\bar x-c)^2$, que es mínimo si y solo si $c=\bar x$.

\medskip

\noindent (ii) La función $g(c)=\sum_i\abs{x_i-c}$ es continua y lineal a trozos. En un punto $c$ distinto de todos los $x_i$, su derivada es $\#\{i:x_i<c\}-\#\{i:x_i>c\}$, porque cada término $\abs{x_i-c}$ tiene derivada $+1$ si $x_i<c$ y $-1$ si $x_i>c$. Esa derivada es $\le0$ mientras a lo sumo la mitad de los datos quede a la izquierda de $c$ y $\ge0$ en cuanto a lo sumo la mitad quede a la derecha. Luego $g$ es no creciente a la izquierda de cualquier mediana y no decreciente a su derecha, y el mínimo se alcanza en las medianas.

\medskip

\noindent (iii) $\sum_i\ind{x_i\ne c}=m-\#\{i:x_i=c\}$, que es mínimo cuando la frecuencia de $c$ es máxima, es decir, en una moda.
\end{proof}

\noindent La proposición justifica la regla de la guía: la media es natural cuando interesa el error cuadrático y la distribución es simétrica (en cuyo caso media y mediana coinciden en la población); con asimetría o atípicos, la mediana es preferible porque, como se verá en la Proposición~\ref{prop:ruptura}, no la arrastra un solo valor extremo; para una característica categórica solo tiene sentido la moda, porque no hay suma ni orden. La imputación tiene además un costo sobre la dispersión que la guía no menciona.

\begin{proposicion}[La imputación por la media reduce la varianza]\label{prop:imput-var}
\noindent Sea una columna con $n$ registros, de los cuales $m$ están observados ($2\le m<n$), con media $\bar x_{\mathrm{obs}}$ y varianza muestral $s^2_{\mathrm{obs}}=\frac{1}{m-1}\sum_{\mathrm{obs}}(x_i-\bar x_{\mathrm{obs}})^2$. Si los $n-m$ huecos se rellenan con $\bar x_{\mathrm{obs}}$, la columna completada tiene media $\bar x_{\mathrm{obs}}$ y varianza
\[
s^2_{\mathrm{imp}}=\frac{m-1}{n-1}\,s^2_{\mathrm{obs}}<s^2_{\mathrm{obs}}\quad\text{si }s^2_{\mathrm{obs}}>0.
\]
\end{proposicion}
\begin{proof}
\noindent La suma de la columna completada es $m\bar x_{\mathrm{obs}}+(n-m)\bar x_{\mathrm{obs}}=n\bar x_{\mathrm{obs}}$, así que su media es $\bar x_{\mathrm{obs}}$. En la suma de cuadrados respecto de esa media, cada valor imputado contribuye $(\bar x_{\mathrm{obs}}-\bar x_{\mathrm{obs}})^2=0$, por lo que la suma total es $\sum_{\mathrm{obs}}(x_i-\bar x_{\mathrm{obs}})^2=(m-1)s^2_{\mathrm{obs}}$. Dividiendo entre $n-1$ se obtiene la fórmula, y $\frac{m-1}{n-1}<1$ porque $m<n$.
\end{proof}

\noindent Con un $30\%$ de huecos ($m=0.7n$) la varianza se reduce aproximadamente al $70\%$. Un modelo que use esa columna ve una característica artificialmente concentrada, y los valores imputados forman una masa puntual en $\bar x_{\mathrm{obs}}$ que no existe en la población. Por eso la guía sugiere acompañar la imputación con un indicador $x^{\mathrm{mis}}=\ind{\text{valor faltante}}$: este indicador permite al modelo tratar distinto a los registros imputados y, en situación MNAR, recoge la información que el hecho de faltar contiene sobre la respuesta. La media y la mediana de imputación son parámetros ajustados en el sentido de la Definición~\ref{def:ajustada} y deben calcularse solo con $\Dtr$.

\medskip

\noindent \emph{Contraste de las tres estrategias.} Eliminar filas conserva la distribución de las filas completas pero pierde datos (el factor $(1-r)^p$) y sesga bajo MNAR; eliminar una columna solo es razonable si concentra los huecos; imputar conserva todas las filas pero reduce la varianza (Proposición~\ref{prop:imput-var}) y sesga bajo MNAR (Proposición~\ref{prop:mcar}). Recolectar de nuevo es la única opción que corrige un mecanismo MNAR, y su costo es el que la guía pide sopesar. Algunas implementaciones de \emph{gradient boosting} evitan la imputación aprendiendo, en cada corte, hacia qué rama enviar los faltantes.

\subsection{Valores atípicos}\label{sec:atipicos}

\noindent \emph{Motivación.} Desde Gauss, el modelo de referencia para los errores de medición es la distribución normal, y bajo ese modelo un valor a más de tres desviaciones estándar de la media es muy improbable. De ahí proviene la regla de las tres sigmas que cita la guía. En el siglo XX se observó que esta regla falla precisamente cuando más se la necesita: los atípicos inflan la media y la desviación estándar con las que se los busca. Tukey (1977) propuso reglas basadas en cuartiles, y la teoría de la robustez de Huber y Hampel dio una medida precisa de cuánto puede un estimador resistir la contaminación. Esta subsección reconstruye ese desarrollo.

\begin{definicion}[Puntuación Z]\label{def:z}
\noindent Dada una muestra $x_1,\dots,x_n$ con media $\bar x$ y desviación estándar $s>0$, la \emph{puntuación Z} de $x_i$ es $z_i=(x_i-\bar x)/s$. La \emph{regla de las tres sigmas} declara atípico a $x_i$ si $\abs{z_i}>3$.
\end{definicion}

\noindent Hay dos convenciones para $s$: la \emph{poblacional} $s_0^2=\frac1n\sum_i(x_i-\bar x)^2$ (la que usa \texttt{scipy.stats.zscore} por defecto) y la \emph{muestral} $s_1^2=\frac1{n-1}\sum_i(x_i-\bar x)^2$. La justificación probabilística de la regla, bajo normalidad, es que $\Prob(\abs{Z}\le3)=0.9973$ para $Z\sim N(0,1)$, así que solo un $0.27\%$ de las observaciones de una normal quedaría marcada (la guía redondea a $0.3\%$). Sin normalidad la cota disponible es mucho más débil.

\begin{proposicion}[Desigualdad de Chebyshev]\label{prop:chebyshev}
\noindent Si $X$ tiene media $\mu$ y varianza $\sigma^2\in(0,\infty)$, entonces para todo $k>0$, $\Prob(\abs{X-\mu}\ge k\sigma)\le 1/k^2$.
\end{proposicion}
\begin{proof}
\noindent Sea $A=\{\abs{X-\mu}\ge k\sigma\}$. En $A$ se cumple $(X-\mu)^2\ge k^2\sigma^2$, y fuera de $A$ se cumple $(X-\mu)^2\ge0$. Por tanto $(X-\mu)^2\ge k^2\sigma^2\,\ind{A}$ en todo punto. Tomando esperanzas, $\sigma^2\ge k^2\sigma^2\Prob(A)$, y dividiendo entre $k^2\sigma^2$ se obtiene la cota.
\end{proof}

\noindent Para $k=3$ la cota es $1/9\approx 11\%$, cuarenta veces el $0.27\%$ normal. La cota se alcanza: la variable que vale $\pm3$ con probabilidad $1/18$ cada uno y $0$ con probabilidad $8/9$ tiene media $0$, varianza $2\cdot 9/18=1$ y $\Prob(\abs{X}\ge3)=1/9$. Así, sin un supuesto de forma, ``tres desviaciones'' no fija ninguna proporción de atípicos. La guía muestra además que, en su ejemplo de diez puntos, la regla no detecta el punto $[25,23]$. Hay una razón puramente aritmética.

\begin{proposicion}[Cota de Samuelson--Shiffler]\label{prop:shiffler}
\noindent Sea $n\ge2$ y $x_1,\dots,x_n$ con $s>0$. Entonces, para todo $i$,
\[
\abs{z_i}\le\sqrt{n-1}\quad\text{si }s=s_0,\qquad\qquad \abs{z_i}\le\frac{n-1}{\sqrt n}\quad\text{si }s=s_1.
\]
\end{proposicion}
\begin{proof}
\noindent Sin pérdida de generalidad $i=1$. Sea $d_j=x_j-\bar x$. Como $\sum_jd_j=0$, se tiene $d_1=-\sum_{j\ge2}d_j$. Por la desigualdad de Cauchy--Schwarz aplicada a los vectores $(1,\dots,1)$ y $(d_2,\dots,d_n)$ de $\R^{n-1}$,
\[
d_1^2=\Big(\sum_{j\ge2}d_j\Big)^2\le (n-1)\sum_{j\ge2}d_j^2.
\]
Entonces $\sum_{j=1}^nd_j^2=d_1^2+\sum_{j\ge2}d_j^2\ge d_1^2+\frac{d_1^2}{n-1}=\frac{n}{n-1}\,d_1^2$. Con $s=s_0$ el miembro izquierdo es $ns_0^2$, de donde $d_1^2\le(n-1)s_0^2$ y $\abs{z_1}\le\sqrt{n-1}$. Con $s=s_1$ el miembro izquierdo es $(n-1)s_1^2$, de donde $d_1^2\le\frac{(n-1)^2}{n}s_1^2$ y $\abs{z_1}\le(n-1)/\sqrt n$.
\end{proof}

\noindent La igualdad en Cauchy--Schwarz exige $d_2=\dots=d_n$, es decir, que todos los demás puntos coincidan. Con $n=10$ y la convención poblacional la cota es $\sqrt9=3$: ningún punto puede tener $\abs{z}>3$, cualquiera que sea su valor, y el valor $3$ solo se alcanza si los otros nueve son idénticos. (Con la convención muestral la cota es $9/\sqrt{10}\approx2.85$.) La regla de las tres sigmas es, por tanto, vacía para $n\le10$. La segunda razón de su fallo es el \emph{enmascaramiento}.

\begin{ejemplo}[Enmascaramiento en el ejemplo de la guía]\label{ej:enmascaramiento}
\noindent Para la segunda característica del ejemplo de la guía, $y=(11,12,12.5,8,8.5,7.5,5,4.7,6,23)$, se tiene $\bar y=9.82$ y $s_0=5.114$, así que el punto $23$ tiene $z=(23-9.82)/5.114=2.58$. Si en cambio se calculan la media y la desviación estándar muestral de los otros nueve puntos ($\bar y_{(-10)}=8.356$, $s_{1,(-10)}=2.926$), la distancia estandarizada de $23$ es $(23-8.356)/2.926\approx5.0$. El atípico infla la $s$ con la que se lo mide y se oculta: la suma de cuadrados pasa de $8\cdot2.926^2\approx68.5$ sin él a $10\cdot5.114^2\approx261.5$ con él, y casi tres cuartas partes de la dispersión total provienen de ese único punto.
\end{ejemplo}

\noindent La herramienta que mide esta fragilidad es el punto de ruptura.

\begin{definicion}[Punto de ruptura finito]\label{def:ruptura}
\noindent Sea $T$ un estimador que asigna un número real a cada muestra de tamaño $n$. El \emph{punto de ruptura finito} de $T$ en la muestra $x$ es la menor fracción $k/n$ tal que, reemplazando $k$ de los $n$ valores por valores arbitrarios, $\abs{T}$ puede hacerse tan grande como se quiera.
\end{definicion}

\noindent La definición captura la propiedad deseada (``unos pocos datos corruptos no deben poder llevar el estimador a cualquier parte'') de forma cuantitativa, y permite comparar estimadores de posición y de escala.

\begin{proposicion}[Punto de ruptura de media, mediana y cuartiles]\label{prop:ruptura}
\noindent Para una muestra de tamaño $n$: (i) la media y la desviación estándar tienen punto de ruptura $1/n$; (ii) la mediana tiene punto de ruptura $\lceil n/2\rceil/n$, que tiende a $1/2$; (iii) los cuartiles $Q_1,Q_3$ y el $\IQR=Q_3-Q_1$ tienen punto de ruptura del orden de $1/4$.
\end{proposicion}
\begin{proof}
\noindent (i) Reemplazando $x_1$ por $t$, la media es $(t+\sum_{j\ge2}x_j)/n$, que tiende a $\infty$ con $t$; un solo valor basta. Para la desviación estándar, $ns_0^2\ge(t-\bar x)^2$ y $t-\bar x=\frac{n-1}{n}t-\frac1n\sum_{j\ge2}x_j\to\infty$.

\medskip

\noindent (ii) Si se reemplazan $k<n/2$ valores, más de la mitad de los datos siguen siendo originales, luego la mediana queda entre el mínimo y el máximo de los datos originales y está acotada. Si se reemplazan $k\ge\lceil n/2\rceil$ valores por un mismo $t$ grande, al menos la mitad de la muestra vale $t$ y la mediana es $\ge t$ (o el promedio de $t$ con un valor acotado, si $n$ es par y $k=n/2$), que tiende a $\infty$.

\medskip

\noindent (iii) El mismo argumento, aplicado al cuantil de orden $1/4$ o $3/4$: mientras menos de aproximadamente un cuarto de los datos se reemplace, en cada cola quedan datos originales suficientes para que el cuantil quede entre valores originales; al reemplazar un cuarto o más en la cola correspondiente, el cuantil puede llevarse a infinito, y con él el $\IQR$.
\end{proof}

\noindent La regla de Tukey sustituye la media y la desviación estándar por cuartiles.

\begin{definicion}[Regla del rango intercuartílico]\label{def:iqr}
\noindent Sean $Q_1$ y $Q_3$ los cuartiles muestrales y $\IQR=Q_3-Q_1$. Un valor $x$ es atípico según la \emph{regla del IQR} (o vallas de Tukey) si $x<Q_1-1.5\,\IQR$ o $x>Q_3+1.5\,\IQR$.
\end{definicion}

\noindent La constante $1.5$ no es arbitraria: fija qué fracción de una normal se marcaría, que resulta comparable a la de la regla de las tres sigmas.

\begin{proposicion}[Vallas de Tukey bajo normalidad]\label{prop:tukey-normal}
\noindent Si $X\sim N(\mu,\sigma^2)$, las vallas poblacionales son $\mu\pm 4z_{0.75}\,\sigma\approx\mu\pm2.698\,\sigma$, donde $z_{0.75}\approx0.6745$ es el cuantil $0.75$ de la normal estándar, y la probabilidad de caer fuera de ellas es aproximadamente $0.70\%$.
\end{proposicion}
\begin{proof}
\noindent Por simetría de la normal, $Q_1=\mu-z_{0.75}\sigma$ y $Q_3=\mu+z_{0.75}\sigma$, luego $\IQR=2z_{0.75}\sigma$. La valla superior es $Q_3+1.5\IQR=\mu+z_{0.75}\sigma+3z_{0.75}\sigma=\mu+4z_{0.75}\sigma$, y la inferior es simétrica. Con $z_{0.75}=0.6745$ resulta $4z_{0.75}=2.698$, y $\Prob(\abs{Z}>2.698)=2(1-\Phi(2.698))\approx0.0070$.
\end{proof}

\noindent Así, bajo normalidad las dos reglas son comparables ($0.27\%$ frente a $0.70\%$), pero la del IQR hereda el punto de ruptura de los cuartiles (Proposición~\ref{prop:ruptura}) en lugar del de la media.

\begin{ejemplo}[Regla del IQR en el ejemplo de la guía]\label{ej:iqr}
\noindent Con la interpolación lineal de pandas, para $y$ se obtiene $Q_1=6.375$, $Q_3=11.75$, $\IQR=5.375$ y valla superior $11.75+1.5\cdot5.375=19.8125$. El valor $23$ la supera y se marca. Para $x=(4,3.8,4.5,8,9,9.5,13,14,13.7,25)$ se obtiene $Q_1=5.375$, $Q_3=13.525$, $\IQR=8.15$ y valla superior $25.75$, así que $x=25$ no se marca. La fila $[25,23]$ se marca porque el código de la guía usa \texttt{.any(axis=1)}: basta que una coordenada caiga fuera.
\end{ejemplo}

\noindent \emph{Contraste Z frente a IQR.} La puntuación Z usa estimadores con punto de ruptura $1/n$, está acotada por $\sqrt{n-1}$ (Proposición~\ref{prop:shiffler}), sufre enmascaramiento (Ejemplo~\ref{ej:enmascaramiento}) y su calibración ($0.27\%$) solo vale bajo normalidad. La regla del IQR usa estimadores con punto de ruptura cercano a $1/4$, no depende de la media y es aplicable con asimetría; su desventaja es que, bajo normalidad y con muestras grandes, marca más puntos legítimos ($0.70\%$) y que, en distribuciones de colas pesadas (lognormal, Pareto), marca como atípicos valores que son parte natural de la cola. Ambas son reglas univariantes: un punto puede ser atípico solo en la combinación de dos coordenadas y pasar ambas reglas. Tanto la media y $s$ como los cuartiles son parámetros ajustados y deben calcularse con $\Dtr$.

\medskip

\noindent \emph{Tratamiento.} La guía enumera tres tratamientos: eliminar, imputar y transformar con logaritmo. Eliminar es correcto si el atípico es un error; si es un valor natural (una compra legítima de 150\,000 en un modelo de fraude) eliminarlo quita justamente la información más valiosa. Imputar hereda los problemas de la Proposición~\ref{prop:imput-var}. La transformación logarítmica se aplica a toda la columna, no solo al atípico (si se aplicara solo al punto, la columna mezclaría dos escalas), y su efecto se estudia en la Sección~\ref{sec:nolineales}: en el ejemplo de la guía la puntuación Z de $23$ pasa de $2.58$ a $2.11$ tras aplicar $\log_{10}$ a la columna, lo que se verifica calculando la media ($0.9436$) y la desviación estándar poblacional ($0.1984$) de los logaritmos: $(\log_{10}23-0.9436)/0.1984=(1.3617-0.9436)/0.1984\approx2.11$.

\subsection{Deduplicación}\label{sec:dedup}

\noindent La guía afirma que los duplicados distorsionan resultados, sesgan el entrenamiento e inflan las métricas. Las tres afirmaciones se siguen de una observación sobre el riesgo empírico.

\begin{proposicion}[Un duplicado es un peso]\label{prop:duplicado}
\noindent Sea $\hat R(f)=\frac1n\sum_{i=1}^n\ell(f(x_i),y_i)$ el riesgo empírico de una función $f$ bajo una pérdida $\ell$. Si el registro $(x_1,y_1)$ aparece $k$ veces en lugar de una, el riesgo empírico sobre los $n+k-1$ registros es
\[
\hat R_k(f)=\frac{1}{n+k-1}\Big(k\,\ell(f(x_1),y_1)+\sum_{i=2}^n\ell(f(x_i),y_i)\Big),
\]
es decir, el riesgo ponderado en el que el registro duplicado tiene peso $k$ y los demás peso $1$.
\end{proposicion}
\begin{proof}
\noindent Las $k$ copias contribuyen cada una $\ell(f(x_1),y_1)$ a la suma, y el total de registros es $n-1+k$.
\end{proof}

\noindent En consecuencia, el minimizador de $\hat R_k$ se desplaza hacia ajustar mejor $(x_1,y_1)$ a costa de los demás, y cualquier estadístico derivado (una media, un conteo como ``transacciones en 24 horas'') trata ese registro como $k$ observaciones independientes, lo que subestima varianzas y sesga conteos. Si una copia cae en $\Dtr$ y otra en $\Dte$, el modelo se evalúa en un punto con el que entrenó, lo que es fuga de datos en el sentido amplio; por eso la deduplicación debe hacerse antes de dividir. Los casi duplicados (``Juan Pérez'' y ``JUAN PEREZ'') son el mismo registro bajo una normalización de texto $\nu$ (pasar a mayúsculas, quitar acentos): se detectan comparando $\nu(z_i)=\nu(z_j)$ en lugar de $z_i=z_j$.

\section{Selección y extracción de características}\label{sec:seleccion}

\noindent \emph{Motivación: la maldición de la dimensionalidad.} La guía afirma que con muchas características ``las observaciones quedan más dispersas''. Un cálculo lo precisa. Si los datos son uniformes en el cubo $[0,1]^p$ y se quiere que un vecindario cúbico de cada punto contenga una fracción $r$ de los datos, su lado debe ser $\ell$ con $\ell^p=r$, es decir, $\ell=r^{1/p}$. Para $r=0.01$ y $p=1$ basta $\ell=0.01$; para $p=10$ hace falta $\ell=0.01^{1/10}\approx0.63$, y para $p=100$, $\ell\approx0.955$. En dimensión alta, un ``vecindario'' que contenga el $1\%$ de los datos abarca casi todo el rango de cada coordenada y deja de ser local. Equivalentemente, para mantener fija la densidad de puntos por celda de lado $\ell$ se necesitan del orden de $\ell^{-p}$ observaciones, un número que crece exponencialmente con $p$. De ahí la necesidad de reducir $p$, que la guía divide en dos familias: \emph{seleccionar} un subconjunto de las columnas o \emph{extraer} columnas nuevas que combinen las originales.

\subsection{Métodos de filtro}\label{sec:filtros}

\noindent Un método de filtro puntúa cada característica $X_j$ por su asociación con la respuesta $Y$, sin ajustar el modelo final, y conserva las de mayor puntuación. La guía menciona tres puntuaciones: correlación, información mutua y la prueba $\chi^2$.

\medskip

\noindent La correlación de Pearson $\Corr(X,Y)=\Cov(X,Y)/\sqrt{\Var X\Var Y}$ es la puntuación más antigua, pero solo mide dependencia \emph{lineal}.

\begin{ejemplo}[Dependencia total con correlación nula]\label{ej:corr-nula}
\noindent Sea $X$ uniforme en $\{-1,0,1\}$ e $Y=X^2$. Entonces $\E[X]=0$ y $\E[X^3]=\frac13(-1+0+1)=0$, así que $\Cov(X,Y)=\E[X\cdot X^2]-\E[X]\E[X^2]=0$ y la correlación es nula, aunque $Y$ es una función exacta de $X$. Un filtro por correlación descartaría una característica perfectamente predictiva.
\end{ejemplo}

\noindent La información mutua, introducida por Shannon (1948) en teoría de la comunicación, detecta cualquier forma de dependencia.

\begin{definicion}[Información mutua]\label{def:im}
\noindent Para variables discretas $X,Y$ con función de masa conjunta $p(x,y)$ y marginales $p(x),p(y)$, la \emph{información mutua} es
\[
I(X;Y)=\sum_{x,y:\,p(x,y)>0}p(x,y)\log\frac{p(x,y)}{p(x)p(y)}.
\]
\end{definicion}

\noindent La definición compara la distribución conjunta con la que tendrían $X$ e $Y$ si fueran independientes, $p(x)p(y)$; el logaritmo del cociente vale cero exactamente cuando ambas coinciden en ese punto.

\begin{teorema}[La información mutua caracteriza la independencia]\label{teo:im}
\noindent $I(X;Y)\ge0$, con igualdad si y solo si $X$ e $Y$ son independientes.
\end{teorema}
\begin{proof}
\noindent Sea $S=\{(x,y):p(x,y)>0\}$. Usamos la desigualdad $\log t\le t-1$ para $t>0$, con igualdad solo en $t=1$ (la función $t-1-\log t$ tiene derivada $1-1/t$, negativa en $(0,1)$ y positiva en $(1,\infty)$, y vale $0$ en $t=1$). Con $t=p(x)p(y)/p(x,y)$,
\begin{align*}
-I(X;Y)&=\sum_{S}p(x,y)\log\frac{p(x)p(y)}{p(x,y)}\le\sum_{S}p(x,y)\Big(\frac{p(x)p(y)}{p(x,y)}-1\Big)\\
&=\sum_Sp(x)p(y)-1\le\sum_{x,y}p(x)p(y)-1=0.
\end{align*}
Luego $I\ge0$. Si hay igualdad, ambas desigualdades son igualdades. La primera exige $p(x)p(y)=p(x,y)$ en todo $(x,y)\in S$; la segunda exige $\sum_{S}p(x)p(y)=1$, es decir, $p(x)p(y)=0$ fuera de $S$, donde $p(x,y)=0$. Juntas dan $p(x,y)=p(x)p(y)$ en todo punto, que es la independencia. Recíprocamente, si hay independencia cada logaritmo vale $\log1=0$.
\end{proof}

\noindent En el Ejemplo~\ref{ej:corr-nula}, $Y$ vale $0$ con probabilidad $1/3$ y $1$ con probabilidad $2/3$, y dado $X$ el valor de $Y$ es determinista. Por la identidad $I(X;Y)=H(Y)-H(Y\mid X)$ (que se obtiene escribiendo $\log\frac{p(x,y)}{p(x)p(y)}=\log\frac{p(y\mid x)}{p(y)}$ y separando la suma), con $H(Y\mid X)=0$ se obtiene $I(X;Y)=H(Y)=-\frac13\log\frac13-\frac23\log\frac23=\log3-\frac23\log2>0$. La información mutua detecta lo que la correlación no ve. Su desventaja es estadística: para variables continuas hay que estimar densidades o discretizar, y el estimador tiene sesgo y varianza apreciables con muestras pequeñas.

\medskip

\noindent Para una característica categórica y una respuesta categórica, la prueba de independencia en una tabla de contingencia es la puntuación clásica, y está ligada a la información mutua.

\begin{proposicion}[Estadístico $G$ e información mutua]\label{prop:g-im}
\noindent Sea una tabla de contingencia con conteos observados $O_{ab}$, total $n$, marginales $O_{a\cdot},O_{\cdot b}$ y conteos esperados bajo independencia $E_{ab}=O_{a\cdot}O_{\cdot b}/n$. Sea $\hat I$ la información mutua (con logaritmo natural) de la distribución empírica $\hat p(a,b)=O_{ab}/n$. Entonces el estadístico de razón de verosimilitudes $G=2\sum_{a,b}O_{ab}\log(O_{ab}/E_{ab})$ cumple $G=2n\hat I$. Además, si $O_{ab}=E_{ab}(1+\delta_{ab})$ con $\max\abs{\delta_{ab}}\to0$, entonces $G=\sum_{a,b}(O_{ab}-E_{ab})^2/E_{ab}+o\big(\sum_{a,b}E_{ab}\delta_{ab}^2\big)$, es decir, $G$ y el estadístico $\chi^2$ de Pearson coinciden a segundo orden.
\end{proposicion}
\begin{proof}
\noindent Como $\hat p(a)=O_{a\cdot}/n$ y $\hat p(b)=O_{\cdot b}/n$, se tiene $\frac{\hat p(a,b)}{\hat p(a)\hat p(b)}=\frac{O_{ab}/n}{O_{a\cdot}O_{\cdot b}/n^2}=\frac{O_{ab}}{E_{ab}}$. Entonces $n\hat I=\sum_{a,b}n\hat p(a,b)\log\frac{O_{ab}}{E_{ab}}=\sum_{a,b}O_{ab}\log\frac{O_{ab}}{E_{ab}}=G/2$.

\medskip

\noindent Para la segunda parte, el desarrollo de Taylor $(1+\delta)\log(1+\delta)=\delta+\frac{\delta^2}{2}+O(\delta^3)$ (producto de $1+\delta$ por $\log(1+\delta)=\delta-\frac{\delta^2}2+O(\delta^3)$) da $O_{ab}\log\frac{O_{ab}}{E_{ab}}=E_{ab}\big(\delta_{ab}+\frac{\delta_{ab}^2}{2}+O(\delta_{ab}^3)\big)$. Sumando, $\sum_{a,b}E_{ab}\delta_{ab}=\sum_{a,b}(O_{ab}-E_{ab})=n-n=0$, porque la suma de los esperados es $\sum_{a,b}O_{a\cdot}O_{\cdot b}/n=n\cdot n/n=n$. Queda $G=\sum_{a,b}E_{ab}\delta_{ab}^2+O(\sum E_{ab}\abs{\delta_{ab}}^3)$, y $E_{ab}\delta_{ab}^2=(O_{ab}-E_{ab})^2/E_{ab}$.
\end{proof}

\noindent Así, ordenar características categóricas por el estadístico $\chi^2$ o por información mutua empírica da, para tablas del mismo tamaño y cerca de la independencia, el mismo orden. La prueba $\chi^2$ añade un umbral de significación con la distribución asintótica $\chi^2_{(r-1)(c-1)}$; la información mutua se extiende a respuestas y características no categóricas.

\subsection{Redundancia, colinealidad y fuga del objetivo}\label{sec:redundancia}

\noindent La guía elimina los ingresos porque ``se mueven en paralelo a las ventas''. El caso extremo de ese paralelismo es la colinealidad exacta, que tiene una consecuencia algebraica precisa.

\begin{proposicion}[Colinealidad exacta]\label{prop:colineal}
\noindent Si alguna columna de $X$ es combinación lineal de otras, entonces $X^\top X$ es singular y el estimador de mínimos cuadrados no es único.
\end{proposicion}
\begin{proof}
\noindent Si $x_{\cdot k}=\sum_{j\ne k}c_jx_{\cdot j}$, el vector $v$ con $v_k=-1$ y $v_j=c_j$ ($j\ne k$) es no nulo y cumple $Xv=0$. Entonces $X^\top Xv=0$, así que $X^\top X$ tiene núcleo no trivial y es singular. Si $\hat\beta$ resuelve las ecuaciones normales $X^\top X\beta=X^\top y$, también las resuelve $\hat\beta+tv$ para todo $t\in\R$, y todas producen las mismas predicciones $X\hat\beta$.
\end{proof}

\noindent Con ingresos $=$ precio $\times$ ventas y precio constante, la colinealidad es exacta y los coeficientes de ventas e ingresos no están identificados. Con colinealidad aproximada, $X^\top X$ es casi singular y la varianza de los coeficientes se dispara. La selección elimina la columna redundante; la extracción (PCA, abajo) la funde con su pareja.

\medskip

\noindent Hay una razón más fuerte para descartar los ingresos del día al predecir las ventas del día: esos ingresos solo se conocen después de las ventas. Formalmente, si la predicción para el instante $t$ se hace con la información disponible en $t$, toda característica debe ser función de esa información; una característica calculada con datos posteriores es \emph{fuga del objetivo} (\emph{target leakage}). Se formaliza con series de tiempo en la Sección~\ref{sec:series}.

\subsection{Extracción lineal: análisis de componentes principales}\label{sec:pca}

\noindent \emph{Motivación.} Pearson (1901) buscaba la recta o el plano que mejor se ajustan a una nube de puntos cuando todas las coordenadas tienen error, y Hotelling (1933) buscaba pocas combinaciones de muchas variables correlacionadas (pruebas psicológicas) que resumieran su variación. En ambos casos las características originales son redundantes y se busca un sistema de coordenadas con menos ejes que conserve lo esencial. La guía propone PCA después de one-hot y para imágenes; el objeto es el mismo.

\begin{definicion}[Componentes principales]\label{def:pca}
\noindent Sea $X\in\R^{n\times p}$ con columnas centradas y $S=\frac1{n-1}X^\top X$ su matriz de covarianza muestral, simétrica y semidefinida positiva, con valores propios $\lambda_1\ge\dots\ge\lambda_p\ge0$ y vectores propios ortonormales $v_1,\dots,v_p$. La \emph{$k$-ésima componente principal} de un registro $x$ es $v_k^\top x$, y la \emph{proporción de varianza explicada} por las $q$ primeras es $\sum_{k\le q}\lambda_k/\sum_{k\le p}\lambda_k$.
\end{definicion}

\noindent La definición es la solución de un problema de optimización, que es lo que la hace natural.

\begin{teorema}[Máxima varianza]\label{teo:pca}
\noindent Para todo $w\in\R^p$ con $\norm{w}=1$, la varianza muestral de la proyección $Xw$ es $w^\top Sw\le\lambda_1$, con igualdad en $w=v_1$. Más en general, entre los $w$ unitarios ortogonales a $v_1,\dots,v_{k-1}$, el máximo de $w^\top Sw$ es $\lambda_k$ y se alcanza en $w=v_k$.
\end{teorema}
\begin{proof}
\noindent La varianza muestral de $Xw$ (columna centrada, porque las de $X$ lo son) es $\frac1{n-1}(Xw)^\top(Xw)=w^\top Sw$. Escribiendo $w=\sum_ia_iv_i$ en la base ortonormal, $\norm{w}^2=\sum_ia_i^2=1$ y, como $Sv_i=\lambda_iv_i$ y $v_i^\top v_j=\ind{i=j}$, $w^\top Sw=\sum_i\lambda_ia_i^2$. Si $w\perp v_1,\dots,v_{k-1}$, entonces $a_1=\dots=a_{k-1}=0$ y $w^\top Sw=\sum_{i\ge k}\lambda_ia_i^2\le\lambda_k\sum_{i\ge k}a_i^2=\lambda_k$, porque $\lambda_i\le\lambda_k$ para $i\ge k$. La igualdad se alcanza con $a_k=1$ y los demás nulos, es decir, $w=v_k$. El caso $k=1$ es el primer enunciado.
\end{proof}

\noindent PCA depende de la escala de las columnas, lo que obliga a estandarizar antes (Sección~\ref{sec:escalado}).

\begin{ejemplo}[PCA sin estandarizar y con características redundantes]\label{ej:pca-escala}
\noindent (a) Si dos características son incorrelacionadas con varianzas $1$ y $10^4$ (por ejemplo, la misma magnitud en metros y otra en centímetros), $S=\mathrm{diag}(1,10^4)$, la primera componente es el eje de la segunda característica y explica $10^4/(10^4+1)=99.99\%$ de la varianza: PCA ``descubre'' la unidad de medida, no una estructura. Tras estandarizar, $S=I$, todo vector unitario tiene varianza $1$ y no hay dirección privilegiada.

\medskip

\noindent (b) Si dos características estandarizadas tienen correlación $\rho>0$ (ventas e ingresos), su matriz de covarianza es $R=\begin{psmallmatrix}1&\rho\\\rho&1\end{psmallmatrix}$. Se verifica que $R(1,1)^\top=(1+\rho)(1,1)^\top$ y $R(1,-1)^\top=(1-\rho)(1,-1)^\top$, así que $\lambda_1=1+\rho$ con $v_1=(1,1)/\sqrt2$ y $\lambda_2=1-\rho$. La primera componente, la suma normalizada de ambas, explica $(1+\rho)/2$; con $\rho=0.98$, el $99\%$. PCA funde las dos columnas redundantes en una.
\end{ejemplo}

\noindent \emph{Contraste selección frente a extracción.} La selección conserva columnas originales, lo que preserva la interpretación y permite dejar de recolectar las descartadas; con filtros univariantes ignora interacciones (una característica inútil por sí sola puede ser útil combinada con otra). La extracción por PCA combina todas las columnas, por lo que sigue necesitando recolectarlas todas y pierde interpretación, pero aprovecha la redundancia (Ejemplo~\ref{ej:pca-escala}(b)). PCA es no supervisado: maximiza varianza de $X$ sin mirar $Y$, y una dirección de poca varianza puede ser la más predictiva. Los vectores $v_k$ son parámetros ajustados y deben estimarse con $\Dtr$.

\section{Escalado de características numéricas}\label{sec:escalado}

\subsection{Por qué importa la escala}\label{sec:por-que-escala}

\noindent La guía afirma que el escalado evita que las características grandes ``dominen'' y acelera el descenso de gradiente. Antes de estudiar cada escalador conviene demostrar cuándo la escala importa y cuándo no, porque la respuesta depende del algoritmo. Se consideran cuatro casos.

\medskip

\noindent \emph{(1) Métodos basados en distancias.} La distancia euclídea $d(x,x')^2=\sum_j(x_j-x'_j)^2$ es una suma de contribuciones por coordenada. Si la característica $j$ se multiplica por $c$, su contribución se multiplica por $c^2$. Con edad en años (diferencias típicas de $10$) e ingreso en pesos (diferencias típicas de $10^4$), la contribución del ingreso es del orden de $10^8$ frente a $10^2$, y los vecinos más cercanos se deciden solo por el ingreso. Esto afecta a $k$ vecinos, a $k$-medias y a las máquinas de vectores de soporte con núcleo radial.

\medskip

\noindent \emph{(2) Descenso de gradiente.} Para la pérdida cuadrática de la regresión lineal con columnas centradas, $L(\beta)=\frac1{2n}\norm{y-X\beta}^2$, el gradiente es $\nabla L(\beta)=\frac1n X^\top(X\beta-y)$ y la matriz hessiana es $H=\frac1nX^\top X$. El siguiente teorema muestra que la velocidad del descenso de gradiente depende del número de condición de $H$, que la escala controla.

\begin{teorema}[Descenso de gradiente en una cuadrática]\label{teo:gd}
\noindent Sea $L(\beta)=\frac12(\beta-\beta^*)^\top H(\beta-\beta^*)+c$ con $H$ simétrica y valores propios $0<\mu=\lambda_{\min}\le\lambda_{\max}=L_H$, y sea $\kappa=L_H/\mu$. La iteración $\beta_{t+1}=\beta_t-\eta\nabla L(\beta_t)$ con $\eta=2/(L_H+\mu)$ cumple
\[
\norm{\beta_{t}-\beta^*}\le\Big(\frac{\kappa-1}{\kappa+1}\Big)^t\norm{\beta_0-\beta^*},
\]
y ningún paso fijo $\eta$ da un factor de contracción uniforme menor.
\end{teorema}
\begin{proof}
\noindent Como $\nabla L(\beta)=H(\beta-\beta^*)$, se tiene $\beta_{t+1}-\beta^*=(I-\eta H)(\beta_t-\beta^*)$. La matriz $I-\eta H$ es simétrica con valores propios $1-\eta\lambda_i$, así que su norma espectral es $\max_i\abs{1-\eta\lambda_i}\le\max_{\lambda\in[\mu,L_H]}\abs{1-\eta\lambda}$. La función $\lambda\mapsto1-\eta\lambda$ es afín, por lo que el máximo de su valor absoluto en el intervalo se alcanza en un extremo: vale $\max(\abs{1-\eta\mu},\abs{1-\eta L_H})$. Este máximo se minimiza cuando $1-\eta\mu=-(1-\eta L_H)$, es decir, $\eta=2/(L_H+\mu)$, y entonces vale $1-\frac{2\mu}{L_H+\mu}=\frac{L_H-\mu}{L_H+\mu}=\frac{\kappa-1}{\kappa+1}$. Cualquier otro $\eta$ aumenta uno de los dos términos, y ese término es el factor real sobre el vector propio correspondiente. Iterando la desigualdad $\norm{\beta_{t+1}-\beta^*}\le\frac{\kappa-1}{\kappa+1}\norm{\beta_t-\beta^*}$ se obtiene la cota.
\end{proof}

\noindent Si dos características incorrelacionadas tienen varianzas $1$ y $10^4$, $H=\mathrm{diag}(1,10^4)$ (con $n$ grande y columnas centradas), $\kappa=10^4$ y el factor es $0.9998$: reducir el error en un factor $10^6$ exige $t\ge\ln(10^6)/(-\ln0.9998)\approx6.9\times10^4$ iteraciones. Tras estandarizar, $H$ es la matriz de correlación, que aquí es $I$; entonces $\kappa=1$, el factor es $0$ y basta una iteración. Esta es la razón real, que la guía enuncia sin demostrar, de que la estandarización acelere el entrenamiento de modelos lineales y redes neuronales: no se debe a que calculen distancias, sino al condicionamiento de la hessiana.

\medskip

\noindent \emph{(3) Mínimos cuadrados ordinarios y regularización.} Sin regularización, el escalado no cambia las predicciones.

\begin{proposicion}[Invariancia de MCO frente a transformaciones afines por columna]\label{prop:mco-inv}
\noindent Sea $\tilde X=XD+\mathbf{1}c^\top$ con $D$ diagonal invertible y $c\in\R^p$ (cada columna se escala y se desplaza). Si el modelo incluye intercepto, los valores ajustados por mínimos cuadrados con $[\mathbf 1,X]$ y con $[\mathbf 1,\tilde X]$ coinciden.
\end{proposicion}
\begin{proof}
\noindent Los valores ajustados son la proyección ortogonal de $y$ sobre el espacio columna de la matriz de diseño. Cada columna de $\tilde X$ es $d_jx_{\cdot j}+c_j\mathbf 1$, que está en el espacio columna de $[\mathbf 1,X]$; recíprocamente $x_{\cdot j}=(\tilde x_{\cdot j}-c_j\mathbf 1)/d_j$ está en el de $[\mathbf 1,\tilde X]$. Los dos espacios columna coinciden, luego también las proyecciones.
\end{proof}

\noindent Con penalización, en cambio, la escala sí importa. En la regresión ridge se minimiza $\norm{y-X\beta}^2+\alpha\norm{\beta}^2$. Si una columna se multiplica por $c$, el coeficiente que produce la misma predicción se divide por $c$ y su penalización se divide por $c^2$: medir una característica en milímetros en vez de metros ($c=1000$) reduce su penalización efectiva en un factor $10^6$. La estandarización hace que la penalización trate a todos los coeficientes en la misma escala. Esta es la segunda razón real por la que la guía recomienda estandarizar para regresión lineal y logística (que en la práctica suelen ajustarse con regularización).

\medskip

\noindent \emph{(4) Árboles de decisión.} Los árboles son invariantes a cualquier transformación estrictamente creciente de cada característica, no solo a las afines.

\begin{proposicion}[Invariancia de los cortes de un árbol]\label{prop:arbol-inv}
\noindent Sea $g:\R\to\R$ estrictamente creciente. Para la muestra $x_{1j},\dots,x_{nj}$ de una característica, la familia de particiones de $\{1,\dots,n\}$ de la forma $\{i:x_{ij}\le t\}$, $t\in\R$, coincide con la familia de la forma $\{i:g(x_{ij})\le t'\}$, $t'\in\R$.
\end{proposicion}
\begin{proof}
\noindent Como $g$ es estrictamente creciente, $x_{ij}\le x_{kj}$ si y solo si $g(x_{ij})\le g(x_{kj})$. Dado $t$, sea $A=\{i:x_{ij}\le t\}$. Si $A$ es vacío, se toma $t'$ menor que todos los $g(x_{ij})$. Si no, sea $a=\max_{i\in A}x_{ij}$ y $t'=g(a)$: entonces $g(x_{ij})\le g(a)$ si y solo si $x_{ij}\le a$, y $x_{ij}\le a$ si y solo si $x_{ij}\le t$ (porque ningún dato cae en $(a,t]$ por maximalidad de $a$). La inclusión recíproca es la misma con $g^{-1}$ en lugar de $g$, que es estrictamente creciente en la imagen de $g$.
\end{proof}

\noindent Como un árbol elige cada corte comparando la impureza de las particiones posibles, y estas coinciden, el árbol ajustado sobre $g(x_{\cdot j})$ define las mismas regiones que sobre $x_{\cdot j}$. Escalar, estandarizar o tomar logaritmos no cambia un árbol (salvo por desempates y por la posición exacta del umbral entre dos datos consecutivos).

\subsection{La familia de escaladores afines}\label{sec:afines}

\noindent Los cuatro escaladores de la guía (MinMax, puntuación Z, robusto y MaxAbs) son casos de una misma construcción, y lo que los distingue es qué estadísticos se usan como posición y escala.

\begin{definicion}[Escalador afín]\label{def:afin}
\noindent Un \emph{escalador afín} de una columna es una transformación ajustada $T_{(c,s)}(x)=(x-c)/s$ con $s>0$, donde $c=\hat c(\Dtr)$ y $s=\hat s(\Dtr)$ son estadísticos de posición y escala calculados con los datos de entrenamiento.
\end{definicion}

\noindent La elección de $(c,s)$ da:
\begin{center}
\begin{tabular}{@{}lll@{}}
\toprule
Escalador & $c$ & $s$\\
\midrule
MinMax (normalización) & $\min_i x_i$ & $\max_ix_i-\min_ix_i$\\
Puntuación Z (estandarización) & $\bar x$ & desviación estándar $s_0$\\
Robusto & mediana & $\IQR$\\
MaxAbs & $0$ & $\max_i\abs{x_i}$\\
\bottomrule
\end{tabular}
\end{center}

\noindent La guía repite en sus puntos esenciales que ninguno de estos escaladores ``corrige el sesgo''. El resultado siguiente lo demuestra para toda la familia a la vez.

\begin{proposicion}[Los escaladores afines conservan la forma]\label{prop:forma}
\noindent Sea $X$ con momentos de tercer y cuarto orden finitos y $\sigma>0$, y sea $Y=(X-c)/s$ con $s>0$. Entonces la asimetría $\gamma_1=\E[(X-\mu)^3]/\sigma^3$ y la curtosis $\gamma_2=\E[(X-\mu)^4]/\sigma^4$ de $Y$ coinciden con las de $X$. Más aún, la función de distribución de $Y$ es $F_Y(y)=F_X(sy+c)$, de modo que $Y$ tiene la misma forma que $X$ salvo cambio de posición y unidad.
\end{proposicion}
\begin{proof}
\noindent $\E[Y]=(\mu-c)/s$ y $Y-\E Y=(X-\mu)/s$, así que $\sigma_Y=\sigma/s$ y $\E[(Y-\E Y)^3]/\sigma_Y^3=\frac{\E[(X-\mu)^3]/s^3}{\sigma^3/s^3}=\gamma_1(X)$; el cálculo para el cuarto momento es idéntico con exponente $4$. Para la distribución, $F_Y(y)=\Prob((X-c)/s\le y)=\Prob(X\le sy+c)$ porque $s>0$.
\end{proof}

\noindent Todo cambio de forma exige una transformación no lineal (Sección~\ref{sec:nolineales}). Lo que distingue a los escaladores afines entre sí es su comportamiento ante atípicos y ceros.

\begin{proposicion}[Propiedades de cada escalador]\label{prop:escaladores}
\noindent Sea una columna $x_1,\dots,x_n$.
\begin{enumerate}[label=(\roman*)]
\item \emph{Puntuación Z.} La columna transformada tiene media $0$ y desviación estándar $s_0$ igual a $1$.
\item \emph{MinMax.} La columna transformada tiene mínimo $0$ y máximo $1$. Si $n-1$ valores están en $[a,b]$, con $a$ igual al mínimo de la columna, y el restante vale $M>b$, los $n-1$ valores transformados quedan en $[0,(b-a)/(M-a)]$, intervalo cuya longitud tiende a $0$ cuando $M\to\infty$.
\item \emph{Robusto.} La mediana transformada es $0$ y el $\IQR$ transformado es $1$; reemplazar un único valor no acotado cambia la mediana y el $\IQR$ en una cantidad acotada (si $n\ge5$).
\item \emph{MaxAbs.} Los valores transformados están en $[-1,1]$, conservan el signo y $T(0)=0$. En consecuencia, si la columna tiene una fracción $\pi$ de ceros, la transformada también, mientras que la puntuación Z y MinMax convierten los ceros en $-\bar x/s_0$ y $-\min x/(\max x-\min x)$, que son no nulos salvo que $\bar x=0$ o $\min x=0$.
\end{enumerate}
\end{proposicion}
\begin{proof}
\noindent (i) La media de $(x_i-\bar x)/s_0$ es $(\bar x-\bar x)/s_0=0$ y su desviación estándar es $s_0/s_0=1$ por el cálculo de la Proposición~\ref{prop:forma}.

\medskip

\noindent (ii) $T$ es creciente, así que lleva el mínimo a $(\min x-\min x)/(\max x-\min x)=0$ y el máximo a $1$. Con mínimo $a$ y máximo $M$, un valor $x\in[a,b]$ se lleva a $(x-a)/(M-a)\in[0,(b-a)/(M-a)]$, porque $0\le x-a\le b-a$. La cota $(b-a)/(M-a)$ tiende a $0$ cuando $M\to\infty$.

\medskip

\noindent (iii) La mediana y el $\IQR$ de la columna transformada son $(\med-\med)/\IQR=0$ y $\IQR/\IQR=1$, porque la transformación es afín creciente y los cuantiles de una transformación afín creciente son la transformación de los cuantiles (la interpolación lineal entre estadísticos de orden conmuta con las funciones afines crecientes). Que un solo valor no los arrastre es la Proposición~\ref{prop:ruptura}, ya que $1/n<1/4$ para $n\ge5$.

\medskip

\noindent (iv) $\abs{x_i/\max_k\abs{x_k}}\le1$, el signo no cambia al dividir entre un positivo y $0/s=0$. Las fórmulas para los otros escaladores son $T(0)=(0-c)/s$.
\end{proof}

\noindent La parte (iv) es la razón por la que la guía recomienda MaxAbs para datos dispersos: una matriz documento--término con $99.9\%$ de ceros se guarda almacenando solo las entradas no nulas; tras centrar, todas las entradas serían no nulas y la memoria se multiplicaría por mil. La parte (ii) es la debilidad de MinMax: en el ejemplo de la guía con $x=25$ como máximo, los otros nueve valores de $x$, que están en $[3.8,14]$, se llevan a $[0,0.48]$.

\medskip

\noindent \emph{Contraste.} La guía resume: normalización cuando importa el rango, estandarización cuando importa la distribución. Con lo demostrado, el contraste preciso es este.
\begin{itemize}
\item \emph{MinMax}: garantiza un rango acotado $[0,1]$ (útil para entradas que deben estar acotadas, como intensidades de píxeles o funciones de activación que se saturan), pero su escala depende de dos estadísticos con punto de ruptura $1/n$; un valor nuevo fuera del rango de entrenamiento se transforma fuera de $[0,1]$.
\item \emph{Puntuación Z}: produce media $0$ y varianza $1$, que es lo que condiciona bien la hessiana (Teorema~\ref{teo:gd}) y homogeneiza las penalizaciones; no acota el rango y su escala se infla con atípicos. No normaliza la distribución (Proposición~\ref{prop:forma}): la frase de la guía ``ideal para algoritmos que asumen normalidad'' debe leerse con ese matiz, y además la regresión lineal supone normalidad de los errores, no de las características.
\item \emph{Robusto}: sus parámetros resisten atípicos, pero el atípico sigue siendo atípico tras la transformación (es afín); en el ejemplo de la guía, $y=23$ se lleva a $(23-8.25)/5.375=2.744$, lejos del resto, cuyo máximo es $0.791$.
\item \emph{MaxAbs}: el único que conserva ceros y dispersión; hereda la sensibilidad de MinMax a un máximo extremo.
\end{itemize}

\section{Transformaciones no lineales}\label{sec:nolineales}

\noindent \emph{Motivación.} La Proposición~\ref{prop:forma} muestra que ningún escalador afín cambia la asimetría. Los modelos lineales clásicos, sin embargo, funcionan mejor cuando la relación con la respuesta es aditiva y la varianza de los errores es constante, y muchas magnitudes reales (ingresos, tiempos, concentraciones, conteos) son asimétricas a la derecha y tienen varianza que crece con la media. Tukey (1957) organizó las transformaciones de potencia $x\mapsto x^\lambda$ en una ``escalera'' según su fuerza, y Box y Cox (1964) propusieron elegir $\lambda$ por máxima verosimilitud para lograr a la vez normalidad aproximada, varianza constante y aditividad. Yeo y Johnson (2000) extendieron la familia a datos con ceros y negativos. Esta sección demuestra por qué estas transformaciones reducen la asimetría, las ordena por fuerza y deriva sus estimadores.

\subsection{Asimetría por cuantiles y transformaciones cóncavas}\label{sec:asimetria-cuantiles}

\noindent Para comparar transformaciones hace falta una medida de asimetría que se comporte bien bajo funciones crecientes. La asimetría por momentos no sirve para esto (los momentos de $g(X)$ no son funciones de los de $X$), pero los cuantiles sí se transportan.

\begin{definicion}[Función cuantil y razón de colas]\label{def:razon-colas}
\noindent Sea $F$ la función de distribución de $X$ y $Q(u)=\inf\{x:F(x)\ge u\}$, $u\in(0,1)$, su función cuantil. Para $u\in(0,\frac12)$ con $Q(u)<Q(\frac12)<Q(1-u)$ se define la \emph{razón de colas}
\[
R_u(X)=\frac{Q(1-u)-Q(\tfrac12)}{Q(\tfrac12)-Q(u)},
\]
y la \emph{asimetría por cuantiles} $B_u(X)=\dfrac{Q(1-u)+Q(u)-2Q(\tfrac12)}{Q(1-u)-Q(u)}$.
\end{definicion}

\noindent $R_u$ compara la distancia de la mediana al cuantil superior con la distancia al cuantil inferior: $R_u>1$ indica cola derecha más larga, $R_u=1$ simetría en esos cuantiles. Dividiendo numerador y denominador de $B_u$ entre $Q(\frac12)-Q(u)$ se obtiene $B_u=(R_u-1)/(R_u+1)$, que es función creciente de $R_u$ (su derivada es $2/(R_u+1)^2>0$) y toma valores en $(-1,1)$; con $u=\frac14$, $B_u$ es el coeficiente de Bowley. Ambas medidas son invariantes a transformaciones afines crecientes, lo que concuerda con la Proposición~\ref{prop:forma}.

\begin{lema}[Los cuantiles se transportan]\label{lem:cuantiles}
\noindent Sea $I$ un intervalo abierto con $\Prob(X\in I)=1$ y $g:I\to\R$ continua y estrictamente creciente. Entonces $Q_{g(X)}(u)=g(Q_X(u))$ para todo $u\in(0,1)$.
\end{lema}
\begin{proof}
\noindent Sea $q=Q_X(u)$, que pertenece a la clausura de $I$; supondremos $q\in I$ (el caso de un extremo se trata igual con límites). Como $g$ es estrictamente creciente, $g(X)\le g(x)$ si y solo si $X\le x$, así que $\Prob(g(X)\le g(x))=F(x)$ para $x\in I$. Por continuidad por la derecha de $F$, $F(q)\ge u$, luego $\Prob(g(X)\le g(q))\ge u$ y $Q_{g(X)}(u)\le g(q)$. Sea ahora $y<g(q)$. Si $y$ es menor que todos los valores de $g$ en $I$, $\Prob(g(X)\le y)=0<u$. En otro caso, por el teorema del valor intermedio (la imagen de $I$ por $g$ continua es un intervalo) existe $x\in I$ con $g(x)=y$ y, como $g$ es creciente, $x<q$; por definición de ínfimo, $F(x)<u$, es decir, $\Prob(g(X)\le y)<u$. Luego ningún $y<g(q)$ está en el conjunto que define $Q_{g(X)}(u)$, y $Q_{g(X)}(u)=g(q)$.
\end{proof}

\begin{lema}[Pendientes de cuerdas de una función cóncava]\label{lem:cuerdas}
\noindent Si $g$ es cóncava en un intervalo que contiene a $a<b<c$, entonces
\[
\frac{g(c)-g(b)}{c-b}\le\frac{g(c)-g(a)}{c-a}\le\frac{g(b)-g(a)}{b-a},
\]
con desigualdades estrictas si $g$ es estrictamente cóncava.
\end{lema}
\begin{proof}
\noindent Sea $t=\frac{c-b}{c-a}\in(0,1)$, de modo que $b=ta+(1-t)c$ (se verifica: $ta+(1-t)c=c-t(c-a)=c-(c-b)=b$). La concavidad da $g(b)\ge tg(a)+(1-t)g(c)$. Restando $g(a)$ a ambos lados, $g(b)-g(a)\ge(1-t)(g(c)-g(a))$, y como $1-t=\frac{b-a}{c-a}$, dividiendo entre $b-a>0$ se obtiene la desigualdad derecha. Restando la desigualdad de concavidad de $g(c)$, $g(c)-g(b)\le t(g(c)-g(a))$; dividiendo entre $c-b=t(c-a)$ se obtiene la izquierda. Con concavidad estricta la primera desigualdad es estricta y también lo son las siguientes.
\end{proof}

\begin{teorema}[Las transformaciones cóncavas reducen la asimetría derecha]\label{teo:concava}
\noindent Sea $g$ continua, estrictamente creciente y cóncava en un intervalo abierto $I$ con $\Prob(X\in I)=1$, y sea $u\in(0,\frac12)$ tal que $Q(u)<Q(\frac12)<Q(1-u)$. Entonces $R_u(g(X))\le R_u(X)$ y $B_u(g(X))\le B_u(X)$, con desigualdades estrictas si $g$ es estrictamente cóncava. Si $g$ es convexa, las desigualdades se invierten.
\end{teorema}
\begin{proof}
\noindent Escribimos $a=Q_X(u)$, $b=Q_X(\frac12)$, $c=Q_X(1-u)$, con $a<b<c$. Por el Lema~\ref{lem:cuantiles}, los cuantiles correspondientes de $g(X)$ son $g(a)<g(b)<g(c)$ (estrictamente, porque $g$ es estrictamente creciente). Por tanto
\[
R_u(g(X))=\frac{g(c)-g(b)}{g(b)-g(a)}=\frac{\frac{g(c)-g(b)}{c-b}}{\frac{g(b)-g(a)}{b-a}}\cdot\frac{c-b}{b-a}=\frac{m_{\mathrm{sup}}}{m_{\mathrm{inf}}}\,R_u(X),
\]
donde $m_{\mathrm{sup}}$ y $m_{\mathrm{inf}}$ son las pendientes de las cuerdas sobre $[b,c]$ y $[a,b]$, ambas positivas. Por el Lema~\ref{lem:cuerdas}, $m_{\mathrm{sup}}\le m_{\mathrm{inf}}$, luego $R_u(g(X))\le R_u(X)$, con desigualdad estricta si $g$ es estrictamente cóncava. Como $B_u$ es función creciente de $R_u$, lo mismo vale para $B_u$. Si $g$ es convexa, $-g$ es cóncava y el Lema~\ref{lem:cuerdas} aplicado a $-g$ invierte la desigualdad entre pendientes.
\end{proof}

\noindent El teorema es la versión precisa de la frase de la guía ``la transformación logarítmica comprime la cola derecha'': cualquier transformación creciente y cóncava acorta la distancia entre la mediana y los cuantiles altos en relación con la distancia a los cuantiles bajos. También explica el riesgo opuesto: aplicada a datos simétricos ($R_u=1$), una transformación estrictamente cóncava produce $R_u<1$, es decir, asimetría a la izquierda.

\begin{ejemplo}[El ejemplo de la guía]\label{ej:bowley}
\noindent Para $y=(4.7,5,6,7.5,8,8.5,11,12,12.5,23)$ (ordenado), con la definición de cuantil del Lema~\ref{lem:cuantiles} ($Q(u)$ es el menor dato con frecuencia acumulada $\ge u$): $Q(\frac14)=6$, $Q(\frac12)=8$, $Q(\frac34)=12$, así que $R_{1/4}(y)=(12-8)/(8-6)=2$. Para $\log_{10}y$, por el Lema~\ref{lem:cuantiles} los cuartiles son $\log_{10}6=0.7782$, $\log_{10}8=0.9031$ y $\log_{10}12=1.0792$, y $R_{1/4}=0.1761/0.1249\approx1.41$. En términos de Bowley, $B$ pasa de $1/3$ a $0.17$. (La asimetría por momentos pasa de $1.49$ a $0.52$.)
\end{ejemplo}

\subsection{La escalera de potencias y la familia de Box--Cox}\label{sec:boxcox}

\begin{definicion}[Transformación de Box--Cox]\label{def:boxcox}
\noindent Para $x>0$ y $\lambda\in\R$,
\[
\BC_\lambda(x)=\begin{cases}\dfrac{x^\lambda-1}{\lambda},&\lambda\ne0,\\[2mm] \log x,&\lambda=0.\end{cases}
\]
\end{definicion}

\noindent La resta de $1$ y la división entre $\lambda$ parecen arbitrarias, pero son lo que hace de la familia un objeto continuo en $\lambda$ en el que el logaritmo ocupa su lugar natural. Sin ellas, $x^\lambda$ tiende a la constante $1$ cuando $\lambda\to0$ y la familia pierde el logaritmo; y para $\lambda<0$, $x^\lambda$ es decreciente e invertiría el orden de los datos.

\begin{proposicion}[Propiedades de Box--Cox]\label{prop:bc-prop}
\noindent Para cada $x>0$: (i) $\lambda\mapsto\BC_\lambda(x)$ es continua, en particular $\lim_{\lambda\to0}\BC_\lambda(x)=\log x$; (ii) para cada $\lambda$, $x\mapsto\BC_\lambda(x)$ es estrictamente creciente con derivada $x^{\lambda-1}$; (iii) es estrictamente cóncava si $\lambda<1$, afín si $\lambda=1$ y estrictamente convexa si $\lambda>1$; (iv) $\BC_\lambda(1)=0$ para todo $\lambda$.
\end{proposicion}
\begin{proof}
\noindent (i) Para $\lambda\ne0$, $\BC_\lambda(x)=\frac{e^{\lambda\log x}-e^{0}}{\lambda-0}$ es el cociente incremental en $0$ de la función $\lambda\mapsto e^{\lambda\log x}$, cuya derivada en $0$ es $\log x\cdot e^0=\log x$. (ii) Para $\lambda\ne0$, $\frac{d}{dx}\frac{x^\lambda-1}\lambda=x^{\lambda-1}$, y para $\lambda=0$, $\frac{d}{dx}\log x=x^{-1}=x^{0-1}$; en ambos casos la derivada es positiva en $x>0$. (iii) La segunda derivada es $(\lambda-1)x^{\lambda-2}$, cuyo signo es el de $\lambda-1$. (iv) $1^\lambda-1=0$ y $\log1=0$.
\end{proof}

\noindent Por (iii) y el Teorema~\ref{teo:concava}, toda $\BC_\lambda$ con $\lambda<1$ reduce la asimetría derecha. La guía dice que la raíz cuadrada tiene un efecto ``menos marcado'' que el logaritmo; el resultado siguiente lo demuestra y ordena toda la familia.

\begin{teorema}[Orden de fuerza en la escalera de potencias]\label{teo:escalera}
\noindent Si $\lambda_2<\lambda_1$, existe $h$ continua, estrictamente creciente y cóncava tal que $\BC_{\lambda_2}=h\circ\BC_{\lambda_1}$ en $(0,\infty)$. En consecuencia, para $X>0$ y $u$ como en el Teorema~\ref{teo:concava}, $R_u(\BC_{\lambda_2}(X))\le R_u(\BC_{\lambda_1}(X))$. En particular, para datos positivos,
\[
R_u(X)\ \ge\ R_u(\sqrt X)\ \ge\ R_u(\sqrt[3]{X})\ \ge\ R_u(\log X).
\]
\end{teorema}
\begin{proof}
\noindent Por la Proposición~\ref{prop:bc-prop}(ii), $\BC_{\lambda_1}$ es una biyección continua estrictamente creciente de $(0,\infty)$ sobre su imagen $J$, que es un intervalo; sea $h=\BC_{\lambda_2}\circ\BC_{\lambda_1}^{-1}:J\to\R$, que es continua y estrictamente creciente por ser composición de funciones con esas propiedades. Por la regla de la cadena y la derivada de la función inversa, para $u\in J$ y $x=\BC_{\lambda_1}^{-1}(u)$,
\[
h'(u)=\frac{\BC_{\lambda_2}'(x)}{\BC_{\lambda_1}'(x)}=\frac{x^{\lambda_2-1}}{x^{\lambda_1-1}}=x^{\lambda_2-\lambda_1}.
\]
Como $\lambda_2-\lambda_1<0$, $x\mapsto x^{\lambda_2-\lambda_1}$ es decreciente, y como $x=\BC_{\lambda_1}^{-1}(u)$ es creciente en $u$, $h'$ es decreciente en $J$. Una función derivable con derivada decreciente es cóncava (por el teorema del valor medio, las pendientes de las cuerdas sobre intervalos consecutivos son valores de $h'$ en puntos ordenados, y por tanto decrecen, que es la condición de cuerdas del Lema~\ref{lem:cuerdas}). Aplicando el Teorema~\ref{teo:concava} a la variable $\BC_{\lambda_1}(X)$ y a la función $h$ se obtiene la desigualdad.

\medskip

\noindent Para la cadena final, $\BC_1(x)=x-1$, $\BC_{1/2}(x)=2(\sqrt x-1)$, $\BC_{1/3}(x)=3(\sqrt[3]x-1)$ y $\BC_0=\log$. Cada una difiere de $x$, $\sqrt x$, $\sqrt[3]x$ en una función afín creciente, y $R_u$ es invariante a tales funciones (sus cuantiles se transforman afínmente por el Lema~\ref{lem:cuantiles} y el cociente de diferencias no cambia). Basta aplicar la primera parte con $1>\frac12>\frac13>0$.
\end{proof}

\noindent La escalera es, pues, una familia totalmente ordenada por concavidad: bajar $\lambda$ comprime más la cola derecha, subirlo por encima de $1$ comprime la izquierda. Esto convierte la elección de transformación en la elección de un número, que es lo que hace posible estimarla.

\medskip

\noindent \emph{Estimación de $\lambda$ por máxima verosimilitud.} Box y Cox supusieron que para algún $\lambda$ las variables $\BC_\lambda(X_i)$ son i.i.d.\ $N(\mu,\sigma^2)$. La densidad de $X_i$ se obtiene por cambio de variable: si $Y=\BC_\lambda(X)$ con $\BC_\lambda$ creciente y derivable, $f_X(x)=f_Y(\BC_\lambda(x))\,\BC_\lambda'(x)=f_Y(\BC_\lambda(x))\,x^{\lambda-1}$. La log-verosimilitud de la muestra es entonces
\begin{equation}\label{eq:bc-lik}
\ell(\lambda,\mu,\sigma^2)=-\frac n2\log(2\pi\sigma^2)-\frac1{2\sigma^2}\sum_{i=1}^n\big(\BC_\lambda(x_i)-\mu\big)^2+(\lambda-1)\sum_{i=1}^n\log x_i.
\end{equation}
Para $\lambda$ fijo, maximizar en $\mu$ y $\sigma^2$ es el problema normal usual: $\hat\mu(\lambda)$ es la media de los $\BC_\lambda(x_i)$ y $\hat\sigma^2(\lambda)=\frac1n\sum_i(\BC_\lambda(x_i)-\hat\mu(\lambda))^2$. Sustituyendo, el segundo término vale $-n/2$, y la \emph{log-verosimilitud perfilada} es
\begin{equation}\label{eq:bc-perfil}
\ell_p(\lambda)=-\frac n2\log\hat\sigma^2(\lambda)+(\lambda-1)\sum_{i=1}^n\log x_i+\text{const}.
\end{equation}
El estimador es $\hat\lambda=\argmax_\lambda\ell_p(\lambda)$, que se calcula numéricamente; es lo que hace \texttt{PowerTransformer} de scikit-learn.

\begin{observacion}[Por qué el término jacobiano es imprescindible]\label{obs:jacobiano}
\noindent Si se omitiera $(\lambda-1)\sum\log x_i$ y se maximizara solo $-\frac n2\log\hat\sigma^2(\lambda)$, el problema sería degenerado. Supóngase que todos los $x_i>1$. Para $\lambda<0$, $\BC_\lambda(x_i)=\frac{1-x_i^{\lambda}}{\abs{\lambda}}\in\big(0,\frac1{\abs\lambda}\big)$, porque $0<x_i^\lambda<1$. Todos los valores transformados quedan en un intervalo de longitud $1/\abs\lambda$, así que $\hat\sigma^2(\lambda)\le1/\lambda^2\to0$ cuando $\lambda\to-\infty$, y $-\frac n2\log\hat\sigma^2(\lambda)\to+\infty$. El ``mejor'' $\lambda$ sería $-\infty$ solo porque encoge la escala. El jacobiano compensa exactamente ese cambio de escala, porque mide cuánto estira o encoge $\BC_\lambda$ el eje en cada dato.
\end{observacion}

\subsection{Yeo--Johnson: ceros y negativos}\label{sec:yj}

\noindent Box--Cox exige $x>0$. El remedio de desplazar, $\BC_\lambda(x+k)$, introduce un segundo parámetro $k$ sin una elección natural. Yeo y Johnson buscaron una familia definida en todo $\R$ que coincidiera con Box--Cox para datos positivos (desplazados en $1$), que tratara la parte negativa con la reflexión adecuada y que conservara la continuidad y el orden por concavidad.

\begin{definicion}[Transformación de Yeo--Johnson]\label{def:yj}
\noindent Para $x\in\R$ y $\lambda\in\R$,
\[
\YJ_\lambda(x)=\begin{cases}\BC_\lambda(x+1)=\dfrac{(x+1)^\lambda-1}{\lambda}\ (\log(x+1)\text{ si }\lambda=0),&x\ge0,\\[3mm]
-\BC_{2-\lambda}(1-x)=-\dfrac{(1-x)^{2-\lambda}-1}{2-\lambda}\ (-\log(1-x)\text{ si }\lambda=2),&x<0.\end{cases}
\]
\end{definicion}

\noindent La rama negativa es la reflexión $x\mapsto-\YJ(-x)$ de una Box--Cox con exponente $2-\lambda$. La elección $2-\lambda$ no es arbitraria: es la que hace que la segunda derivada tenga el mismo signo en ambas ramas, como se verá en (iii).

\begin{proposicion}[Propiedades de Yeo--Johnson]\label{prop:yj}
\noindent (i) $\YJ_\lambda$ es continua y derivable en $0$, con $\YJ_\lambda(0)=0$ y $\YJ_\lambda'(0)=1$. (ii) Es estrictamente creciente en $\R$, con derivada $(x+1)^{\lambda-1}$ si $x\ge0$ y $(1-x)^{1-\lambda}$ si $x<0$. (iii) Es estrictamente cóncava en $\R$ si $\lambda<1$, la identidad si $\lambda=1$ y estrictamente convexa si $\lambda>1$. (iv) El logaritmo del jacobiano de la muestra es $(\lambda-1)\sum_i\sgn(x_i)\log(\abs{x_i}+1)$, con $\sgn(0)=1$.
\end{proposicion}
\begin{proof}
\noindent (i) Ambas ramas valen $0$ en $x=0$: $\BC_\lambda(1)=0$ y $-\BC_{2-\lambda}(1)=0$ por la Proposición~\ref{prop:bc-prop}(iv). (ii) Para $x>0$, la derivada es $\BC_\lambda'(x+1)=(x+1)^{\lambda-1}$. Para $x<0$, por la regla de la cadena, $\frac{d}{dx}\big[-\BC_{2-\lambda}(1-x)\big]=\BC_{2-\lambda}'(1-x)=(1-x)^{1-\lambda}$. Ambas son positivas y ambas tienden a $1$ en $x=0$, lo que junto con la continuidad da (i). (iii) Para $x>0$ la segunda derivada es $(\lambda-1)(x+1)^{\lambda-2}$; para $x<0$ es $\frac{d}{dx}(1-x)^{1-\lambda}=-(1-\lambda)(1-x)^{-\lambda}=(\lambda-1)(1-x)^{-\lambda}$. En ambas ramas el signo es el de $\lambda-1$, y como la derivada primera es continua en $0$, la derivada es estrictamente monótona en todo $\R$, lo que da concavidad o convexidad estricta global. Si $\lambda=1$: para $x\ge0$, $(x+1)-1=x$; para $x<0$, $-((1-x)-1)=x$. (iv) El jacobiano es el producto de las derivadas de (ii); su logaritmo es $\sum_{x_i\ge0}(\lambda-1)\log(x_i+1)+\sum_{x_i<0}(1-\lambda)\log(1-x_i)$, que es la expresión enunciada porque $1-x_i=\abs{x_i}+1$ si $x_i<0$.
\end{proof}

\noindent Por (iii) y el Teorema~\ref{teo:concava}, $\lambda<1$ reduce la asimetría derecha y $\lambda>1$ la izquierda, en todo $\R$. La estimación de $\lambda$ es la de Box--Cox con el jacobiano de (iv) en lugar de $(\lambda-1)\sum\log x_i$. En el ejemplo de la guía con datos simétricos ($\{1,3,5,7,-1,-3\}$ y $\{2,4,6,8,-2,-4\}$), las $\hat\lambda$ estimadas son $0.915$ y $1.023$, cercanas a $1$: el estimador reconoce que no hay asimetría que corregir. \texttt{PowerTransformer} estandariza además la salida por defecto, por lo que el resultado tiene media $0$ y varianza $1$.

\subsection{Estabilización de la varianza}\label{sec:estabilizacion}

\noindent La otra motivación de Box y Cox era la varianza no constante. Si la varianza de una medición crece con su media, un modelo con errores homocedásticos no es adecuado. El método delta indica qué transformación la estabiliza.

\begin{teorema}[Método delta]\label{teo:delta}
\noindent Sean $T_n$ variables aleatorias con $\sqrt n(T_n-\theta)\xrightarrow{d}N(0,v(\theta))$ y $g$ con derivada continua en un entorno de $\theta$. Entonces $\sqrt n\,(g(T_n)-g(\theta))\xrightarrow{d}N\big(0,g'(\theta)^2v(\theta)\big)$.
\end{teorema}
\begin{proof}
\noindent Por el teorema del valor medio, $g(T_n)-g(\theta)=g'(\xi_n)(T_n-\theta)$ con $\xi_n$ entre $T_n$ y $\theta$ (cuando $T_n$ está en el entorno donde $g$ es derivable, lo que ocurre con probabilidad que tiende a $1$). Como $\sqrt n(T_n-\theta)$ converge en distribución, $T_n-\theta=n^{-1/2}\cdot\sqrt n(T_n-\theta)\to0$ en probabilidad; luego $\xi_n\to\theta$ en probabilidad y, por continuidad de $g'$, $g'(\xi_n)\to g'(\theta)$ en probabilidad. Por el teorema de Slutsky, $\sqrt n(g(T_n)-g(\theta))=g'(\xi_n)\sqrt n(T_n-\theta)\xrightarrow{d}g'(\theta)\,N(0,v(\theta))$, que es $N(0,g'(\theta)^2v(\theta))$.
\end{proof}

\noindent Para que la varianza límite no dependa de $\theta$ se necesita $g'(\theta)^2v(\theta)=\text{const}$, es decir, $g'(\theta)\propto v(\theta)^{-1/2}$.

\begin{corolario}[Raíz para conteos, logaritmo para magnitudes multiplicativas]\label{cor:estab}
\noindent (i) Si $T_n$ es la media de $n$ variables Poisson$(\theta)$, $v(\theta)=\theta$ y $g(t)=2\sqrt t$ da varianza límite $1$. (ii) Si $T_n$ es la media de $n$ exponenciales de media $\theta$, $v(\theta)=\theta^2$ y $g(t)=\log t$ da varianza límite $1$.
\end{corolario}
\begin{proof}
\noindent Por el teorema central del límite, $\sqrt n(T_n-\theta)\to N(0,v(\theta))$ con $v$ la varianza de una observación: $\theta$ para la Poisson y $\theta^2$ para la exponencial. (i) $g'(\theta)=\theta^{-1/2}$ y $g'(\theta)^2\theta=1$. (ii) $g'(\theta)=\theta^{-1}$ y $g'(\theta)^2\theta^2=1$. En ambos casos $g$ tiene derivada continua en $\theta>0$ y se aplica el Teorema~\ref{teo:delta}.
\end{proof}

\noindent Este es el fundamento de la afirmación de la guía de que la raíz y las transformaciones de potencia ``estabilizan la varianza'': la raíz cuadrada corresponde a varianza proporcional a la media (conteos), el logaritmo a desviación estándar proporcional a la media (magnitudes con error relativo constante, como ingresos o precios).

\subsection{Logaritmo y raíces: casos particulares y contraste}\label{sec:log-raices}

\noindent \emph{Linealización.} El logaritmo convierte relaciones multiplicativas en aditivas: si $y=a\,b^{x}$ con $a,b>0$, entonces $\log_{10}y=\log_{10}a+x\log_{10}b$, que es afín en $x$. En el ejemplo de la guía, $y=10^{x-3}$ para $x=3,4,5,6$ da $y=1,10,100,1000$, y $\log_{10}y=x-3$ es una recta exacta. Del mismo modo, un modelo multiplicativo $y=\beta_0x_1^{\beta_1}x_2^{\beta_2}\varepsilon$ se vuelve lineal en $\log y$, $\log x_1$ y $\log x_2$ con error aditivo $\log\varepsilon$.

\medskip

\noindent \emph{Ceros y negativos.} $\log x$ solo está definido para $x>0$. Para conteos con ceros se usa $\log(1+x)$, que coincide con $\YJ_0(x)$ para $x\ge0$ y con $\BC_0(x+1)$: la Proposición~\ref{prop:yj} lo sitúa en la familia de Yeo--Johnson. La raíz cuadrada está definida en $[0,\infty)$ y es $\BC_{1/2}$ salvo una función afín. La raíz cúbica está definida en todo $\R$ y es impar y estrictamente creciente, pero es cóncava en $(0,\infty)$ y convexa en $(-\infty,0)$ (su segunda derivada es $-\frac29x^{-5/3}$, de signo opuesto al de $x$). Por el Teorema~\ref{teo:concava}, en datos de un solo signo reduce la asimetría hacia la mediana; en datos con ambos signos comprime \emph{ambas} colas hacia cero y no corrige una asimetría en una dirección determinada. Además su derivada $\frac13x^{-2/3}$ no está acotada en $0$, así que separa mucho valores pequeños cercanos a cero, lo que puede amplificar ruido de medición.

\medskip

\noindent \emph{Contraste.}
\begin{center}
\begin{tabularx}{\linewidth}{@{}lYYY@{}}
\toprule
Transformación & Dominio & Parámetro ajustado & Efecto y limitaciones\\
\midrule
$\log x$, $\log(1+x)$ & $x>0$; $x>-1$ & ninguno & la más fuerte de las fijas; lineariza modelos multiplicativos; estabiliza varianza $\propto$ media$^2$; puede sobrecorregir datos casi simétricos\\
$\sqrt x$ & $x\ge0$ & ninguno & más débil que el logaritmo (Teorema~\ref{teo:escalera}); estabiliza varianza de conteos\\
$\sqrt[3]x$ & $\R$ & ninguno & entre la raíz cuadrada y el logaritmo en datos positivos; con ambos signos comprime las dos colas\\
Box--Cox & $x>0$ & $\lambda$ por MV & elige la fuerza con los datos; requiere positividad\\
Yeo--Johnson & $\R$ & $\lambda$ por MV & como Box--Cox en $\R$; $\lambda$ difícil de interpretar si hay ambos signos\\
\bottomrule
\end{tabularx}
\end{center}

\noindent Todas son estrictamente crecientes, por lo que (Proposición~\ref{prop:arbol-inv}) no cambian nada en un árbol y solo tienen efecto en modelos que usan los valores de forma aditiva o métrica. Los parámetros $\lambda$ se estiman con $\Dtr$. Frente al escalado robusto, que resiste atípicos pero no cambia la forma, estas transformaciones cambian la forma y acercan el atípico al resto sin eliminarlo.

\section{Discretización}\label{sec:binning}

\noindent \emph{Motivación.} La relación entre una característica continua y la respuesta puede no ser lineal ni monótona (la edad frente al riesgo de un seguro, que es alto en jóvenes y en mayores). Antes de los modelos flexibles, la forma más sencilla de capturarla con un modelo lineal era partir el rango en intervalos y estimar un nivel por intervalo. La guía presenta la discretización como paso de lo numérico a lo categórico; su contenido matemático es la aproximación por funciones escalonadas.

\begin{definicion}[Discretización]\label{def:binning}
\noindent Dados cortes $-\infty=t_0<t_1<\dots<t_{B-1}<t_B=\infty$, la \emph{discretización} de $x$ es el índice $b(x)=k$ si $x\in(t_{k-1},t_k]$. Es de \emph{ancho igual} si $t_k-t_{k-1}=h$ constante para $2\le k\le B-1$, y \emph{por cuantiles} si $t_k=\hat Q(k/B)$ para los cuantiles empíricos de $\Dtr$. La codificación \emph{one-hot} de $b(x)$ es el vector $(\ind{b(x)=1},\dots,\ind{b(x)=B})$.
\end{definicion}

\begin{proposicion}[Modelo lineal sobre intervalos: medias por celda]\label{prop:celdas}
\noindent Sea $G:\{1,\dots,n\}\to\{1,\dots,B\}$ una asignación de registros a grupos, todos no vacíos, y sea $D\in\R^{n\times B}$ la matriz de indicadores $D_{ik}=\ind{G(i)=k}$. El ajuste por mínimos cuadrados de $y\approx D\gamma$ tiene solución única $\hat\gamma_k=\bar y_k$, la media de $y$ en el grupo $k$.
\end{proposicion}
\begin{proof}
\noindent Cada fila de $D$ tiene exactamente un $1$, así que $(D^\top D)_{kl}=\sum_iD_{ik}D_{il}=n_k\ind{k=l}$, donde $n_k>0$ es el tamaño del grupo $k$; $D^\top D$ es diagonal invertible. Además $(D^\top y)_k=\sum_{i:G(i)=k}y_i$. Las ecuaciones normales $D^\top D\gamma=D^\top y$ dan $n_k\gamma_k=\sum_{G(i)=k}y_i$, es decir, $\gamma_k=\bar y_k$.
\end{proof}

\noindent Así, discretizar y codificar en one-hot convierte cualquier modelo lineal en uno que ajusta un nivel constante por intervalo. La calidad de esa aproximación depende del ancho de los intervalos.

\begin{proposicion}[Error de aproximación escalonada]\label{prop:escalon}
\noindent Si $f$ es $L$-Lipschitz en un intervalo $[a,a+h]$, la constante $c=f(a+h/2)$ cumple $\sup_{x\in[a,a+h]}\abs{f(x)-c}\le Lh/2$. La cota no puede mejorarse: para $f(x)=Lx$, $\min_c\sup_{[a,a+h]}\abs{f(x)-c}=Lh/2$.
\end{proposicion}
\begin{proof}
\noindent Para $x\in[a,a+h]$, $\abs{x-(a+h/2)}\le h/2$ y, por la condición de Lipschitz, $\abs{f(x)-f(a+h/2)}\le L\,h/2$. Para $f(x)=Lx$, la imagen del intervalo es $[La,La+Lh]$, de longitud $Lh$; cualquier constante $c$ dista al menos $Lh/2$ de uno de los dos extremos (si distara menos de $Lh/2$ de ambos, por la desigualdad triangular la longitud sería menor que $Lh$), y $c=La+Lh/2$ alcanza exactamente $Lh/2$.
\end{proof}

\noindent Las dos proposiciones exhiben el compromiso de la técnica: intervalos estrechos reducen el sesgo de aproximación (Proposición~\ref{prop:escalon}), pero cada nivel $\bar y_k$ se estima con $n_k$ datos y tiene varianza $\sigma^2/n_k$ si los errores son homocedásticos con varianza $\sigma^2$, que crece al estrechar los intervalos. Además, toda la información dentro de un intervalo se pierde: dos valores en el mismo intervalo reciben la misma predicción.

\medskip

\noindent \emph{Robustez.} La guía afirma que la discretización hace al modelo más robusto frente a atípicos y ruido. Con cortes por cuantiles, esto es exacto: cada corte $\hat Q(k/B)$ es un estadístico de orden de $\Dtr$, así que $b(x)$ depende de $x$ solo a través de sus comparaciones con esos estadísticos de orden; un atípico, por extremo que sea, supera a todos y recibe la etiqueta $B$, la misma que cualquier valor apenas superior a $\hat Q((B-1)/B)$. Los propios cortes, por la Proposición~\ref{prop:ruptura}, se mueven poco si se contamina una fracción pequeña de los datos. Además, por el Lema~\ref{lem:cuantiles}, la discretización por cuantiles es invariante a transformaciones estrictamente crecientes, lo que hace irrelevante tomar logaritmos antes. La discretización de ancho igual no tiene esta propiedad: un atípico estira el rango y vacía los intervalos intermedios.

\medskip

\noindent \emph{Contraste.} Frente a las transformaciones no lineales de la Sección~\ref{sec:nolineales}, la discretización permite relaciones no monótonas pero pierde resolución y añade $B-1$ parámetros; frente a un árbol, que elige los cortes mirando la respuesta, los cortes aquí son no supervisados. El ejemplo de la guía (peso del vehículo en cinco categorías y después cinco indicadores) es exactamente la Definición~\ref{def:binning} seguida de one-hot, y la Proposición~\ref{prop:celdas} dice qué estima un modelo lineal con ellos. Los cortes por cuantiles se ajustan con $\Dtr$.

\section{Codificación de variables categóricas}\label{sec:categoricas}

\noindent \emph{Motivación.} Una variable categórica toma valores en un conjunto finito $\mathcal{C}=\{c_1,\dots,c_K\}$ sin estructura aritmética: no tiene sentido sumar ``Rojo'' y ``Azul'' ni, si es nominal, decir que uno es mayor que otro. Los algoritmos, en cambio, operan en $\R^p$. Codificar es elegir una función $\phi:\mathcal{C}\to\R^d$, y toda elección introduce en $\mathcal{C}$ la geometría de $\R^d$ (distancias, orden, combinaciones lineales). El problema es que esa geometría no debe inventar relaciones inexistentes ni consumir demasiadas dimensiones. Las cuatro codificaciones de la guía son cuatro compromisos distintos entre esas dos exigencias. En esta sección $K=\abs{\mathcal{C}}$ es el número de categorías \emph{distintas} (la \emph{cardinalidad}); la guía llama a veces ``categorías'' a las observaciones, que pueden repetirse: en su ejemplo hay $10$ observaciones de $K=8$ colores.

\subsection{Codificación por etiquetas u ordinal}\label{sec:label}

\begin{definicion}[Codificación ordinal]\label{def:ordinal}
\noindent Es una biyección $\phi:\mathcal{C}\to\{0,1,\dots,K-1\}$, usada como una sola columna numérica.
\end{definicion}

\noindent La biyección induce un orden total y distancias $\abs{\phi(c)-\phi(c')}$ en $\mathcal{C}$. \texttt{LabelEncoder} de scikit-learn elige el orden alfabético (en $\{\text{White},\text{Black},\text{Red}\}$ asigna Black $=0$, Red $=1$, White $=2$), mientras que \texttt{OrdinalEncoder} con \texttt{categories} permite fijar el orden de una variable ordinal (Small $<$ Medium $<$ Large $<$ X-Large). Las consecuencias de ese orden dependen del modelo.

\begin{proposicion}[Codificación ordinal en un modelo lineal]\label{prop:ordinal-lineal}
\noindent Si la única representación de la categoría en un modelo lineal es la columna $\phi(c)$, las predicciones por categoría son $f(c)=\beta_0+\beta_1\phi(c)$. El vector de efectos $(f(c_1),\dots,f(c_K))$ pertenece al subespacio de dimensión $2$ generado por $\mathbf 1$ y $(\phi(c_1),\dots,\phi(c_K))$. Para $K\ge3$ no todo vector de efectos es representable, y el conjunto de los representables depende de la biyección elegida.
\end{proposicion}
\begin{proof}
\noindent El vector de efectos es $\beta_0\mathbf1+\beta_1(\phi(c_k))_k$, combinación lineal de dos vectores fijos, luego está en un subespacio de dimensión a lo sumo $2<K$ y no puede ser todo $\R^K$. Para la dependencia de la biyección, sea $K=3$ y sea $(0,\tfrac12,1)$ el vector de efectos de las categorías $(c_1,c_2,c_3)$. Con la codificación $\phi=(0,1,2)$ es representable, con $\beta_0=0$ y $\beta_1=\frac12$. Con la codificación $\phi=(1,0,2)$, que solo intercambia los códigos de $c_1$ y $c_2$, se necesitaría $\beta_0+\beta_1=0$, $\beta_0=\frac12$ y $\beta_0+2\beta_1=1$; las dos primeras dan $\beta_1=-\frac12$ y entonces $\beta_0+2\beta_1=-\frac12\ne1$, así que no es representable.
\end{proof}

\noindent Para una variable nominal, el orden alfabético no tiene relación con la respuesta, y el modelo lineal queda forzado a efectos equiespaciados en ese orden arbitrario; esta es la razón precisa de la advertencia de la guía (``Red sería el doble de Black''). Para una variable ordinal con efecto aproximadamente equiespaciado, la codificación ordinal es parsimoniosa: $2$ parámetros en lugar de $K$.

\medskip

\noindent Con árboles la situación es otra. Un corte sobre la columna ordinal separa $\mathcal{C}$ en $\{c:\phi(c)\le t\}$ y su complemento.

\begin{proposicion}[Cortes de un árbol sobre una codificación ordinal]\label{prop:ordinal-arbol}
\noindent (i) Un único corte produce a lo sumo $K-1$ biparticiones no triviales de $\mathcal{C}$, de las $2^{K-1}-1$ posibles. (ii) Con a lo sumo dos cortes sucesivos se puede aislar cualquier categoría.
\end{proposicion}
\begin{proof}
\noindent (i) Las biparticiones inducidas son $\{\phi\le j\}$ frente a $\{\phi>j\}$ para $j=0,\dots,K-2$, es decir, $K-1$. El total de biparticiones en dos partes no vacías se cuenta así: hay $2^K$ subconjuntos $A$, se excluyen $A=\emptyset$ y $A=\mathcal{C}$, y cada bipartición $\{A,A^c\}$ se contó dos veces, lo que da $(2^K-2)/2=2^{K-1}-1$. (ii) Para aislar la categoría de código $j$, el corte $\phi\le j$ deja a un lado los códigos $\{0,\dots,j\}$ y, dentro de él, el corte $\phi\le j-1$ separa $\{j\}$ (si $j=0$ o $j=K-1$ basta un corte).
\end{proof}

\noindent Así, un árbol puede representar cualquier efecto por categoría con la codificación ordinal, al precio de más cortes y de que las agrupaciones ``naturales'' que no son intervalos en el orden elegido requieren varios. Esa es la justificación de la regla de la guía (ordinal para árboles), que por la Proposición~\ref{prop:arbol-inv} tampoco depende de la escala de los códigos.

\subsection{Codificación one-hot}\label{sec:onehot}

\noindent \emph{Motivación.} Las variables indicadoras (\emph{dummy variables}) se introdujeron en la econometría de mediados del siglo XX para incluir factores cualitativos en regresiones sin imponerles orden. La idea es asignar a cada categoría su propio eje.

\begin{definicion}[Codificación one-hot]\label{def:onehot}
\noindent Es $\phi(c_k)=e_k\in\R^K$, el $k$-ésimo vector canónico.
\end{definicion}

\begin{proposicion}[Geometría y álgebra de one-hot]\label{prop:onehot}
\noindent (i) $\ip{e_k}{e_l}=\ind{k=l}$ y $\norm{e_k-e_l}=\sqrt2$ para $k\ne l$: todas las categorías son equidistantes. (ii) Un modelo lineal $f(c)=\gamma^\top\phi(c)$ puede representar cualquier vector de efectos en $\R^K$. (iii) Si además hay intercepto, las columnas de la matriz de diseño $[\mathbf1,D]$ son linealmente dependientes (``trampa de las variables ficticias'') y el estimador de mínimos cuadrados no es único. (iv) Si se elimina la columna de una categoría de referencia $r$, el ajuste por mínimos cuadrados con intercepto y las $K-1$ restantes es único, sus valores ajustados son las medias por categoría y los coeficientes son $\hat\beta_0=\bar y_r$ y $\hat\beta_k=\bar y_k-\bar y_r$.
\end{proposicion}
\begin{proof}
\noindent (i) El producto interno de vectores canónicos es la delta de Kronecker, y $\norm{e_k-e_l}^2=\norm{e_k}^2-2\ip{e_k}{e_l}+\norm{e_l}^2=1-0+1=2$. (ii) $f(c_k)=\gamma^\top e_k=\gamma_k$, así que el vector de efectos es $\gamma$, arbitrario. (iii) Cada fila de $D$ tiene un único $1$, luego $\sum_kD_{\cdot k}=\mathbf1$; la columna del intercepto es combinación lineal de las demás y se aplica la Proposición~\ref{prop:colineal}. (iv) El espacio columna de $[\mathbf1,D_{-r}]$ contiene $D_{\cdot r}=\mathbf1-\sum_{k\ne r}D_{\cdot k}$, luego coincide con el de $D$; las $K$ columnas de $[\mathbf1,D_{-r}]$ son independientes porque generan un espacio de dimensión $K$ (el de $D$, cuyas columnas tienen soportes disjuntos no vacíos). Por la Proposición~\ref{prop:celdas}, los valores ajustados, que son la proyección sobre ese espacio común, son las medias por categoría. En un registro de la categoría $r$ el valor ajustado es $\hat\beta_0$, luego $\hat\beta_0=\bar y_r$; en uno de la categoría $k\ne r$ es $\hat\beta_0+\hat\beta_k=\bar y_k$.
\end{proof}

\noindent La parte (i) es la ventaja conceptual de one-hot para variables nominales: no inventa orden ni parecidos. La parte (ii) es su flexibilidad: con un modelo lineal equivale al modelo de medias por celda del análisis de varianza. Sus costos son la dimensión $K$ (más de $40\,000$ columnas para los códigos postales de EE.\,UU., como señala la guía), el hecho de que las categorías raras tienen coeficientes estimados con muy pocos datos, y que una categoría nueva en producción no tiene columna (hay que decidir, por ejemplo, codificarla como el vector nulo). La matriz resultante es dispersa (una entrada no nula por fila y variable), lo que conecta con la Proposición~\ref{prop:escaladores}(iv): si se escala, debe hacerse con MaxAbs o no hacerse.

\subsection{Codificación binaria}\label{sec:binaria}

\begin{definicion}[Codificación binaria]\label{def:binaria}
\noindent Se numeran las categorías con $\iota:\mathcal{C}\to\{1,\dots,K\}$ biyectiva y se representa $\iota(c)$ en base $2$ con $d$ dígitos: $\phi(c)=(b_{d-1},\dots,b_0)\in\{0,1\}^d$ con $\iota(c)=\sum_{s=0}^{d-1}b_s2^s$, donde $d$ es el menor entero con $2^d>K$.
\end{definicion}

\noindent La numeración empieza en $1$ en \texttt{category\_encoders} (en orden de aparición), lo que reserva el código $0\cdots0$; por eso $8$ categorías producen $4$ columnas en el ejemplo de la guía y no $3$.

\begin{proposicion}[Dimensión e inyectividad]\label{prop:binaria}
\noindent (i) $d=\lfloor\log_2K\rfloor+1=\lceil\log_2(K+1)\rceil$. (ii) $\phi$ es inyectiva: la codificación no pierde información.
\end{proposicion}
\begin{proof}
\noindent (i) Sea $d$ el entero con $2^{d-1}\le K<2^d$; es el menor con $2^d>K$ y cumple $d-1\le\log_2K<d$, luego $d=\lfloor\log_2K\rfloor+1$. De $K<2^d$ se sigue $K+1\le2^d$, y de $2^{d-1}\le K$ se sigue $K+1>2^{d-1}$; así $d-1<\log_2(K+1)\le d$ y $d=\lceil\log_2(K+1)\rceil$. (ii) Todo entero en $\{0,\dots,2^d-1\}$ tiene una única representación binaria con $d$ dígitos (existencia por divisiones sucesivas entre $2$; unicidad porque si dos representaciones distintas difieren por primera vez, desde arriba, en el dígito $s$, sus valores difieren al menos en $2^s-\sum_{t<s}2^t=1$). Como $\iota$ es biyectiva, $\phi$ es composición de inyectivas.
\end{proof}

\noindent La guía afirma que la codificación binaria ``puede provocar una pérdida de información''. Por (ii) no se pierde información en sentido estricto. Lo que ocurre es de otra naturaleza y depende del modelo.

\begin{proposicion}[Capacidad de un modelo lineal sobre bits]\label{prop:binaria-lineal}
\noindent Con un modelo lineal $f(c)=\beta_0+\sum_{s}\beta_sb_s(c)$, los vectores de efectos representables forman un subespacio de $\R^K$ de dimensión a lo sumo $d+1$. Si $K>d+1$, lo que ocurre para todo $K\ge5$, no todo vector de efectos es representable. Además, la distancia de Hamming entre códigos, $\#\{s:b_s(c)\ne b_s(c')\}$, varía entre $1$ y $d$ según la numeración, sin relación con las categorías.
\end{proposicion}
\begin{proof}
\noindent La aplicación $(\beta_0,\dots,\beta_{d-1})\mapsto(f(c_1),\dots,f(c_K))$ es lineal de $\R^{d+1}$ en $\R^K$, y su imagen tiene dimensión a lo sumo $d+1$. Por la Proposición~\ref{prop:binaria}, $d+1\le\log_2K+2$. La función $K\mapsto K-\log_2K-2$ vale $5-\log_25-2\approx0.68>0$ en $K=5$ y es creciente para $K\ge2$ (su derivada es $1-\frac1{K\ln2}>0$), así que $d+1<K$ para todo $K\ge5$ y la imagen es un subespacio propio. (En $K=4$ se tiene $d=3$ y $d+1=K$.) En el ejemplo de la guía, $K=8$ y $d+1=5$. Para las distancias: en el ejemplo, Red ($\iota=3$, \texttt{0011}) y Pink ($\iota=7$, \texttt{0111}) difieren en un bit, mientras que White ($\iota=1$, \texttt{0001}) y Brown ($\iota=8$, \texttt{1000}) difieren en dos; cambiar la numeración cambia estas distancias.
\end{proof}

\noindent Esto es lo que la guía llama ``difuminar las distinciones'': un modelo lineal o basado en distancias ve parecidos entre categorías que comparten bits y no puede asignar efectos arbitrarios. Un árbol, en cambio, puede aislar cualquier categoría fijando sus $d$ bits con a lo sumo $d$ cortes, porque cada corte sobre una columna binaria separa según un bit.

\subsection{Hashing de características}\label{sec:hashing}

\noindent \emph{Motivación.} En aplicaciones de gran escala (filtros de spam personalizados, publicidad en línea), el número de categorías es enorme y crece con el tiempo: palabras, identificadores de usuario, combinaciones usuario--palabra. One-hot exige un diccionario $\mathcal{C}\to\{1,\dots,K\}$ que no cabe en memoria, que hay que reconstruir cuando aparecen categorías nuevas y cuya dimensión $K$ no se conoce de antemano. Weinberger, Dasgupta, Langford, Smola y Attenberg (2009) formalizaron el ``truco del hashing'': reemplazar el diccionario por una función fija que envía cada categoría a una de $m$ columnas, aceptando colisiones, y mostraron que con un signo aleatorio los productos internos se conservan en esperanza.

\begin{definicion}[Mapa de hashing con signo]\label{def:hash}
\noindent Sea $\mathcal{U}$ el universo de categorías (posiblemente infinito), $h:\mathcal{U}\to\{1,\dots,m\}$ y $\xi:\mathcal{U}\to\{-1,+1\}$. Para un vector $x\in\R^{\mathcal U}$ con finitas coordenadas no nulas (por ejemplo, conteos de palabras o el indicador de una categoría), el vector \emph{hasheado} $\phi(x)\in\R^m$ es
\[
\phi_j(x)=\sum_{i\in\mathcal{U}:\,h(i)=j}\xi(i)\,x_i,\qquad j=1,\dots,m.
\]
Para una única categoría $c$ ($x=e_c$), $\phi(e_c)=\xi(c)\,e_{h(c)}$. Dos categorías $c\ne c'$ \emph{colisionan} si $h(c)=h(c')$.
\end{definicion}

\noindent En la implementación, $h$ y $\xi$ son funciones hash deterministas (\texttt{FeatureHasher} usa MurmurHash3: el índice sale del valor absoluto del hash módulo $m$ y el signo, del signo de ese mismo valor). Para analizarlas se modelan como aleatorias: $h$ uniforme e independiente por pares, y los $\xi(i)$ independientes, con $\Prob(\xi(i)=\pm1)=\frac12$ e independientes de $h$. Nada se ajusta con los datos, y por eso no hay diccionario que guardar ni riesgo de fuga en esta etapa.

\begin{proposicion}[Colisiones]\label{prop:colisiones}
\noindent Sean $K$ categorías distintas. (i) Si $K>m$, hay al menos una colisión (principio del palomar). (ii) Bajo el modelo aleatorio, dos categorías fijas colisionan con probabilidad $1/m$, el número esperado de pares que colisionan es $\binom K2/m$ y, si $h$ es además totalmente independiente, la probabilidad de no tener colisiones es $\prod_{i=1}^{K-1}(1-i/m)\le\exp\!\big(-\tfrac{K(K-1)}{2m}\big)$.
\end{proposicion}
\begin{proof}
\noindent (i) $K$ objetos en $m<K$ cajas no pueden quedar uno por caja. (ii) $\Prob(h(c)=h(c'))=\sum_j\Prob(h(c)=j)\Prob(h(c')=j)=m\cdot m^{-2}=1/m$ por independencia por pares y uniformidad. El número de pares que colisionan es $\sum_{\{c,c'\}}\ind{h(c)=h(c')}$ y, por linealidad de la esperanza, su media es $\binom K2/m$. Con independencia total, la probabilidad de que la categoría $i+1$ evite las $i$ cajas ya ocupadas por las anteriores (sin colisiones previas) es $1-i/m$, y el producto de estas probabilidades condicionales da la probabilidad de ninguna colisión. Por $1-t\le e^{-t}$, el producto es $\le\exp(-\sum_{i=1}^{K-1}i/m)=\exp(-K(K-1)/(2m))$.
\end{proof}

\noindent En el ejemplo de la guía, $K=8>m=3$ y las colisiones son inevitables; la salida real muestra White y Blue con el mismo vector, y Green, Pink y Brown también. Con $K=10^4$ categorías y $m=2^{20}\approx10^6$ columnas, el número esperado de pares en colisión es $\binom{10^4}2/2^{20}\approx47.7$: casi seguro hay alguna, pero afectan a una fracción ínfima de las categorías. El signo $\xi$ hace que las colisiones no sesguen la geometría.

\begin{teorema}[Weinberger et al.: insesgadez del producto interno]\label{teo:hash}
\noindent Bajo el modelo aleatorio, para $x,x'\in\R^{\mathcal U}$ con soporte finito,
\[
\E\big[\ip{\phi(x)}{\phi(x')}\big]=\ip{x}{x'},\qquad \Var\big(\ip{\phi(x)}{\phi(x')}\big)=\frac1m\sum_{i\ne j}\big(x_i^2x_j'^2+x_ix_i'x_jx_j'\big).
\]
\end{teorema}
\begin{proof}
\noindent Escribimos $\delta_{ij}=\ind{h(i)=h(j)}$. Desarrollando el producto interno de las sumas que definen $\phi$,
\[
\ip{\phi(x)}{\phi(x')}=\sum_{j=1}^m\sum_{h(i)=j}\sum_{h(k)=j}\xi(i)\xi(k)x_ix'_k=\sum_{i,k}\xi(i)\xi(k)\,\delta_{ik}\,x_ix'_k.
\]
Los términos con $i=k$ valen $\xi(i)^2x_ix'_i=x_ix'_i$, porque $\xi(i)^2=1$ y $\delta_{ii}=1$; su suma es $\ip x{x'}$. Sea $S=\sum_{i\ne k}\xi(i)\xi(k)\delta_{ik}x_ix'_k$ el resto. Para $i\ne k$, $\E[\xi(i)\xi(k)\delta_{ik}]=\E[\xi(i)]\E[\xi(k)]\E[\delta_{ik}]=0$ por independencia, así que $\E S=0$ y el producto interno es insesgado.

\medskip

\noindent Su varianza es $\E[S^2]=\sum_{i\ne k}\sum_{a\ne b}\E[\xi(i)\xi(k)\xi(a)\xi(b)]\,\E[\delta_{ik}\delta_{ab}]\,x_ix'_kx_ax'_b$. Como los $\xi$ son independientes con media $0$ y $\xi^2=1$, $\E[\xi(i)\xi(k)\xi(a)\xi(b)]$ es $1$ si cada índice aparece un número par de veces y $0$ en otro caso; con $i\ne k$ y $a\ne b$ esto ocurre exactamente cuando $(a,b)=(i,k)$ o $(a,b)=(k,i)$. En ambos casos $\E[\delta_{ik}\delta_{ab}]=\E[\delta_{ik}^2]=\E[\delta_{ik}]=1/m$ (Proposición~\ref{prop:colisiones}). Sumando las dos contribuciones, $\E[S^2]=\frac1m\sum_{i\ne k}\big(x_ix'_kx_ix'_k+x_ix'_kx_kx'_i\big)=\frac1m\sum_{i\ne k}\big(x_i^2x_k'^2+x_ix'_ix_kx'_k\big)$.
\end{proof}

\noindent Para vectores unitarios la varianza es $O(1/m)$: el producto interno hasheado se concentra en torno al verdadero al crecer $m$. Sin el signo ($\xi\equiv1$), el término $S$ tendría esperanza $\frac1m\sum_{i\ne k}x_ix'_k$, que para vectores de conteos no negativos es positiva: las colisiones harían parecer más similares a todos los pares. Para dos categorías distintas, $\ip{\phi(e_c)}{\phi(e_{c'})}=\xi(c)\xi(c')\delta_{cc'}$ vale $0$ si no colisionan y $\pm1$ con igual probabilidad si colisionan, con media $0=\ip{e_c}{e_{c'}}$.

\subsection{Contraste de las codificaciones}\label{sec:contraste-cat}

\begin{center}
\begin{tabularx}{\linewidth}{@{}lcccY@{}}
\toprule
Codificación & Dimensión & Inyectiva & Diccionario & Geometría inducida y uso\\
\midrule
Ordinal & $1$ & sí & sí & orden y distancias arbitrarios si la variable es nominal (Prop.~\ref{prop:ordinal-lineal}); adecuada para ordinales y para árboles (Prop.~\ref{prop:ordinal-arbol})\\
One-hot & $K$ & sí & sí & equidistancia, efectos arbitrarios (Prop.~\ref{prop:onehot}); inviable con $K$ muy grande; categorías nuevas sin columna\\
Binaria & $\lceil\log_2(K+1)\rceil$ & sí & sí & parecidos artificiales por bits compartidos; capacidad lineal limitada (Prop.~\ref{prop:binaria-lineal}); adecuada para árboles\\
Hashing & $m$ fija & no si $K>m$ & no & colisiones indistinguibles; productos internos insesgados (Teo.~\ref{teo:hash}); memoria fija, categorías nuevas sin reentrenar, no interpretable\\
\bottomrule
\end{tabularx}
\end{center}

\noindent La regla de la guía se deduce de la tabla: ordinal para variables ordinales y modelos de árboles; one-hot para cardinalidad pequeña y modelos no basados en árboles; binaria para reducir la dimensión de one-hot cuando el modelo tolera los parecidos artificiales; hashing para alta cardinalidad y recursos limitados. PCA después de one-hot, que la guía menciona como alternativa, reduce la dimensión pero, por el Teorema~\ref{teo:pca}, conserva las direcciones de mayor varianza, que para indicadores son las de las categorías frecuentes, y pierde la interpretación por categoría.

\section{Series de tiempo}\label{sec:series}

\noindent \emph{Motivación.} En una serie de tiempo $(y_t)_{t\in\mathbb Z}$ las observaciones no son intercambiables: $y_t$ depende de $y_{t-1},y_{t-2},\dots$, y el orden es parte de la información. Un modelo tabular (regresión, árboles, redes) no sabe nada de ese orden; para que lo aproveche hay que escribir el pasado como columnas. Al mismo tiempo, el orden introduce una restricción nueva que no existe en datos i.i.d.: al predecir $y_t$ solo puede usarse lo que se conocía antes de $t$. Las técnicas de la guía (descomponer la fecha, codificarla, rezagos y ventanas móviles) son formas de construir columnas que respeten esa restricción.

\begin{definicion}[Característica admisible]\label{def:admisible}
\noindent Sea $\mathcal{I}_{t}$ la información disponible al final del instante $t$: los valores $y_s$ y las covariables $w_s$ con $s\le t$, más las funciones deterministas del calendario (fecha, día de la semana, feriados), que se conocen de antemano. Una característica para predecir $y_t$ con horizonte $1$ es \emph{admisible} si es función de $\mathcal{I}_{t-1}$. Toda característica no admisible es fuga del objetivo.
\end{definicion}

\noindent La definición convierte la advertencia de la guía (``un rezago solo puede usar valores anteriores al momento de la predicción'') en un criterio verificable columna por columna. Es la misma condición que garantizan las consultas \emph{point-in-time} de Feature Store: el valor de una característica para un evento con marca de tiempo $t$ se toma del último registro con marca $\le t$.

\subsection{Descomposición de la fecha y codificación cíclica}\label{sec:fecha}

\noindent Las características de calendario (mes, día del mes, semana del año, día del año, trimestre, día de la semana) son funciones deterministas de la marca de tiempo y, por la Definición~\ref{def:admisible}, siempre admisibles. Su utilidad es permitir que el modelo capture patrones periódicos; su dificultad es que son periódicas, y ninguna codificación en una recta representa bien un ciclo.

\begin{definicion}[Codificación cíclica]\label{def:ciclica}
\noindent Para una característica periódica $m\in\{0,1,\dots,P-1\}$ de periodo $P$ (mes con $P=12$ numerado desde $0$, hora con $P=24$, día de la semana con $P=7$), la \emph{codificación cíclica} es $\psi(m)=\big(\cos\frac{2\pi m}{P},\,\sin\frac{2\pi m}{P}\big)\in\R^2$.
\end{definicion}

\noindent La codificación coloca los $P$ valores en los vértices de un polígono regular inscrito en la circunferencia unidad. La propiedad que se busca es que la distancia entre dos valores dependa solo de cuántos pasos los separan en el ciclo.

\begin{proposicion}[Distancia circular]\label{prop:ciclica}
\noindent Sea $d_P(m,m')=\min(\abs{m-m'},P-\abs{m-m'})$ la distancia circular. Entonces
\[
\norm{\psi(m)-\psi(m')}=2\sin\Big(\frac{\pi\,d_P(m,m')}{P}\Big),
\]
que es una función estrictamente creciente de $d_P(m,m')\in\{0,\dots,\lfloor P/2\rfloor\}$. En cambio, con la codificación ordinal $m\mapsto m$, la distancia es $\abs{m-m'}$.
\end{proposicion}
\begin{proof}
\noindent Sea $\alpha=2\pi(m-m')/P$. Por la identidad $\cos a\cos b+\sin a\sin b=\cos(a-b)$,
\[
\norm{\psi(m)-\psi(m')}^2=2-2\Big(\cos\tfrac{2\pi m}{P}\cos\tfrac{2\pi m'}{P}+\sin\tfrac{2\pi m}{P}\sin\tfrac{2\pi m'}{P}\Big)=2-2\cos\alpha=4\sin^2\frac\alpha2,
\]
usando $1-\cos\alpha=2\sin^2(\alpha/2)$. Entonces $\norm{\psi(m)-\psi(m')}=2\abs{\sin(\pi(m-m')/P)}$. La función $k\mapsto\abs{\sin(\pi k/P)}$ es par y cumple $\abs{\sin(\pi(P-k)/P)}=\abs{\sin(\pi-\pi k/P)}=\abs{\sin(\pi k/P)}$, así que su valor en $k=\abs{m-m'}$ coincide con su valor en $d_P(m,m')$. Para $d\in[0,P/2]$ el argumento $\pi d/P$ está en $[0,\pi/2]$, donde el seno es no negativo y estrictamente creciente.
\end{proof}

\noindent Con meses numerados $0$ (enero) a $11$ (diciembre), la codificación ordinal pone a diciembre y enero a distancia $11$, la máxima, cuando en el ciclo son vecinos; la cíclica les da la distancia $2\sin(\pi/12)\approx0.518$, la misma que a enero y febrero. Un segundo argumento, para modelos lineales, muestra qué forma de efecto estacional permite la codificación.

\begin{proposicion}[Efecto lineal de la codificación cíclica]\label{prop:armonico}
\noindent Un modelo lineal $f(m)=\beta_0+a\cos\frac{2\pi m}P+b\sin\frac{2\pi m}P$ es exactamente una sinusoide de periodo $P$ con amplitud y fase arbitrarias: $f(m)=\beta_0+A\cos\big(\frac{2\pi m}{P}-\varphi\big)$ con $A=\sqrt{a^2+b^2}$ y $(\cos\varphi,\sin\varphi)=(a,b)/A$ si $A>0$.
\end{proposicion}
\begin{proof}
\noindent Por la fórmula del coseno de una diferencia, $A\cos(\theta-\varphi)=A\cos\varphi\cos\theta+A\sin\varphi\sin\theta=a\cos\theta+b\sin\theta$ con $\theta=2\pi m/P$.
\end{proof}

\noindent \emph{Contraste.} La codificación cíclica usa $2$ columnas y fuerza un efecto sinusoidal suave con un solo máximo en el ciclo; la codificación one-hot del mes (o del trimestre, como sugiere la guía) usa $P$ columnas y, por la Proposición~\ref{prop:onehot}, permite cualquier perfil estacional, a costa de estimar $P$ parámetros; la ordinal usa $1$ columna e introduce la discontinuidad diciembre--enero. Esto explica la recomendación de Data Wrangler: codificación cíclica para modelos lineales y redes, ordinal para árboles, a los que la discontinuidad no afecta porque pueden cortar en cualquier punto (Proposición~\ref{prop:ordinal-arbol}). Si hace falta un perfil más rico que una sinusoide, se añaden armónicos $\cos\frac{2\pi km}{P},\sin\frac{2\pi km}{P}$ para $k=2,3,\dots$, que es un desarrollo de Fourier del efecto estacional.

\subsection{Rezagos y autocorrelación}\label{sec:rezagos}

\begin{definicion}[Rezago y autocorrelación]\label{def:rezago}
\noindent La característica de \emph{rezago} $k\ge1$ para predecir $y_t$ es $y_{t-k}$; un rezago \emph{estacional} es $y_{t-s}$ con $s$ el periodo (por ejemplo, la misma hora del día anterior, $s=24$). Si $(y_t)$ es estacionaria en sentido débil (media $\mu$ y covarianzas $\gamma(h)=\Cov(y_t,y_{t-h})$ que no dependen de $t$), su \emph{función de autocorrelación} es $\rho(h)=\gamma(h)/\gamma(0)$.
\end{definicion}

\noindent Los rezagos son admisibles por construcción, porque $y_{t-k}\in\mathcal{I}_{t-1}$. La autocorrelación mide cuánto informa cada rezago de forma lineal, como muestra el resultado siguiente.

\begin{proposicion}[Mejor predictor lineal con un rezago]\label{prop:lag1}
\noindent Si $(y_t)$ es estacionaria débil con media $0$ y $\gamma(0)>0$, entonces $\E[(y_t-ay_{t-k})^2]$ se minimiza en $a=\rho(k)$, con valor mínimo $\gamma(0)(1-\rho(k)^2)$.
\end{proposicion}
\begin{proof}
\noindent Desarrollando, $\E[(y_t-ay_{t-k})^2]=\gamma(0)-2a\gamma(k)+a^2\gamma(0)$, una parábola en $a$ con coeficiente principal $\gamma(0)>0$, minimizada en $a=\gamma(k)/\gamma(0)=\rho(k)$. Sustituyendo, el valor es $\gamma(0)-2\gamma(k)^2/\gamma(0)+\gamma(k)^2/\gamma(0)=\gamma(0)(1-\rho(k)^2)$.
\end{proof}

\noindent Así, la fracción de varianza que un rezago explica linealmente es $\rho(k)^2$, y los rezagos útiles se identifican en el gráfico de autocorrelación. En el modelo más simple la estructura es explícita.

\begin{ejemplo}[Proceso AR(1)]\label{ej:ar1}
\noindent Sea $y_t=\phi y_{t-1}+\varepsilon_t$ con $\abs\phi<1$ y $\varepsilon_t$ i.i.d.\ de media $0$ y varianza $\sigma^2$, independientes del pasado $y_{t-1},y_{t-2},\dots$, y supongamos la solución estacionaria. Para $h\ge1$, $\gamma(h)=\Cov(\phi y_{t-1}+\varepsilon_t,\,y_{t-h})=\phi\gamma(h-1)+0$, porque $\varepsilon_t$ es independiente de $y_{t-h}$. Iterando, $\gamma(h)=\phi^h\gamma(0)$ y $\rho(h)=\phi^h$. De $\gamma(0)=\Var(\phi y_{t-1}+\varepsilon_t)=\phi^2\gamma(0)+\sigma^2$ se obtiene $\gamma(0)=\sigma^2/(1-\phi^2)$. Obsérvese que $\rho(2)=\phi^2\ne0$ si $\phi\ne0$: el rezago $2$ está correlacionado con $y_t$. La pregunta relevante para construir características es otra: si aporta algo \emph{una vez incluido} el rezago $1$.
\end{ejemplo}

\begin{definicion}[Autocorrelación parcial de orden $2$]\label{def:pacf}
\noindent Sea $(y_t)$ estacionaria débil con media $0$ y tal que la matriz de covarianzas $\Gamma=\begin{psmallmatrix}\gamma(0)&\gamma(1)\\\gamma(1)&\gamma(0)\end{psmallmatrix}$ de $(y_{t-1},y_{t-2})$ es invertible. Sean $(a_1^*,a_2^*)$ los coeficientes que minimizan $\E[(y_t-a_1y_{t-1}-a_2y_{t-2})^2]$. La \emph{autocorrelación parcial de orden $2$} es $\alpha(2)=a_2^*$: el peso del rezago $2$ en el mejor predictor lineal que ya usa el rezago $1$.
\end{definicion}

\noindent La función a minimizar es una cuadrática en $(a_1,a_2)$ con matriz hessiana $2\Gamma$, definida positiva por ser una matriz de covarianzas invertible; su minimizador es único y está caracterizado por anular el gradiente, es decir, por las condiciones de ortogonalidad $\E[(y_t-a_1y_{t-1}-a_2y_{t-2})\,y_{t-k}]=0$ para $k=1,2$ (la derivada respecto de $a_k$ es $-2$ veces esa esperanza).

\begin{proposicion}[El segundo rezago no aporta en un AR(1)]\label{prop:pacf-ar1}
\noindent En el proceso AR(1) del Ejemplo~\ref{ej:ar1}, $(a_1^*,a_2^*)=(\phi,0)$; en particular $\alpha(2)=0$.
\end{proposicion}
\begin{proof}
\noindent $\Gamma$ es invertible porque su determinante es $\gamma(0)^2-\gamma(1)^2=\gamma(0)^2(1-\phi^2)>0$, usando $\gamma(1)=\phi\gamma(0)$ y $\abs\phi<1$. Con $(a_1,a_2)=(\phi,0)$, el residuo es $y_t-\phi y_{t-1}=\varepsilon_t$, que es independiente de $y_{t-1}$ y de $y_{t-2}$ y tiene media $0$, así que $\E[\varepsilon_ty_{t-k}]=\E[\varepsilon_t]\E[y_{t-k}]=0$ para $k=1,2$. Se cumplen las condiciones de ortogonalidad, que caracterizan al minimizador único.
\end{proof}

\noindent Así, en un AR(1) la característica $y_{t-2}$ es redundante dada $y_{t-1}$, aunque su correlación con $y_t$ no sea nula: la autocorrelación simple no basta para decidir qué rezagos incluir, y la parcial sí. Es un caso particular del fenómeno de redundancia de la Sección~\ref{sec:redundancia}.

\noindent \emph{Fuga en rezagos.} Un rezago mal alineado ($k=0$) usa el propio $y_t$ y produce un ajuste perfecto que no existe en producción. Lo mismo ocurre, de forma más sutil, con estadísticas calculadas en ventanas centradas o con toda la serie (una media global, una normalización con el máximo de toda la serie), porque incluyen valores posteriores a $t$ y violan la Definición~\ref{def:admisible}.

\subsection{Ventanas móviles}\label{sec:ventanas}

\begin{definicion}[Estadísticas de ventana móvil]\label{def:ventana}
\noindent Para una ventana de tamaño $w$, las características de ventana móvil de $y_t$ son estadísticas de $(y_{t-w},\dots,y_{t-1})$: la media $\bar y_t^{(w)}=\frac1w\sum_{k=1}^wy_{t-k}$, el mínimo, el máximo, la desviación estándar, etc.
\end{definicion}

\noindent La ventana termina en $t-1$ para que la característica sea admisible (con $w=3$, la fila $t$ recibe estadísticas de $t-3,t-2,t-1$, como describe la guía). El tamaño $w$ regula un compromiso entre ruido y retraso.

\begin{proposicion}[Compromiso de la media móvil]\label{prop:media-movil}
\noindent (i) Si $y_t=\mu+\varepsilon_t$ con $\varepsilon_t$ incorrelacionados de varianza $\sigma^2$, entonces $\E\bar y_t^{(w)}=\mu$ y $\Var\bar y_t^{(w)}=\sigma^2/w$. (ii) Si $y_t=a+bt$ (tendencia lineal sin ruido), entonces $\bar y_t^{(w)}=y_t-b\,\frac{w+1}2$: la media móvil va retrasada $\frac{w+1}{2}$ periodos.
\end{proposicion}
\begin{proof}
\noindent (i) Por linealidad, $\E\bar y_t^{(w)}=\mu$, y por incorrelación $\Var\bar y_t^{(w)}=w^{-2}\sum_{k=1}^w\Var\varepsilon_{t-k}=\sigma^2/w$. (ii) $\bar y_t^{(w)}=a+\frac bw\sum_{k=1}^w(t-k)=a+bt-\frac bw\cdot\frac{w(w+1)}2=y_t-b\frac{w+1}2$.
\end{proof}

\noindent Ventanas grandes reducen el ruido en proporción $1/w$ pero reaccionan tarde a cambios de nivel; ventanas pequeñas siguen los cambios pero son ruidosas. Por eso es habitual incluir varias ventanas a la vez. La extracción automática (Data Wrangler ofrece \emph{Rolling window features} y \emph{Extract features}) calcula muchas estadísticas de este tipo y deja al modelo o a una selección posterior (Sección~\ref{sec:filtros}) decidir cuáles sirven.

\section{Imágenes}\label{sec:imagenes}

\noindent Esta sección construye con precisión los objetos que la guía menciona de pasada (píxeles, normalización, bordes, redes convolucionales, vectores de características, cuadros delimitadores) y demuestra las propiedades en que se apoya cada técnica.

\subsection{La imagen digital como objeto matemático}\label{sec:imagen-digital}

\noindent \emph{Motivación.} Una fotografía es, físicamente, una función continua que asigna a cada punto de un rectángulo una intensidad de luz en varias bandas de color. Un sensor digital no puede registrar infinitos valores: mide la intensidad solo en una rejilla finita de posiciones (\emph{muestreo}) y redondea cada medición a uno de finitos niveles (\emph{cuantización}). El resultado es un objeto finito, y es ese objeto, no la escena, el que llega al algoritmo.

\begin{definicion}[Imagen digital, píxel, canal]\label{def:imagen}
\noindent Sean $H,W,C$ enteros positivos y $b$ un entero positivo (la \emph{profundidad} en bits). Sea $\Omega=\{0,\dots,H-1\}\times\{0,\dots,W-1\}$ la \emph{rejilla}. Una \emph{imagen digital} de tamaño $H\times W$ con $C$ canales es una función
\[
I:\Omega\times\{1,\dots,C\}\to\{0,1,\dots,2^b-1\}.
\]
Un \emph{píxel} es un punto $(u,v)\in\Omega$ ($u$ es la fila, $v$ la columna), y su valor es el vector de sus $C$ intensidades, $\big(I(u,v,c)\big)_{c=1}^{C}$. El \emph{canal} $c$ es la función $(u,v)\mapsto I(u,v,c)$ definida en $\Omega$. Cuando se permiten valores reales (tras normalizar o dentro de una red), $I$ toma valores en $\R$, y el conjunto de tales imágenes es el espacio vectorial $\R^{H\times W\times C}$ con las operaciones punto a punto.
\end{definicion}

\noindent \emph{Alcance.} Con $C=1$ se tiene una imagen en escala de grises; con $C=3$ y canales rojo, verde y azul, una imagen RGB; las imágenes multiespectrales de drones del documento original tienen $C>3$. Con $b=8$, los valores van de $0$ a $255$, el rango que cita la guía. Un \emph{tensor} de orden $3$ es simplemente un elemento de $\R^{H\times W\times C}$, es decir, un arreglo con tres índices. La definición fija $H$, $W$ y $C$: dos imágenes de tamaños distintos pertenecen a espacios distintos y no pueden compararse ni entrar al mismo modelo sin antes llevarlas a un tamaño común (Sección~\ref{sec:resolucion}).

\medskip

\noindent Para usar una imagen en un modelo tabular hay que escribirla como vector de $\R^p$. La forma estándar es recorrer los índices en un orden fijo.

\begin{proposicion}[Vectorización]\label{prop:vec}
\noindent La aplicación $\mathrm{vec}:\R^{H\times W\times C}\to\R^{HWC}$ definida por $\mathrm{vec}(I)_{k}=I(u,v,c+1)$ con $k=(uW+v)C+c$, para $u\in\{0,\dots,H-1\}$, $v\in\{0,\dots,W-1\}$, $c\in\{0,\dots,C-1\}$, es una biyección lineal que conserva la norma: $\norm{\mathrm{vec}(I)}^2=\sum_{u,v,c}I(u,v,c)^2$.
\end{proposicion}
\begin{proof}
\noindent Basta ver que $(u,v,c)\mapsto k=(uW+v)C+c$ es una biyección de $\{0,\dots,H-1\}\times\{0,\dots,W-1\}\times\{0,\dots,C-1\}$ sobre $\{0,\dots,HWC-1\}$. Dado $k$ en ese rango, el algoritmo de la división da de forma única $c=k\bmod C$ y $q=\lfloor k/C\rfloor\in\{0,\dots,HW-1\}$; de nuevo de forma única, $v=q\bmod W$ y $u=\lfloor q/W\rfloor\in\{0,\dots,H-1\}$. Recíprocamente, esos $(u,v,c)$ cumplen $(uW+v)C+c=qC+c=k$. Así, $\mathrm{vec}$ solo reordena coordenadas: es lineal, y la suma de cuadrados no cambia al reordenar sumandos.
\end{proof}

\noindent La dimensión es $p=HWC$: con $224\times224\times3$ son $150\,528$ coordenadas. El resultado siguiente muestra que, además de ser muchas, estas coordenadas tienen una geometría inadecuada para comparar imágenes. Para una imagen de un canal y $(s,t)\in\mathbb{Z}^2$, la \emph{traslación circular} es $(\tau_{s,t}I)(u,v)=I\big((u-s)\bmod H,\,(v-t)\bmod W\big)$: desplaza el contenido $s$ filas hacia abajo y $t$ columnas hacia la derecha, y lo que sale por un borde reaparece por el opuesto.

\begin{proposicion}[La distancia entre píxeles no percibe desplazamientos]\label{prop:dist-pixeles}
\noindent Sea $\delta\in\R^{H\times W}$ la imagen de un canal que vale $1$ en el píxel $(0,0)$ y $0$ en los demás. Para toda traslación $(s,t)$ con $(s\bmod H,t\bmod W)\ne(0,0)$, se tiene $\norm{\delta-\tau_{s,t}\delta}^2=2$.
\end{proposicion}
\begin{proof}
\noindent $\tau_{s,t}\delta$ vale $1$ solo en el píxel $(s\bmod H,t\bmod W)$, que es distinto de $(0,0)$ por hipótesis. La diferencia $\delta-\tau_{s,t}\delta$ vale $1$ en $(0,0)$, $-1$ en ese otro píxel y $0$ en los demás, así que su norma al cuadrado es $1+1=2$.
\end{proof}

\noindent Un punto brillante desplazado un píxel está exactamente igual de lejos del original que uno desplazado al extremo opuesto de la imagen. La distancia euclídea entre vectores de píxeles, que usan $k$ vecinos o $k$-medias, no tiene noción de ``el mismo objeto un poco movido''. Las características de las subsecciones siguientes se construyen para remediarlo.

\subsection{Normalización de píxeles}\label{sec:pixeles}

\begin{definicion}[Normalización y estandarización por canal]\label{def:norm-pix}
\noindent La \emph{normalización a $[0,1]$} de una imagen de profundidad $b$ es $I\mapsto I/(2^b-1)$, aplicada a cada valor. La \emph{estandarización por canal} con parámetros $\mu_c\in\R$ y $\sigma_c>0$ es $I(u,v,c)\mapsto\big(I(u,v,c)/(2^b-1)-\mu_c\big)/\sigma_c$.
\end{definicion}

\begin{proposicion}[Propiedades de la normalización de píxeles]\label{prop:norm-pix}
\noindent (i) La normalización a $[0,1]$ es el escalador MinMax de la Definición~\ref{def:afin} con $c=0$ y $s=2^b-1$ fijados de antemano; su imagen está contenida en $[0,1]$ para toda imagen posible, y no es una transformación ajustada (no depende de los datos), por lo que no puede producir fuga en el sentido de la Definición~\ref{def:fuga}. (ii) Si $\mu_c$ y $\sigma_c$ son la media y la desviación estándar (poblacional) de los valores normalizados del canal $c$ sobre todos los píxeles de $\Dtr$, el canal $c$ estandarizado tiene media $0$ y desviación estándar $1$ sobre $\Dtr$. (iii) Si una red se entrenó con los parámetros $(\mu_c,\sigma_c)$ y se le aplican otros $(\mu'_c,\sigma'_c)$, la entrada que recibe es la imagen correcta sometida a la transformación afín $x\mapsto\frac{\sigma_c}{\sigma'_c}x+\frac{\mu_c-\mu'_c}{\sigma'_c}$ en cada canal.
\end{proposicion}
\begin{proof}
\noindent (i) Como $0\le I(u,v,c)\le2^b-1$, se tiene $0\le I(u,v,c)/(2^b-1)\le1$; los parámetros $c=0$ y $s=2^b-1$ son constantes. (ii) Es la Proposición~\ref{prop:escaladores}(i) aplicada a la lista de todos los valores normalizados del canal $c$ en $\Dtr$. (iii) Sea $x$ un valor normalizado del canal $c$ y $z=(x-\mu_c)/\sigma_c$ la entrada esperada por la red. Entonces $x=\sigma_cz+\mu_c$, y la entrada recibida es $(x-\mu'_c)/\sigma'_c=\frac{\sigma_c}{\sigma'_c}z+\frac{\mu_c-\mu'_c}{\sigma'_c}$.
\end{proof}

\noindent La parte (i) explica por qué la normalización de imágenes no tiene la debilidad de MinMax frente a atípicos (Proposición~\ref{prop:escaladores}(ii)): los extremos no se estiman, se conocen. La parte (iii) es la razón de usar, con un modelo preentrenado, las medias y desviaciones del conjunto con que se preentrenó y no las de los datos propios: cualquier otra elección altera sistemáticamente las entradas de la red.

\subsection{Características diseñadas a mano: bordes}\label{sec:bordes}

\noindent \emph{Motivación.} Antes de las redes neuronales, la visión por computadora construía características explícitas con conocimiento del problema. Las más básicas son los bordes, porque el contorno de un objeto suele coincidir con un cambio brusco de intensidad. Formalmente, si la intensidad fuera una función derivable $f(u,v)$ en el plano, un borde sería un lugar donde la norma de su gradiente $\nabla f=(\partial_uf,\partial_vf)$ es grande. La imagen digital solo da valores en la rejilla, así que las derivadas deben aproximarse con diferencias.

\begin{definicion}[Diferencia central]\label{def:dif-central}
\noindent Para una función $f$ de una variable y $h>0$, la \emph{diferencia central} es $D_hf(u)=\dfrac{f(u+h)-f(u-h)}{2h}$.
\end{definicion}

\begin{proposicion}[Error de la diferencia central]\label{prop:dif-central}
\noindent Si $f$ tiene tercera derivada continua en $[u-h,u+h]$, entonces
\[
D_hf(u)=f'(u)+\frac{h^2}{12}\big(f'''(\xi_+)+f'''(\xi_-)\big)
\]
para ciertos $\xi_+\in(u,u+h)$ y $\xi_-\in(u-h,u)$. En particular, $\abs{D_hf(u)-f'(u)}\le\frac{h^2}{6}\max_{[u-h,u+h]}\abs{f'''}$, y $D_hf(u)=f'(u)$ exactamente si $f$ es un polinomio de grado a lo sumo $2$.
\end{proposicion}
\begin{proof}
\noindent Por el teorema de Taylor con resto de Lagrange, $f(u\pm h)=f(u)\pm hf'(u)+\frac{h^2}2f''(u)\pm\frac{h^3}6f'''(\xi_\pm)$. Restando, los términos de orden $0$ y $2$ se cancelan: $f(u+h)-f(u-h)=2hf'(u)+\frac{h^3}{6}\big(f'''(\xi_+)+f'''(\xi_-)\big)$. Dividiendo entre $2h$ se obtiene la fórmula. La cota se sigue de acotar cada $\abs{f'''(\xi_\pm)}$ por el máximo. Si $f$ es un polinomio de grado $\le2$, $f'''\equiv0$ y el error es nulo.
\end{proof}

\noindent La diferencia \emph{hacia adelante} $(f(u+h)-f(u))/h$ tiene error de orden $h$ (el término $\frac h2f''$ no se cancela); la central es preferible porque su error es de orden $h^2$. En una imagen $h=1$ (un píxel), y el error depende de la suavidad de la intensidad. El problema práctico es el ruido: una diferencia amplifica las fluctuaciones aleatorias de cada píxel. Sobel (1968) combinó la diferencia en una dirección con un promedio en la perpendicular.

\begin{definicion}[Filtro de Sobel]\label{def:sobel}
\noindent Sea $I$ una imagen de un canal. En un píxel interior ($1\le u\le H-2$, $1\le v\le W-2$), las respuestas de Sobel horizontal y vertical son
\[
G_v(u,v)=\sum_{a=-1}^{1}\sum_{b=-1}^{1}s(a)\,d(b)\,I(u+a,v+b),\qquad G_u(u,v)=\sum_{a=-1}^{1}\sum_{b=-1}^{1}d(a)\,s(b)\,I(u+a,v+b),
\]
con $s(-1)=s(1)=1$, $s(0)=2$ (suavizado) y $d(-1)=-1$, $d(0)=0$, $d(1)=1$ (diferencia). La \emph{magnitud del gradiente} es $G(u,v)=\sqrt{G_u(u,v)^2+G_v(u,v)^2}$. En forma matricial, el núcleo de $G_v$ es el producto exterior $s\,d^\top=\begin{psmallmatrix}-1&0&1\\-2&0&2\\-1&0&1\end{psmallmatrix}$.
\end{definicion}

\noindent La restricción a píxeles interiores evita tener que decidir qué valor usar fuera de la rejilla; esa decisión (el relleno) se trata en la Definición~\ref{def:relleno}.

\begin{proposicion}[Qué calcula el filtro de Sobel]\label{prop:sobel}
\noindent (i) Si $I(u,v)=\alpha v+\beta u+\gamma$ en un entorno $3\times3$ del píxel, entonces $G_v(u,v)=8\alpha$ y $G_u(u,v)=8\beta$: el filtro dividido entre $8$ recupera exactamente el gradiente de una intensidad afín. (ii) Si $I(u,v)=\mu+\varepsilon(u,v)$ con $\varepsilon(u,v)$ incorrelacionados de media $0$ y varianza $\sigma^2$, el promedio ponderado $\frac14\sum_as(a)I(u+a,v)$ tiene media $\mu$ y varianza $\frac38\sigma^2$.
\end{proposicion}
\begin{proof}
\noindent (i) Separando las sumas, $G_v(u,v)=\sum_as(a)\sum_bd(b)\big(\alpha(v+b)+\beta(u+a)+\gamma\big)$. Como $\sum_bd(b)=0$, los términos que no dependen de $b$ desaparecen de la suma interior, que vale $\alpha\sum_bd(b)\,b=\alpha(1+1)=2\alpha$. Queda $G_v=2\alpha\sum_as(a)=2\alpha\cdot4=8\alpha$. El cálculo para $G_u$ es el mismo con los papeles de $u$ y $v$ intercambiados. (ii) Los pesos $\frac14(1,2,1)$ suman $1$, así que la media es $\mu$; por incorrelación, la varianza es $\frac1{16}(1^2+2^2+1^2)\sigma^2=\frac6{16}\sigma^2$.
\end{proof}

\noindent La parte (i) justifica el diseño (la diferencia mide la pendiente sin sesgo en intensidades afines) y la parte (ii) el suavizado (reduce el ruido a menos de la mitad sin cambiar la media). La guía menciona otra clase de características manuales: indicadores como \texttt{windshield\_present}. Formalmente, son funciones $\phi:\R^{H\times W\times C}\to\{0,1\}$ que valen $1$ si un detector encuentra la parte buscada. Su limitación es de fondo: cada problema nuevo exige diseñar sus propios detectores, y el diseño requiere conocimiento del dominio. Las redes convolucionales automatizan ese diseño.

\subsection{Convolución}\label{sec:conv}

\noindent \emph{Motivación.} El filtro de Sobel aplica los mismos $9$ pesos en cada píxel. Esa idea (un filtro pequeño que se desliza) tiene dos virtudes: usa pocos parámetros y trata igual a todas las posiciones. El resultado principal de esta subsección (Teorema~\ref{teo:caract-conv}) muestra que la segunda virtud no es accidental: las operaciones lineales que tratan igual a todas las posiciones son \emph{exactamente} las convoluciones. En esta subsección los índices de fila y columna se toman módulo $H$ y $W$, es decir, $\Omega$ se identifica con $\mathbb{Z}_H\times\mathbb{Z}_W$.

\begin{definicion}[Convolución circular]\label{def:conv}
\noindent Sea $I\in\R^{H\times W}$ una imagen de un canal y $k:\mathbb Z_H\times\mathbb Z_W\to\R$ un \emph{núcleo}. La convolución (en la convención de las redes neuronales, que no invierte el núcleo) es
\[
(k*I)(u,v)=\sum_{(a,b)\in\mathbb Z_H\times\mathbb Z_W}k(a,b)\,I(u+a,v+b).
\]
El núcleo tiene \emph{radio} $r$ si $k(a,b)=0$ salvo cuando $a,b\in\{-r,\dots,r\}$ (módulo $H$ y $W$); entonces solo intervienen $(2r+1)^2$ pesos. El operador de traslación es el $\tau_{s,t}$ de la Proposición~\ref{prop:dist-pixeles}.
\end{definicion}

\noindent La suma sobre todo $\mathbb Z_H\times\mathbb Z_W$ incluye como caso particular los núcleos pequeños ($r=1$ da los núcleos $3\times3$, como el de Sobel). La convolución es lineal en $I$ para $k$ fijo, porque cada valor de salida es una combinación lineal de valores de entrada con coeficientes que no dependen de $I$.

\begin{teorema}[Equivariancia por traslaciones]\label{teo:equiv}
\noindent Para todo núcleo $k$ y toda traslación, $k*(\tau_{s,t}I)=\tau_{s,t}(k*I)$.
\end{teorema}
\begin{proof}
\noindent Para cada $(u,v)$,
\[
\big(k*(\tau_{s,t}I)\big)(u,v)=\sum_{a,b}k(a,b)\,I(u+a-s,v+b-t)=(k*I)(u-s,v-t)=\big(\tau_{s,t}(k*I)\big)(u,v),
\]
donde la segunda igualdad es la definición de $k*I$ evaluada en $(u-s,v-t)$.
\end{proof}

\begin{teorema}[Caracterización de la convolución]\label{teo:caract-conv}
\noindent Sea $L:\R^{H\times W}\to\R^{H\times W}$ lineal. Entonces $L$ conmuta con todas las traslaciones ($L\tau_{s,t}=\tau_{s,t}L$ para todo $(s,t)$) si y solo si existe un núcleo $k$ tal que $LI=k*I$ para toda $I$.
\end{teorema}
\begin{proof}
\noindent Si $L$ es una convolución, conmuta con las traslaciones por el Teorema~\ref{teo:equiv}. Recíprocamente, sea $\delta$ la imagen de la Proposición~\ref{prop:dist-pixeles} y $g=L\delta$. La imagen $\tau_{s,t}\delta$ vale $1$ en $(s,t)$ y $0$ en el resto, así que toda imagen se escribe como $I=\sum_{(s,t)}I(s,t)\,\tau_{s,t}\delta$. Por linealidad y conmutación,
\[
LI=\sum_{(s,t)}I(s,t)\,L\tau_{s,t}\delta=\sum_{(s,t)}I(s,t)\,\tau_{s,t}g,\quad\text{es decir,}\quad (LI)(u,v)=\sum_{(s,t)}I(s,t)\,g(u-s,v-t).
\]
Definiendo $k(a,b)=g(-a,-b)$ y cambiando el índice de suma a $(a,b)=(s-u,t-v)$, que recorre $\mathbb Z_H\times\mathbb Z_W$ cuando $(s,t)$ lo hace, se obtiene $(LI)(u,v)=\sum_{(a,b)}k(a,b)\,I(u+a,v+b)=(k*I)(u,v)$.
\end{proof}

\noindent Una aplicación lineal general de $\R^{H\times W}$ en sí mismo tiene $(HW)^2$ parámetros; el teorema dice que exigir que trate igual todas las posiciones los reduce a los $HW$ valores de $k$, y exigir además que cada salida dependa solo de un entorno (radio $r$) los reduce a $(2r+1)^2$. Traslación y localidad son las dos hipótesis que definen una capa convolucional.

\medskip

\noindent En la práctica la rejilla no es circular, y hay que decidir qué valores usar fuera de ella.

\begin{definicion}[Relleno y paso]\label{def:relleno}
\noindent El \emph{relleno circular} usa los índices módulo $H$ y $W$ (Definición~\ref{def:conv}); el \emph{relleno con ceros} define $I(u,v)=0$ fuera de $\Omega$. Una convolución con \emph{paso} $\varsigma\in\{1,2,\dots\}$ conserva solo las salidas en posiciones $(\varsigma u,\varsigma v)$.
\end{definicion}

\begin{ejemplo}[Relleno con ceros y paso rompen la equivariancia]\label{ej:rompe}
\noindent Se usan imágenes de una fila y la traslación circular $(\tau x)(v)=x(v-1\bmod W)$. (a) Con $W=3$, relleno con ceros y núcleo $k(-1)=1$, $k(0)=k(1)=0$, se tiene $(k*x)(v)=x(v-1)$ con $x(-1)=0$. Para $x=(0,0,1)$: $k*x=(0,0,0)$, luego $\tau(k*x)=(0,0,0)$; pero $\tau x=(1,0,0)$ y $k*\tau x=(0,1,0)$. (b) Con $W=4$ y paso $2$ sin filtrar, el submuestreo es $S(x)=(x(0),x(2))$. Para $x=(1,0,0,0)$, $S(x)=(1,0)$ y $S(\tau x)=S(0,1,0,0)=(0,0)$, que no es traslación de $(1,0)$ por ningún desplazamiento, porque tiene distinta suma.
\end{ejemplo}

\noindent Por eso la equivariancia del Teorema~\ref{teo:equiv} es exacta solo con relleno circular y paso $1$; en las redes reales es aproximada, con errores concentrados en los bordes y en las capas con paso.

\subsection{Redes convolucionales y vectores de características}\label{sec:cnn}

\noindent \emph{Motivación.} Un detector de bordes es una convolución fija. LeCun y colaboradores (1989) propusieron aprender los núcleos a partir de datos etiquetados, apilando varias capas para que las primeras detecten patrones simples (bordes) y las siguientes combinaciones de ellos. Para definir estas redes se necesitan cuatro piezas: capa convolucional, activación, reducción espacial y capa de clasificación.

\begin{definicion}[Capa convolucional]\label{def:capa-conv}
\noindent Una \emph{capa convolucional} de radio $r$ con $C_{\mathrm{in}}$ canales de entrada y $C_{\mathrm{out}}$ de salida tiene parámetros $K(a,b,c,c')$, con $a,b\in\{-r,\dots,r\}$, $c\le C_{\mathrm{in}}$, $c'\le C_{\mathrm{out}}$, y sesgos $\beta_{c'}$. Transforma $I\in\R^{H\times W\times C_{\mathrm{in}}}$ en
\[
L(I)(u,v,c')=\beta_{c'}+\sum_{c=1}^{C_{\mathrm{in}}}\sum_{a,b=-r}^{r}K(a,b,c,c')\,I(u+a,v+b,c).
\]
\end{definicion}

\noindent Cada canal de salida suma las convoluciones de todos los canales de entrada con núcleos propios, más una constante. Con $C_{\mathrm{in}}=C_{\mathrm{out}}=1$ y $\beta=0$ es la Definición~\ref{def:conv}.

\begin{proposicion}[Número de parámetros]\label{prop:param-conv}
\noindent Una capa convolucional tiene $(2r+1)^2C_{\mathrm{in}}C_{\mathrm{out}}+C_{\mathrm{out}}$ parámetros, cualquiera que sea $H\times W$. Una capa \emph{densa} $x\mapsto Ax+\beta$ de $\R^p$ en $\R^m$ tiene $pm+m$.
\end{proposicion}
\begin{proof}
\noindent $K$ tiene un valor por cada combinación de $a$ y $b$ ($2r+1$ valores cada uno), $c$ ($C_{\mathrm{in}}$ valores) y $c'$ ($C_{\mathrm{out}}$ valores), y hay un sesgo por canal de salida. La matriz $A$ tiene $pm$ entradas y $\beta$ tiene $m$.
\end{proof}

\noindent Con $r=1$, $C_{\mathrm{in}}=3$ y $C_{\mathrm{out}}=64$: $9\cdot3\cdot64+64=1792$ parámetros. Una capa densa sobre los $150\,528$ valores de una imagen $224\times224\times3$ con $m=1000$ salidas tiene $150\,528\,000+1000\approx1.5\times10^8$.

\begin{definicion}[Activación, reducción y red convolucional]\label{def:cnn}
\noindent (i) Una \emph{activación} es una función $\sigma:\R\to\R$ aplicada a cada valor; la más usada es $\mathrm{ReLU}(x)=\max(0,x)$. (ii) La \emph{reducción por promedio} de $2\times2$ con paso $2$ lleva $I\in\R^{H\times W\times C}$ ($H,W$ pares) a $\R^{\frac H2\times\frac W2\times C}$ con $P(I)(u,v,c)=\frac14\sum_{a,b\in\{0,1\}}I(2u+a,2v+b,c)$; la \emph{reducción por máximo} reemplaza el promedio por el máximo. (iii) El \emph{promedio global} es $\mathrm{GAP}:\R^{H\times W\times C}\to\R^C$, $\mathrm{GAP}(I)_c=\frac1{HW}\sum_{(u,v)\in\Omega}I(u,v,c)$. (iv) Una \emph{red convolucional} es una composición $\Phi=\mathrm{GAP}\circ\Lambda_J\circ\dots\circ\Lambda_1$, donde cada $\Lambda_j$ es una capa convolucional seguida de una activación (y, opcionalmente, de una reducción). (v) Un \emph{clasificador} en $K$ clases sobre $\Phi$ es $\hat y(I)=\argmax_k\big(w_k^\top\Phi(I)+\beta_k\big)$, con $w_k\in\R^{d}$ y $d$ la dimensión de $\Phi(I)$; las probabilidades asociadas son $\mathrm{softmax}(z)_k=e^{z_k}/\sum_je^{z_j}$ aplicada al vector de puntuaciones $z_k=w_k^\top\Phi(I)+\beta_k$.
\end{definicion}

\noindent El vector $\Phi(I)\in\R^d$ es el \emph{vector de características} o \emph{embedding} de la imagen que describe la guía: la salida de la red sin la última capa. La softmax es estrictamente creciente en cada puntuación y sus valores son positivos y suman $1$, así que puede leerse como distribución de probabilidad sobre las clases, y su clase de máxima probabilidad es la de máxima puntuación. Los resultados siguientes establecen qué propiedades tiene $\Phi$ por construcción, antes de cualquier entrenamiento.

\begin{teorema}[Invariancia del vector de características]\label{teo:invariancia-cnn}
\noindent Sea $\Phi=\mathrm{GAP}\circ\Lambda_J\circ\dots\circ\Lambda_1$ con capas convolucionales de relleno circular y paso $1$, activaciones aplicadas valor a valor y sin reducciones intermedias. Si $\tau_{s,t}$ actúa sobre imágenes de varios canales trasladando cada canal, entonces $\Phi(\tau_{s,t}I)=\Phi(I)$ para toda imagen $I$ y toda traslación.
\end{teorema}
\begin{proof}
\noindent Primero, cada capa convolucional $L$ conmuta con $\tau_{s,t}$: cada canal de salida es una suma de convoluciones de canales de entrada más una constante; las convoluciones conmutan con $\tau_{s,t}$ por el Teorema~\ref{teo:equiv}, la suma de funciones que conmutan con $\tau_{s,t}$ también conmuta (por linealidad de $\tau_{s,t}$), y una imagen constante no cambia al trasladarla. Segundo, una activación aplicada valor a valor conmuta con $\tau_{s,t}$, porque $\sigma\big(I(u-s,v-t,c)\big)$ es tanto $(\sigma\circ\tau_{s,t}I)(u,v,c)$ como $(\tau_{s,t}\,\sigma(I))(u,v,c)$. Tercero, por inducción en $j$, la composición $\Lambda_j\circ\dots\circ\Lambda_1$ conmuta con $\tau_{s,t}$: si $M$ y $N$ conmutan con $\tau_{s,t}$, entonces $MN\tau_{s,t}=M\tau_{s,t}N=\tau_{s,t}MN$. Por último, $\mathrm{GAP}(\tau_{s,t}J)_c=\frac1{HW}\sum_{(u,v)}J(u-s,v-t,c)=\mathrm{GAP}(J)_c$, porque $(u,v)\mapsto(u-s,v-t)$ es una biyección de $\mathbb Z_H\times\mathbb Z_W$ y reordenar una suma no la cambia. Aplicado a $J=\Lambda_J\circ\dots\circ\Lambda_1(I)$ se obtiene el resultado.
\end{proof}

\noindent El contraste con la Proposición~\ref{prop:dist-pixeles} es total: allí la representación por píxeles no distinguía un desplazamiento pequeño de uno grande; aquí el vector de características es idéntico para todos los desplazamientos. El Ejemplo~\ref{ej:rompe} muestra que las hipótesis de relleno circular y paso $1$ no pueden omitirse; en las redes reales, que usan relleno con ceros y reducciones, la invariancia es solo aproximada.

\begin{proposicion}[Regiones de decisión en el espacio de características]\label{prop:regiones}
\noindent Para el clasificador de la Definición~\ref{def:cnn}(v), el conjunto $R_k=\{z\in\R^d:w_k^\top z+\beta_k\ge w_j^\top z+\beta_j\ \forall j\}$ de vectores asignados a la clase $k$ (con empates resueltos de cualquier modo fijo) es una intersección de $K-1$ semiespacios, y por tanto convexo. En consecuencia, si el clasificador acierta en todos los registros de $\Dtr$, los vectores $\Phi(I_i)$ de clases distintas quedan separados por hiperplanos: para cada par de clases $k\ne j$, el hiperplano $(w_k-w_j)^\top z+\beta_k-\beta_j=0$ deja todos los vectores de la clase $k$ en el lado no negativo y todos los de la clase $j$ en el no positivo.
\end{proposicion}
\begin{proof}
\noindent La condición $w_k^\top z+\beta_k\ge w_j^\top z+\beta_j$ equivale a $(w_k-w_j)^\top z+(\beta_k-\beta_j)\ge0$, un semiespacio cerrado; $R_k$ es la intersección de estas condiciones para $j\ne k$, y la intersección de convexos es convexa. Si el registro $i$ de clase $k$ está bien clasificado, $\Phi(I_i)\in R_k$ y cumple cada una de las $K-1$ desigualdades; si el registro $l$ de clase $j$ está bien clasificado, cumple $(w_j-w_k)^\top\Phi(I_l)+\beta_j-\beta_k\ge0$, es decir, está en el lado no positivo del mismo hiperplano.
\end{proof}

\noindent La proposición precisa en qué sentido una red entrenada ``organiza'' sus vectores de características: la última capa es lineal, así que un buen ajuste en entrenamiento implica separación lineal de las clases \emph{de entrenamiento} en el espacio de $\Phi$. No implica nada sobre otras tareas.

\begin{definicion}[Aprendizaje por transferencia]\label{def:transfer}
\noindent Sea $\Phi$ la parte convolucional de una red entrenada en una tarea de origen (por ejemplo, clasificar un millón de fotografías en mil categorías). Para una tarea nueva con datos $(I_i,y_i)$, el \emph{aprendizaje por transferencia con características fijas} ajusta solo un modelo $g$ (lineal, un árbol, etc.) sobre los pares $(\Phi(I_i),y_i)$, sin modificar $\Phi$. El \emph{ajuste fino} (\emph{fine-tuning}) modifica también los parámetros de $\Phi$, partiendo de los de la tarea de origen.
\end{definicion}

\noindent La justificación de la transferencia es una \emph{hipótesis}, no un teorema: que los patrones útiles para la tarea de origen (bordes, texturas, partes de objetos) también lo sean para la nueva. Su ventaja cuantificable es el número de parámetros: con características fijas de dimensión $d=2048$ y un clasificador lineal de $K$ clases se estiman $2048K+K$ parámetros (Proposición~\ref{prop:param-conv} con $p=d$, $m=K$), en lugar de los millones de la red completa, lo que permite ajustar con pocos datos. Si la hipótesis falla (imágenes médicas o multiespectrales muy distintas de las fotografías de origen), el error de validación lo revela, y el ajuste fino o el entrenamiento completo son las alternativas. Es lo que la guía describe con los modelos preentrenados de SageMaker JumpStart.

\medskip

\noindent \emph{Contraste.} Los píxeles crudos conservan toda la información pero tienen dimensión enorme y una geometría inadecuada (Proposición~\ref{prop:dist-pixeles}); las características manuales son interpretables y baratas pero exigen diseño por problema; las características aprendidas son invariantes por construcción (Teorema~\ref{teo:invariancia-cnn}) y reutilizables, pero no son interpretables y dependen de la hipótesis de transferencia.

\subsection{Reducción de resolución y de dimensión}\label{sec:resolucion}

\begin{proposicion}[Reducción por promedio]\label{prop:pool}
\noindent La reducción por promedio $P$ de la Definición~\ref{def:cnn}(ii) es lineal, divide la dimensión entre $4$ y es una convolución con núcleo constante $\frac14$ en $\{0,1\}^2$ seguida de paso $2$. Si los valores de un canal son $\mu+\varepsilon(u,v)$ con errores incorrelacionados de varianza $\sigma^2$, cada valor reducido tiene media $\mu$ y varianza $\sigma^2/4$.
\end{proposicion}
\begin{proof}
\noindent Cada salida es un promedio de cuatro entradas con pesos fijos, luego $P$ es lineal; la dimensión pasa de $HWC$ a $\frac H2\frac W2C$. El promedio de $I(2u+a,2v+b,c)$ sobre $a,b\in\{0,1\}$ es la convolución con ese núcleo evaluada en $(2u,2v)$, que es la definición de paso $2$. La media y la varianza se calculan como en la Proposición~\ref{prop:media-movil}(i) con $w=4$.
\end{proof}

\noindent Reducir la resolución es la forma de llevar imágenes de distintos tamaños a uno común (Definición~\ref{def:imagen}) y la ``reducción de la resolución'' que la guía menciona como filtro de dimensionalidad; el costo es perder los detalles menores de dos píxeles. PCA sobre los vectores de píxeles o sobre los vectores de características sigue el Teorema~\ref{teo:pca} y debe ajustarse con $\Dtr$.

\subsection{Características producidas por servicios de visión}\label{sec:bbox}

\begin{definicion}[Cuadro delimitador y puntuación de confianza]\label{def:bbox}
\noindent Un \emph{cuadro delimitador} en coordenadas normalizadas es una cuádrupla $(x_0,y_0,w,h)$ con $x_0,y_0\ge0$, $w,h>0$, $x_0+w\le1$ e $y_0+h\le1$: el rectángulo $[x_0,x_0+w]\times[y_0,y_0+h]$, con las coordenadas horizontal y vertical divididas entre el ancho y el alto de la imagen. Una \emph{detección} es un triple (etiqueta, puntuación de confianza en $[0,1]$, cuadro).
\end{definicion}

\noindent Las coordenadas normalizadas no dependen de la resolución: si la imagen se reescala por un factor $\lambda>0$ en ambos ejes, las coordenadas en píxeles se multiplican por $\lambda$ y las normalizadas no cambian, porque numerador y denominador se multiplican por $\lambda$. A partir de una lista de detecciones se derivan características tabulares: el número de detecciones de cada etiqueta con confianza mayor que un umbral, el área relativa $wh\in(0,1]$ de un objeto, su centro $(x_0+\frac w2,\,y_0+\frac h2)$. Es el uso que la guía describe para Amazon Rekognition, cuyo ejemplo (``Car, confianza 98\%, en el rectángulo que empieza en $(0.12,0.40)$ con ancho $0.35$ y alto $0.20$'') da área relativa $0.07$ y centro $(0.295,0.50)$.

\section{Texto}\label{sec:texto}

\noindent \emph{Motivación.} Un documento es una cadena de caracteres de longitud variable; un modelo necesita un vector de dimensión fija. El camino que recorre la guía (tokenizar, eliminar palabras vacías, reducir a raíces, formar n-gramas y usar embeddings) reproduce la historia del procesamiento de lenguaje natural: primero representaciones de conteo, dispersas y sin noción de significado; después representaciones densas aprendidas que sí la tienen. Cada paso responde a una limitación del anterior que aquí se demuestra. Primero hay que fijar qué es un texto.

\subsection{Cadenas, documentos y corpus}\label{sec:cadenas}

\begin{definicion}[Alfabeto, cadena, documento, corpus]\label{def:cadena}
\noindent Un \emph{alfabeto} es un conjunto finito no vacío $\Sigma$ (los caracteres Unicode, o los $256$ valores de un byte). Una \emph{cadena} sobre $\Sigma$ es una sucesión finita $s=(\sigma_1,\dots,\sigma_\ell)$ de elementos de $\Sigma$; su \emph{longitud} es $\abs s=\ell$, y la cadena vacía $\varepsilon$ tiene longitud $0$. El conjunto de todas las cadenas es $\Sigma^*=\bigcup_{\ell\ge0}\Sigma^\ell$. La \emph{concatenación} de $s=(\sigma_1,\dots,\sigma_\ell)$ y $s'=(\sigma'_1,\dots,\sigma'_{m})$ es $ss'=(\sigma_1,\dots,\sigma_\ell,\sigma'_1,\dots,\sigma'_m)$. Un \emph{documento} es una cadena, y un \emph{corpus} es una sucesión finita de documentos.
\end{definicion}

\noindent La misma construcción se aplica a cualquier conjunto $A$: $A^*$ denota las sucesiones finitas de elementos de $A$. El conjunto $\Sigma^*$ es infinito numerable y no tiene estructura de espacio vectorial (no hay una suma de cadenas con sentido), así que ningún algoritmo de los capítulos anteriores puede recibir un documento directamente. Toda la sección construye aplicaciones de $\Sigma^*$ en algún $\R^p$.

\subsection{Tokens y tokenización}\label{sec:tokens}

\noindent \emph{Motivación.} La unidad mínima de $\Sigma^*$, el carácter, lleva poca información por sí sola (la letra ``e'' no dice nada del tema de un documento), y la unidad máxima, el documento completo, casi nunca se repite. Hace falta una unidad intermedia que se repita entre documentos y tenga significado: la palabra, o un fragmento de palabra. Esa unidad es el token.

\begin{definicion}[Vocabulario, token, tokenización]\label{def:token}
\noindent Un \emph{vocabulario} es un conjunto finito $V\subset\Sigma^*$ de cadenas no vacías; sus elementos se llaman \emph{tokens} (o tipos de token). Una \emph{tokenización} es una función $\tau:\Sigma^*\to V^*$ que asigna a cada documento una sucesión de tokens $\tau(d)=(t_1,\dots,t_L)$. Cada posición $\ell$ con su token $t_\ell$ es una \emph{ocurrencia}. La tokenización es \emph{sin pérdida} si existe $\rho:V^*\to\Sigma^*$ con $\rho(\tau(d))=d$ para todo $d$.
\end{definicion}

\noindent La distinción entre token y ocurrencia importa para contar: ``the cat sat on the mat'' tiene $6$ ocurrencias y $5$ tokens distintos, porque ``the'' ocurre dos veces. Una tokenización sin pérdida es inyectiva (si $\tau(d)=\tau(d')$, aplicando $\rho$ se obtiene $d=d'$), así que no confunde documentos distintos.

\begin{ejemplo}[Tokenización por espacios y por caracteres]\label{ej:tokenizadores}
\noindent (a) La \emph{tokenización por espacios} asigna a $d$ la sucesión, en orden, de sus subcadenas contiguas maximales que no contienen el carácter espacio. Para ``Machine learning is fascinating'' produce las $4$ ocurrencias \texttt{["Machine", "learning", "is", "fascinating"]} de la guía. No es sin pérdida: ``a b'' y ``a\ \ b'' (con dos espacios) tienen la misma tokenización, luego $\tau$ no es inyectiva. Además, ``fun'' y ``fun.'' son tokens distintos, lo que obliga a normalizar la puntuación antes. (b) La \emph{tokenización por caracteres}, con $V=\Sigma$ y $\tau(\sigma_1\cdots\sigma_\ell)=(\sigma_1,\dots,\sigma_\ell)$, es sin pérdida con $\rho$ igual a la concatenación, pero sus tokens tienen poco significado y las sucesiones son largas.
\end{ejemplo}

\noindent Con tokens de palabra completa y un vocabulario finito aparece un problema estructural. Como $V$ se construye con un corpus de entrenamiento, habrá palabras que no están en él.

\begin{proposicion}[Palabras fuera del vocabulario]\label{prop:oov}
\noindent Sea $\tau$ una tokenización por palabras con vocabulario finito $V$ que reemplaza cada palabra fuera de $V$ por un token especial \texttt{\textless unk\textgreater}. Entonces $\tau$ no es inyectiva.
\end{proposicion}
\begin{proof}
\noindent Supóngase que $\Sigma$ contiene algún carácter $\sigma$ distinto del espacio. Las cadenas $\sigma,\sigma\sigma,\sigma\sigma\sigma,\dots$ son infinitas palabras distintas sin espacios, y $V$ es finito, así que existen dos de ellas, $p\ne p'$, fuera de $V$. Los documentos de una palabra $p$ y $p'$ se tokenizan ambos como $(\texttt{\textless unk\textgreater})$.
\end{proof}

\noindent La solución moderna son los tokens de subpalabra. El algoritmo BPE (\emph{byte pair encoding}) nació como método de compresión (Gage, 1994) y Sennrich, Haddow y Birch (2016) lo adaptaron para construir vocabularios.

\begin{definicion}[BPE]\label{def:bpe}
\noindent Sea un corpus de entrenamiento representado como una lista de palabras, cada una escrita como sucesión de símbolos, inicialmente sus caracteres, y sea $V_0=\Sigma$. Para $j=1,\dots,J$: se cuenta, para cada par ordenado de símbolos $(x,y)$, cuántas veces aparecen adyacentes dentro de una palabra; se elige el par $(x_j,y_j)$ más frecuente (con un criterio fijo de desempate); se añade el símbolo concatenado $x_jy_j$ al vocabulario, $V_j=V_{j-1}\cup\{x_jy_j\}$; y se reemplaza cada aparición de $x_j$ seguido de $y_j$ por $x_jy_j$, de izquierda a derecha y sin solapamientos. Para tokenizar un texto nuevo se parte de sus caracteres y se aplican las fusiones $j=1,\dots,J$ en el mismo orden.
\end{definicion}

\begin{proposicion}[Propiedades de BPE]\label{prop:bpe}
\noindent (i) Todo documento sobre $\Sigma$ se tokeniza con $V_J$: no hay palabras fuera del vocabulario. (ii) La tokenización es sin pérdida dentro de cada palabra: la concatenación de los tokens es la palabra original. (iii) Si en la fusión $j$ se reemplazan $f_j\ge1$ apariciones, el número total de ocurrencias del corpus baja exactamente en $f_j$, mientras que $\abs{V_j}\le\abs{V_{j-1}}+1$.
\end{proposicion}
\begin{proof}
\noindent (i) El proceso de tokenización parte de los caracteres, que están en $V_0=\Sigma\subset V_J$, y cada fusión solo produce símbolos de $V_J$; nunca hace falta un símbolo ajeno. (ii) Por inducción en el número de fusiones aplicadas. Antes de la primera, la concatenación de los símbolos es la palabra. Si la concatenación es la palabra antes de una fusión, reemplazar dos símbolos adyacentes $x,y$ por el símbolo $xy$ no cambia la concatenación total, porque la concatenación es asociativa: $(\cdots)\,x\,y\,(\cdots)=(\cdots)\,(xy)\,(\cdots)$. (iii) Cada reemplazo convierte dos ocurrencias en una, así que $f_j$ reemplazos reducen el total en $f_j$. Se añade a lo sumo un símbolo nuevo (ninguno si $x_jy_j$ ya estaba en el vocabulario).
\end{proof}

\begin{ejemplo}[Dos fusiones]\label{ej:bpe}
\noindent Sea el corpus de cuatro palabras ``ab'', ``ab'', ``ab'', ``ac'', con $8$ ocurrencias iniciales. Los pares adyacentes son $(a,b)$ con frecuencia $3$ y $(a,c)$ con frecuencia $1$; la primera fusión crea ``ab'' y deja $[ab],[ab],[ab],[a,c]$, con $8-3=5$ ocurrencias. La segunda fusión crea ``ac'' y deja $4$. La palabra frecuente se convirtió en un token en la primera fusión; la rara, en la segunda.
\end{ejemplo}

\noindent La parte (iii) muestra el compromiso que controla $J$: más fusiones dan vocabularios más grandes y sucesiones más cortas. Por eso el número de tokens de un texto no coincide con su número de palabras: una palabra frecuente suele ser un solo token y una rara se parte en varios. La regla que cita la guía (un token equivale a unos $\frac34$ de palabra en inglés) es una observación empírica sobre tokenizadores concretos, no una propiedad matemática, y cambia con el idioma y el vocabulario. En Amazon Bedrock, el costo de un modelo de texto en modalidad bajo demanda es lineal en las ocurrencias: con $T_{\mathrm{in}}$ tokens de entrada, $T_{\mathrm{out}}$ de salida y precios por token $p_{\mathrm{in}},p_{\mathrm{out}}$, el costo es $p_{\mathrm{in}}T_{\mathrm{in}}+p_{\mathrm{out}}T_{\mathrm{out}}$. Con la regla empírica, una página de $500$ palabras son unos $500/0.75\approx667$ tokens. Los modelos que generan imágenes suelen cobrarse por imagen.

\subsection{Bolsa de palabras}\label{sec:bow}

\begin{definicion}[Bolsa de palabras]\label{def:bow}
\noindent Dada una tokenización con vocabulario $V$, la \emph{bolsa de palabras} de un documento $d$ con $\tau(d)=(t_1,\dots,t_L)$ es el vector de conteos $x(d)\in\N^{V}$ con $x_v(d)=\#\{\ell:t_\ell=v\}$. Equivalentemente, $x(d)=\sum_{\ell=1}^Le_{t_\ell}$, donde $e_v$ es el vector canónico de $v$ (la codificación one-hot de la Definición~\ref{def:onehot}).
\end{definicion}

\noindent La equivalencia se verifica coordenada a coordenada: la coordenada $v$ de $\sum_\ell e_{t_\ell}$ suma $1$ por cada $\ell$ con $t_\ell=v$. El vocabulario es un parámetro ajustado: se construye con $\Dtr$. La bolsa de palabras es la primera aplicación de esta sección de $\Sigma^*$ en un espacio vectorial, y hereda de one-hot la dimensión $\abs V$ (decenas o cientos de miles) y la dispersión. Tiene además una pérdida de información específica.

\begin{proposicion}[La bolsa de palabras ignora el orden]\label{prop:bow-orden}
\noindent Si $\tau(d')$ es una permutación de $\tau(d)$, entonces $x(d)=x(d')$. En particular, $x$ no es inyectiva sobre documentos de dos o más tokens distintos.
\end{proposicion}
\begin{proof}
\noindent Si $\tau(d')=(t_{\pi(1)},\dots,t_{\pi(L)})$ con $\pi$ una permutación, $x(d')=\sum_\ell e_{t_{\pi(\ell)}}=\sum_\ell e_{t_\ell}=x(d)$, porque una suma finita no cambia al reordenar sus sumandos.
\end{proof}

\noindent ``dog bites man'' y ``man bites dog'' tienen la misma bolsa de palabras y significados distintos. Los n-gramas (Sección~\ref{sec:ngramas}) recuperan parte del orden.

\subsection{Palabras vacías y ponderación TF-IDF}\label{sec:stop}

\noindent \emph{Motivación.} Luhn (1958), al construir resúmenes automáticos, observó que las palabras más frecuentes de un texto (artículos, preposiciones, auxiliares) no indican su tema y dominan los conteos. La solución más directa es descartarlas.

\begin{definicion}[Eliminación de palabras vacías]\label{def:stop}
\noindent Dada una lista $S\subset V$ de \emph{palabras vacías}, su eliminación es la aplicación $P_S:\R^V\to\R^V$, $(P_Sx)_v=x_v\,\ind{v\notin S}$.
\end{definicion}

\begin{proposicion}[La eliminación es una proyección ortogonal]\label{prop:stop}
\noindent $P_S$ es lineal, idempotente ($P_S^2=P_S$) y simétrica ($\ip{P_Sx}{y}=\ip{x}{P_Sy}$); es la proyección ortogonal sobre el subespacio de vectores nulos en las coordenadas de $S$.
\end{proposicion}
\begin{proof}
\noindent $P_S$ multiplica cada coordenada por una constante $\ind{v\notin S}\in\{0,1\}$, luego es lineal. Como $\ind{v\notin S}^2=\ind{v\notin S}$, aplicarla dos veces es aplicarla una. $\ip{P_Sx}{y}=\sum_vx_v\ind{v\notin S}y_v=\ip x{P_Sy}$. Una aplicación lineal idempotente y simétrica es la proyección ortogonal sobre su imagen, que es el subespacio indicado.
\end{proof}

\noindent La proyección reduce la dimensión efectiva en $\abs S$, pero es una decisión fija que no mira la tarea, y puede eliminar información esencial.

\begin{ejemplo}[Negaciones]\label{ej:negacion}
\noindent Si ``not'' $\in S$, los documentos ``not good'' y ``good'' cumplen $P_Sx(\text{``not good''})=P_Sx(\text{``good''})=e_{\text{good}}$. En análisis de sentimiento, dos textos de polaridad opuesta reciben la misma representación. Es la razón del ``a menudo'' de la guía.
\end{ejemplo}

\noindent Spärck Jones (1972) propuso en su lugar ponderar cada palabra según su rareza en la colección, decidiendo con los datos cuánto pesa.

\begin{definicion}[TF-IDF]\label{def:tfidf}
\noindent Sea una colección de $N$ documentos y $\mathrm{df}(v)$ el número de documentos que contienen el token $v$ al menos una vez (se consideran solo los $v$ con $\mathrm{df}(v)\ge1$). La \emph{frecuencia inversa en documentos} es $\idf(v)=\log\frac N{\mathrm{df}(v)}$ y la representación TF-IDF de $d$ es el vector con coordenadas $x_v(d)\,\idf(v)$.
\end{definicion}

\begin{proposicion}[Propiedades de la IDF]\label{prop:idf}
\noindent (i) $0\le\idf(v)\le\log N$. (ii) $\idf$ es estrictamente decreciente en $\mathrm{df}(v)$. (iii) $\idf(v)=0$ si y solo si $v$ aparece en todos los documentos. (iv) La variante de scikit-learn, $\log\frac{1+N}{1+\mathrm{df}(v)}+1$, es siempre $\ge1$.
\end{proposicion}
\begin{proof}
\noindent (i) De $1\le\mathrm{df}(v)\le N$ se sigue $1\le N/\mathrm{df}(v)\le N$, y el logaritmo es creciente. (ii) $N/\mathrm{df}$ es estrictamente decreciente en $\mathrm{df}$ y el logaritmo es estrictamente creciente. (iii) $\log(N/\mathrm{df})=0$ si y solo si $\mathrm{df}=N$. (iv) $\mathrm{df}\le N$ implica $\frac{1+N}{1+\mathrm{df}}\ge1$, cuyo logaritmo es $\ge0$.
\end{proof}

\noindent \emph{Contraste.} Por (iii), las palabras vacías más extremas (presentes en todos los documentos) se anulan sin lista, y por (ii) las intermedias se atenúan en lugar de eliminarse; la decisión depende de la colección y no de una lista fija. Pero TF-IDF sigue siendo una bolsa de palabras reponderada: hereda la Proposición~\ref{prop:bow-orden}, y no resuelve el Ejemplo~\ref{ej:negacion} si ``not'' es frecuente. Los $\mathrm{df}(v)$ son parámetros ajustados con $\Dtr$. La transformación \emph{Vectorize} de Data Wrangler produce vectores TF-IDF.

\subsection{Stemming y lematización}\label{sec:stemming}

\noindent \emph{Motivación.} ``run'', ``runs'' y ``running'' son tokens distintos en la bolsa de palabras, así que un modelo aprende por separado lo que dice cada uno, con menos datos para cada uno. Reducirlos a una forma común agrupa sus conteos.

\begin{definicion}[Normalización morfológica]\label{def:normalizacion}
\noindent Una \emph{normalización morfológica} es una función $\pi:V\to V'$ que asigna a cada token una forma base. Induce la aplicación lineal $A_\pi:\R^V\to\R^{V'}$, $(A_\pi x)_{v'}=\sum_{v:\,\pi(v)=v'}x_v$. El \emph{stemming} define $\pi$ mediante reglas sobre la cadena (recorte de sufijos), sin diccionario; el más usado es el algoritmo de Porter (1980). La \emph{lematización} asigna a cada token su \emph{lema}, la forma de diccionario de la palabra, usando un diccionario y la categoría gramatical del token en su contexto (formalmente, $\pi$ se define sobre pares (token, categoría)).
\end{definicion}

\begin{proposicion}[Qué conserva y qué pierde una normalización]\label{prop:normalizacion}
\noindent (i) $A_\pi$ conserva el número total de ocurrencias: $\sum_{v'}(A_\pi x)_{v'}=\sum_vx_v$. (ii) $\abs{\pi(V)}\le\abs V$. (iii) $A_\pi x=A_\pi y$ si y solo si, para cada $v'$, los conteos de $x$ e $y$ sumados sobre la fibra $\pi^{-1}(v')$ coinciden; en particular, $A_\pi$ es inyectiva si y solo si $\pi$ lo es.
\end{proposicion}
\begin{proof}
\noindent (i) Las fibras $\pi^{-1}(v')$ forman una partición de $V$, así que sumar por fibras y luego sobre $v'$ es sumar sobre $V$. (ii) La imagen de un conjunto finito por una función tiene a lo sumo tantos elementos como el conjunto. (iii) La primera parte es la definición de $A_\pi$ coordenada a coordenada. Si $\pi$ es inyectiva, cada fibra tiene a lo sumo un elemento y $A_\pi$ solo reordena coordenadas. Si $\pi(v)=\pi(w)$ con $v\ne w$, entonces $A_\pi e_v=A_\pi e_w$ con $e_v\ne e_w$.
\end{proof}

\noindent Matemáticamente, la normalización es la misma operación que el hashing de la Definición~\ref{def:hash} (sin signo): una aplicación lineal que suma coordenadas que ``colisionan''. La diferencia es que aquí las colisiones se eligen con criterio lingüístico para que sean deseables, mientras que en el hashing son arbitrarias. La calidad de $\pi$ se mide contra la función de lema verdadera $\lambda$.

\begin{definicion}[Sobrerreducción y subreducción]\label{def:sobre}
\noindent Sea $\lambda$ la función que asigna a cada token su lema correcto. Una normalización $\pi$ comete \emph{sobrerreducción} en $v,w$ si $\pi(v)=\pi(w)$ pero $\lambda(v)\ne\lambda(w)$, y \emph{subreducción} si $\lambda(v)=\lambda(w)$ pero $\pi(v)\ne\pi(w)$.
\end{definicion}

\begin{ejemplo}[Errores del algoritmo de Porter]\label{ej:porter}
\noindent (Verificado con la implementación de Porter de NLTK.) Sobrerreducción: ``universe'', ``university'' y ``universal'' van todos a ``univers'', y ``organization'' va a ``organ'', la misma forma que ``organ''. Subreducción: ``run'', ``runs'' y ``running'' van a ``run'', pero ``ran'' queda como ``ran'', porque la relación entre ``ran'' y ``run'' no es un sufijo. Además, el stemming puede producir cadenas que no son palabras: ``studies'' $\mapsto$ ``studi''. La lematización, con diccionario y categoría gramatical, lleva ``ran'' a ``run'', ``studies'' a ``study'' y el adjetivo ``better'' a ``good'', y evita las sobrerreducciones anteriores, a costa de un diccionario y de analizar la categoría de cada token.
\end{ejemplo}

\noindent \emph{Contraste.} Ambas agrupan conteos y reducen la dimensión (Proposición~\ref{prop:normalizacion}), lo que mejora la estimación de las formas raras; ambas pierden la información morfológica que distinguía los tokens de una misma fibra (tiempo verbal, número), que en algunas tareas importa. El stemming es rápido y comete los dos tipos de error; la lematización comete menos y es más lenta.

\subsection{N-gramas}\label{sec:ngramas}

\noindent \emph{Motivación.} Shannon (1948) aproximó el inglés mediante modelos en que cada letra o palabra depende de las $n-1$ anteriores, y mostró que al crecer $n$ los textos generados se parecen cada vez más al idioma. Las mismas unidades, secuencias de $n$ tokens consecutivos, sirven como características que conservan el orden local que la bolsa de palabras pierde.

\begin{definicion}[N-grama]\label{def:ngrama}
\noindent Sea $\tau(d)=(t_1,\dots,t_L)$ y $1\le n\le L$. Los \emph{n-gramas} de $d$ son las sucesiones $(t_\ell,\dots,t_{\ell+n-1})\in V^n$ para $\ell=1,\dots,L-n+1$. La \emph{bolsa de n-gramas} es el vector de conteos de n-gramas, indexado por $V^n$. Con $n=2$ se llaman bigramas.
\end{definicion}

\begin{proposicion}[Número y poder de separación de los n-gramas]\label{prop:ngramas}
\noindent (i) Un documento de $L\ge n$ tokens tiene exactamente $L-n+1$ ocurrencias de n-gramas, y el espacio de n-gramas posibles tiene $\abs V^n$ elementos. (ii) Existen documentos con la misma bolsa de palabras y distinta bolsa de bigramas.
\end{proposicion}
\begin{proof}
\noindent (i) El n-grama que empieza en la posición $\ell$ existe si y solo si $\ell+n-1\le L$, es decir, $\ell\in\{1,\dots,L-n+1\}$. Un elemento de $V^n$ es una elección de un token para cada una de $n$ posiciones, con $\abs V$ opciones cada una. (ii) ``dog bites man'' y ``man bites dog'' tienen la misma bolsa de palabras por la Proposición~\ref{prop:bow-orden}. Los bigramas del primero son $(\text{dog},\text{bites})$ y $(\text{bites},\text{man})$; los del segundo, $(\text{man},\text{bites})$ y $(\text{bites},\text{dog})$. El bigrama $(\text{dog},\text{bites})$ tiene conteo $1$ en el primero y $0$ en el segundo.
\end{proof}

\noindent En ``Machine learning is fun'' ($L=4$) hay $4-2+1=3$ bigramas, como en la guía. Los bigramas también resuelven el Ejemplo~\ref{ej:negacion} si no se eliminan palabras vacías: ``not good'' contiene el bigrama $(\text{not},\text{good})$. El costo está en (i): con $\abs V=5\times10^4$ hay $2.5\times10^9$ bigramas posibles; aunque solo aparezca una fracción, cada uno se observa pocas veces, y se vuelve a la maldición de la dimensionalidad (Sección~\ref{sec:seleccion}). En la práctica se usan $n\le3$, con un umbral de frecuencia mínima o con hashing.

\subsection{Embeddings de palabras}\label{sec:embeddings}

\noindent \emph{Motivación.} Todas las representaciones anteriores tratan los tokens como símbolos sin relación entre sí.

\begin{definicion}[Similitud coseno]\label{def:coseno}
\noindent Para $a,b\in\R^D$ no nulos, $\cos(a,b)=\ip ab/(\norm a\norm b)\in[-1,1]$, el coseno del ángulo entre ellos (el rango se sigue de la desigualdad de Cauchy--Schwarz).
\end{definicion}

\begin{proposicion}[Ortogonalidad de las representaciones one-hot]\label{prop:ortogonal}
\noindent Para tokens $v\ne v'$, $\cos(e_v,e_{v'})=0$.
\end{proposicion}
\begin{proof}
\noindent $\ip{e_v}{e_{v'}}=0$ para $v\ne v'$ y $\norm{e_v}=\norm{e_{v'}}=1$.
\end{proof}

\noindent ``king'' es tan distinto de ``queen'' como de ``banana'', y lo que un modelo aprende sobre una palabra no se transfiere a sus sinónimos. La idea que resolvió el problema es la \emph{hipótesis distribucional} (Harris, 1954; Firth, 1957): palabras que aparecen en contextos parecidos tienen significados parecidos. Es una hipótesis lingüística, no un teorema; lo que sí puede demostrarse es qué representación produce cada método si se la acepta. Para eso hay que precisar qué es un contexto.

\begin{definicion}[Ventana de contexto y conteos de coocurrencia]\label{def:contexto}
\noindent Sea $c\ge1$ el \emph{radio de la ventana}. En un documento tokenizado $(t_1,\dots,t_L)$, los \emph{contextos} de la ocurrencia en la posición $\ell$ son los tokens $t_{\ell+j}$ con $1\le\abs j\le c$ y $1\le\ell+j\le L$. El multiconjunto $\mathcal D$ de \emph{pares de coocurrencia} contiene un par $(t_\ell,t_{\ell+j})$ por cada ocurrencia y cada uno de sus contextos, en todo el corpus. Se escribe $\#(w,c)$ para el número de pares con palabra $w$ y contexto $c$, $\#(w)=\sum_c\#(w,c)$, $\#(c)=\sum_w\#(w,c)$ y $\abs{\mathcal D}=\sum_{w,c}\#(w,c)$. La \emph{información mutua puntual} empírica, para $\#(w,c)>0$, es
\[
\PMI(w,c)=\log\frac{\#(w,c)\,\abs{\mathcal D}}{\#(w)\,\#(c)}=\log\frac{\hat p(w,c)}{\hat p(w)\,\hat p(c)},\qquad \hat p(w,c)=\frac{\#(w,c)}{\abs{\mathcal D}},\ \hat p(w)=\frac{\#(w)}{\abs{\mathcal D}},\ \hat p(c)=\frac{\#(c)}{\abs{\mathcal D}}.
\]
\end{definicion}

\noindent La igualdad de las dos expresiones de la PMI se obtiene sustituyendo las proporciones empíricas: $\frac{\#(w,c)/\abs{\mathcal D}}{\#(w)\#(c)/\abs{\mathcal D}^2}=\frac{\#(w,c)\abs{\mathcal D}}{\#(w)\#(c)}$. La PMI es el sumando de la información mutua de la Definición~\ref{def:im} para la distribución empírica de pares: es positiva cuando $w$ y $c$ coocurren más de lo que lo harían si fueran independientes.

\begin{definicion}[Embedding de palabras]\label{def:embedding}
\noindent Un \emph{embedding} de dimensión $D$ es una función $v:V\to\R^D$, $w\mapsto v_w$, equivalentemente una matriz de $\abs V\times D$ cuyas filas son los $v_w$, con $D$ mucho menor que $\abs V$ (típicamente del orden de $10^2$).
\end{definicion}

\noindent \emph{Word2Vec} (Mikolov et al., 2013) aprende dos embeddings, $v_w$ para palabras y $u_c$ para contextos. En su variante \emph{skip-gram} modela la probabilidad de un contexto dada una palabra como $p(c\mid w)=\mathrm{softmax}\big((u_{c'}^\top v_w)_{c'\in V}\big)_c=\exp(u_c^\top v_w)/\sum_{c'\in V}\exp(u_{c'}^\top v_w)$ y maximiza $\sum_{(w,c)\in\mathcal D}\log p(c\mid w)$. El denominador suma sobre todo el vocabulario, así que evaluar cada término cuesta del orden de $\abs V$ operaciones. El muestreo negativo lo reemplaza por un problema de clasificación binaria.

\begin{definicion}[Skip-gram con muestreo negativo]\label{def:sgns}
\noindent Sea $\sigma(x)=1/(1+e^{-x})$ la función logística y $k\ge1$ un entero. Por cada par observado $(w,c)\in\mathcal D$ se muestrean $k$ contextos \emph{negativos} $c_N$ independientes con la distribución unigrama $P(c_N=c)=\#(c)/\abs{\mathcal D}$. El objetivo (SGNS) a maximizar es
\[
\sum_{(w,c)\in\mathcal D}\Big(\log\sigma(u_c^\top v_w)+\sum_{\text{negativos }c_N}\log\sigma(-u_{c_N}^\top v_w)\Big).
\]
\end{definicion}

\noindent Como $\sigma(x)\in(0,1)$ y $\sigma(-x)=1-\sigma(x)$, $\sigma(u_c^\top v_w)$ se interpreta como la probabilidad de que el par sea observado y no muestreado al azar; el objetivo es la log-verosimilitud de una regresión logística que distingue unos de otros. Levy y Goldberg (2014) mostraron qué calcula este procedimiento.

\begin{teorema}[SGNS factoriza la PMI desplazada]\label{teo:sgns}
\noindent La esperanza del objetivo de la Definición~\ref{def:sgns}, respecto del muestreo negativo, es $\sum_{w,c}\ell_{wc}(u_c^\top v_w)$ con
\[
\ell_{wc}(x)=\#(w,c)\log\sigma(x)+k\,\frac{\#(w)\#(c)}{\abs{\mathcal D}}\log\sigma(-x).
\]
Si $\#(w,c)>0$, $\ell_{wc}$ es estrictamente cóncava y se maximiza únicamente en $x^*=\PMI(w,c)-\log k$. En consecuencia, si la dimensión $D$ permite elegir los productos $u_c^\top v_w$ de forma independiente para cada par, el máximo de la esperanza del objetivo se alcanza con $u_c^\top v_w=\PMI(w,c)-\log k$ para todo par con $\#(w,c)>0$.
\end{teorema}
\begin{proof}
\noindent El primer sumando de la Definición~\ref{def:sgns} aporta $\log\sigma(u_c^\top v_w)$ una vez por cada una de las $\#(w,c)$ apariciones del par. Para el segundo, la palabra $w$ aparece como primer elemento de $\#(w)$ pares; en cada uno se muestrean $k$ negativos, y cada negativo es igual a $c$ con probabilidad $\#(c)/\abs{\mathcal D}$. Por linealidad de la esperanza, el número esperado de veces que se suma $\log\sigma(-u_c^\top v_w)$ es $k\#(w)\#(c)/\abs{\mathcal D}$. Agrupando por pares $(w,c)$ se obtiene la expresión con $\ell_{wc}$.

\medskip

\noindent De $\sigma(x)=1/(1+e^{-x})$ se obtiene $\frac{d}{dx}\log\sigma(x)=\frac{e^{-x}}{1+e^{-x}}=\sigma(-x)$ y, por la regla de la cadena, $\frac d{dx}\log\sigma(-x)=-\sigma(x)$; además $\sigma'(x)=\sigma(x)\sigma(-x)$. Con $B=k\#(w)\#(c)/\abs{\mathcal D}>0$,
\[
\ell_{wc}'(x)=\#(w,c)\,\sigma(-x)-B\,\sigma(x),\qquad \ell_{wc}''(x)=-\big(\#(w,c)+B\big)\sigma(x)\sigma(-x)<0.
\]
La concavidad estricta garantiza que un punto crítico, si existe, es el máximo único. La condición $\ell'_{wc}(x)=0$ equivale a $\sigma(x)/\sigma(-x)=\#(w,c)/B$, y $\frac{\sigma(x)}{\sigma(-x)}=\frac{1+e^{x}}{1+e^{-x}}=e^x$ (multiplicando numerador y denominador por $e^x$). Luego $x^*=\log\frac{\#(w,c)}{B}=\log\frac{\#(w,c)\abs{\mathcal D}}{\#(w)\#(c)}-\log k$. Si los productos pueden elegirse de forma independiente, la suma se maximiza maximizando cada sumando.
\end{proof}

\noindent La hipótesis sobre $D$ es fuerte (con $D\ll\abs V$ no se cumple en general): el teorema describe el objetivo al que el entrenamiento se aproxima, y los embeddings reales son una factorización aproximada de rango $D$ de la matriz de PMI desplazada. Aun así, permite dar una versión exacta de la hipótesis distribucional.

\begin{corolario}[Contextos iguales, vectores iguales]\label{cor:contextos}
\noindent Supóngase que en el óptimo del Teorema~\ref{teo:sgns} dos palabras $w,w'$ tienen las mismas frecuencias condicionales de contexto, es decir, $\#(w,c)/\#(w)=\#(w',c)/\#(w')$ para todo $c$, con $\#(w,c)>0$ si y solo si $\#(w',c)>0$, y que los vectores $u_c$ con $\#(w,c)>0$ generan $\R^D$. Entonces $v_w=v_{w'}$.
\end{corolario}
\begin{proof}
\noindent Por la Definición~\ref{def:contexto}, $\PMI(w,c)=\log\big(\frac{\#(w,c)}{\#(w)}\cdot\frac{\abs{\mathcal D}}{\#(c)}\big)$, que por hipótesis coincide con $\PMI(w',c)$ para todo $c$ con conteo positivo. Por el teorema, $u_c^\top v_w=u_c^\top v_{w'}$, es decir, $u_c^\top(v_w-v_{w'})=0$, para todos esos $c$. Un vector ortogonal a un sistema generador de $\R^D$ es nulo, porque es ortogonal a toda combinación lineal de ese sistema, en particular a sí mismo.
\end{proof}

\noindent \emph{GloVe} (Pennington, Socher y Manning, 2014) parte de la matriz de coocurrencias $X_{ij}=\#(i,j)$ y define los embeddings como minimizadores del error cuadrático ponderado
\[
J=\sum_{i,j:\,X_{ij}>0}f(X_{ij})\big(w_i^\top\tilde w_j+b_i+\tilde b_j-\log X_{ij}\big)^2,\qquad f(x)=\min\big(1,(x/x_{\max})^{\alpha}\big),
\]
con $w_i,\tilde w_j\in\R^D$, sesgos escalares $b_i,\tilde b_j$ y, en el artículo original, $\alpha=\frac34$. El peso $f$ es creciente y acotado por $1$: impide que los pares muy frecuentes dominen la suma, y los pares no observados no entran (el logaritmo de $0$ no está definido).

\begin{proposicion}[GloVe y la PMI]\label{prop:glove}
\noindent Si un minimizador de $J$ ajusta exactamente ($J=0$) y sus sesgos son $b_i=\log X_{i\cdot}$ y $\tilde b_j=\log X_{\cdot j}-\log X_{\cdot\cdot}$, con $X_{i\cdot}=\sum_jX_{ij}$, $X_{\cdot j}=\sum_iX_{ij}$ y $X_{\cdot\cdot}=\sum_{i,j}X_{ij}$, entonces $w_i^\top\tilde w_j=\PMI(i,j)$ para todo par con $X_{ij}>0$.
\end{proposicion}
\begin{proof}
\noindent $J=0$ con pesos positivos exige que cada término sea nulo: $w_i^\top\tilde w_j=\log X_{ij}-b_i-\tilde b_j=\log X_{ij}-\log X_{i\cdot}-\log X_{\cdot j}+\log X_{\cdot\cdot}=\log\frac{X_{ij}X_{\cdot\cdot}}{X_{i\cdot}X_{\cdot j}}$, que es la PMI de la Definición~\ref{def:contexto} con $\#(w)=X_{i\cdot}$, $\#(c)=X_{\cdot j}$ y $\abs{\mathcal D}=X_{\cdot\cdot}$.
\end{proof}

\noindent Así, las dos técnicas que la guía presenta como ``una red pequeña'' y ``una factorización'' aproximan el mismo objeto, la matriz de PMI de las coocurrencias; difieren en la función de pérdida (logística frente a cuadrática) y en el peso de cada par. La guía dice que ``king'' y ``queen'' ``podrían'' quedar cerca: es una observación empírica sobre embeddings entrenados en corpus grandes (medida con la similitud coseno), no una consecuencia de las definiciones; el Corolario~\ref{cor:contextos} es lo que puede garantizarse.

\medskip

\noindent \emph{Contraste de representaciones de texto.}
\begin{center}
\begin{tabularx}{\linewidth}{@{}lcY@{}}
\toprule
Representación & Dimensión & Ventajas y limitaciones\\
\midrule
Bolsa de palabras & $\abs V$, dispersa & simple, interpretable; ignora el orden (Prop.~\ref{prop:bow-orden}) y la similitud (Prop.~\ref{prop:ortogonal})\\
Sin palabras vacías & $\abs V-\abs S$ & proyección fija (Prop.~\ref{prop:stop}); puede borrar negaciones (Ej.~\ref{ej:negacion})\\
Stemming / lematización & $\abs{\pi(V)}$ & agrupa conteos (Prop.~\ref{prop:normalizacion}); sobre y subreducción (Ej.~\ref{ej:porter})\\
TF-IDF & $\abs V$, dispersa & pondera por rareza con los datos (Prop.~\ref{prop:idf}); mismas limitaciones de orden y similitud\\
N-gramas & hasta $\abs V^n$ & recuperan orden local (Prop.~\ref{prop:ngramas}); explosión combinatoria\\
Embeddings & $D\approx10^2$, densa & aproximan la PMI (Teo.~\ref{teo:sgns}, Prop.~\ref{prop:glove}); requieren un corpus grande; un vector por palabra aunque tenga varios sentidos; sin vector para palabras fuera de $V$ (Prop.~\ref{prop:oov})\\
\bottomrule
\end{tabularx}
\end{center}

\section{Desbalance de clases}\label{sec:desbalance}

\noindent \emph{Motivación.} En clasificación binaria con una clase positiva rara (enfermedad, fraude) de proporción $\pi$, los algoritmos que minimizan el error promedio encuentran rentable ignorarla. El resultado siguiente, que la guía ilustra con el $99\%$ de personas sanas, lo hace preciso.

\begin{definicion}[Clasificador, exactitud y sensibilidad]\label{def:exactitud}
\noindent Un \emph{clasificador} binario es una función $\hat y:\R^p\to\{0,1\}$. Si $(X,Y)$ es un registro aleatorio, su \emph{exactitud} (\emph{accuracy}) es $\Prob(\hat y(X)=Y)$ y su \emph{sensibilidad} (\emph{recall}) es $\Prob(\hat y(X)=1\mid Y=1)$, la fracción de positivos que detecta.
\end{definicion}

\noindent La exactitud trata igual los dos tipos de error; la sensibilidad mide solo lo que ocurre con la clase positiva.

\begin{proposicion}[Paradoja de la exactitud y umbral de Bayes]\label{prop:paradoja}
\noindent Sea $\eta(x)=\Prob(Y=1\mid X=x)$ y $\pi=\Prob(Y=1)$. (i) El clasificador constante $\hat y\equiv0$ tiene exactitud $1-\pi$ y sensibilidad $0$. (ii) Si un error sobre la clase $1$ cuesta $w_1>0$ y uno sobre la clase $0$ cuesta $w_0>0$, el clasificador que minimiza el costo esperado predice $1$ si y solo si $\eta(x)>\dfrac{w_0}{w_0+w_1}$. Con $w_0=w_1$ el umbral es $\frac12$.
\end{proposicion}
\begin{proof}
\noindent (i) Acierta exactamente cuando $Y=0$, con probabilidad $1-\pi$, y nunca predice $1$. (ii) Condicionando a $X=x$, predecir $1$ tiene costo esperado $w_0(1-\eta(x))$ (se equivoca si $Y=0$) y predecir $0$ tiene costo $w_1\eta(x)$. Conviene predecir $1$ si $w_0(1-\eta(x))<w_1\eta(x)$, es decir, $w_0<(w_0+w_1)\eta(x)$. Minimizar punto a punto minimiza la esperanza total, porque el costo esperado es la integral de los costos condicionales.
\end{proof}

\noindent Si la clase positiva es rara, $\eta(x)$ rara vez supera $\frac12$ y el clasificador óptimo para la exactitud casi nunca predice la clase de interés. Las técnicas de mitigación de la guía cambian, de maneras distintas, los pesos efectivos $w_0,w_1$ o los datos.

\subsection{Ponderación y remuestreo}\label{sec:ponderacion}

\begin{definicion}[Pesos balanceados]\label{def:pesos}
\noindent Con $K$ clases de tamaños $n_1,\dots,n_K$ y $n=\sum n_c$, los \emph{pesos balanceados} son $w_c=\dfrac{n}{K\,n_c}$, y el riesgo ponderado es $\hat R_w(f)=\frac1n\sum_iw_{y_i}\ell(f(x_i),y_i)$.
\end{definicion}

\noindent La fórmula es la que usa \texttt{class\_weight="balanced"} en scikit-learn. Su justificación es que iguala la contribución total de cada clase: la suma de los pesos de la clase $c$ es $n_c\cdot\frac{n}{Kn_c}=\frac nK$, la misma para todas. En el caso binario, el cociente $w_1/w_0=n_0/n_1$ es el que usa \texttt{scale\_pos\_weight} en XGBoost, y por la Proposición~\ref{prop:paradoja}(ii) el umbral de Bayes pasa de $\frac12$ a $\frac{w_0}{w_0+w_1}=\frac{n_1}{n_0+n_1}=\hat\pi$: se predice la clase rara en cuanto su probabilidad condicional supera su frecuencia marginal.

\begin{proposicion}[Remuestreo y ponderación coinciden en esperanza]\label{prop:remuestreo}
\noindent (i) \emph{Sobremuestreo aleatorio}: si cada ejemplo de la clase $1$ se incluye $M_i$ veces, con $M_i$ aleatorios de media $r$ independientes de los datos, el riesgo empírico (sin normalizar) tiene esperanza $\sum_{y_i=0}\ell_i+r\sum_{y_i=1}\ell_i$. (ii) \emph{Submuestreo aleatorio}: si cada ejemplo de la clase $0$ se conserva con probabilidad $\beta$, independientemente, la esperanza es $\beta\sum_{y_i=0}\ell_i+\sum_{y_i=1}\ell_i$. En ambos casos coincide con un riesgo ponderado con $w_1/w_0=r$ o $w_1/w_0=1/\beta$.
\end{proposicion}
\begin{proof}
\noindent Por la Proposición~\ref{prop:duplicado}, un ejemplo incluido $M$ veces contribuye $M\ell_i$. (i) $\E[\sum_{y_i=0}\ell_i+\sum_{y_i=1}M_i\ell_i]=\sum_{y_i=0}\ell_i+r\sum_{y_i=1}\ell_i$ por linealidad. (ii) Con $B_i\sim\mathrm{Bernoulli}(\beta)$, $\E[\sum_{y_i=0}B_i\ell_i]=\beta\sum_{y_i=0}\ell_i$.
\end{proof}

\noindent La proposición muestra que las tres técnicas persiguen el mismo cambio de pesos, y sus diferencias están en la varianza y en la información. La ponderación lo logra de forma determinista y sin descartar datos. El sobremuestreo añade aleatoriedad y repite ejemplos, lo que favorece el sobreajuste a esos puntos. El submuestreo descarta en promedio una fracción $1-\beta$ de la clase mayoritaria: con $0.1\%$ de fraudes y balance total ($\beta=n_1/n_0\approx0.001$) se descarta el $99.9\%$ de las transacciones legítimas, y la varianza de todo lo estimado sobre esa clase se multiplica por el orden de $1/\beta$. El submuestreo cambia además las probabilidades que aprende el modelo.

\begin{proposicion}[Corrección de probabilidades tras submuestreo]\label{prop:correccion}
\noindent Si la clase $0$ se submuestrea con probabilidad $\beta$ independiente de $x$, la probabilidad condicional en los datos submuestreados es $\eta_\beta(x)=\dfrac{\eta(x)}{\eta(x)+\beta(1-\eta(x))}$, y la original se recupera como $\eta(x)=\dfrac{\beta\,\eta_\beta(x)}{1-\eta_\beta(x)+\beta\,\eta_\beta(x)}$.
\end{proposicion}
\begin{proof}
\noindent Sea $S$ el indicador de que un registro se conserva, con $\Prob(S=1\mid Y=1,X=x)=1$ y $\Prob(S=1\mid Y=0,X=x)=\beta$. Por la regla de Bayes,
\[
\eta_\beta(x)=\Prob(Y=1\mid X=x,S=1)=\frac{1\cdot\eta(x)}{1\cdot\eta(x)+\beta(1-\eta(x))}.
\]
Despejando: $\eta_\beta(\eta+\beta-\beta\eta)=\eta$, es decir, $\eta_\beta\beta=\eta(1-\eta_\beta+\beta\eta_\beta)$, que da la segunda fórmula (el paréntesis es positivo: $1-\eta_\beta+\beta\eta_\beta=1-(1-\beta)\eta_\beta\ge\beta>0$ porque $\eta_\beta\le1$).
\end{proof}

\noindent Un modelo entrenado sobre datos submuestreados sobrestima la probabilidad de la clase rara; si se necesitan probabilidades calibradas (por ejemplo, para fijar un umbral de costo), hay que aplicar la corrección.

\subsection{SMOTE}\label{sec:smote}

\noindent \emph{Motivación.} El sobremuestreo aleatorio repite puntos exactos, y un modelo flexible puede aislarlos con regiones diminutas alrededor de cada uno. Chawla, Bowyer, Hall y Kegelmeyer (2002) propusieron generar puntos nuevos entre ejemplos minoritarios cercanos, para que la región de decisión de la clase rara se ensanche en lugar de fragmentarse.

\begin{definicion}[SMOTE]\label{def:smote}
\noindent Sea $M\subset\R^p$ el conjunto finito de ejemplos de la clase minoritaria en $\Dtr$, con $\abs M>k$. Los \emph{$k$ vecinos más cercanos} de $x\in M$ son los $k$ puntos de $M\setminus\{x\}$ con menor distancia euclídea a $x$ (con un criterio fijo de desempate). Para generar un punto sintético: se elige un ejemplo $x\in M$, se elige al azar uno $x'$ entre sus $k$ vecinos más cercanos, se toma $u\sim U(0,1)$ independiente y se define $z=x+u(x'-x)$.
\end{definicion}

\begin{proposicion}[Propiedades de SMOTE]\label{prop:smote}
\noindent (i) $z$ está en el segmento $[x,x']$ y, por tanto, en la envolvente convexa de los ejemplos minoritarios; SMOTE no puede generar puntos fuera de ella. (ii) Si, como modelo simplificado, $x$ y $x'$ fueran independientes con la misma distribución, de media $\mu$ y covarianza $\Sigma$, entonces $\E z=\mu$ y $\Cov z=\frac23\Sigma$.
\end{proposicion}
\begin{proof}
\noindent (i) $z=(1-u)x+ux'$ con $u\in[0,1]$ es combinación convexa de $x$ y $x'$, que pertenecen a la envolvente convexa, y esta es cerrada por combinaciones convexas. (ii) $z-\mu=(1-u)(x-\mu)+u(x'-\mu)$. Por independencia de $u$, $x$, $x'$, la esperanza es $\E[1-u]\cdot0+\E[u]\cdot0=0$, y la covarianza es $\E[(1-u)^2]\Sigma+\E[u^2]\Sigma+2\E[u(1-u)]\Cov(x,x')$, donde el último término es nulo por independencia. Como $\E[u^2]=\int_0^1t^2dt=\frac13$ y $\E[(1-u)^2]=\frac13$ por simetría, $\Cov z=\frac23\Sigma$.
\end{proof}

\noindent La parte (i) es a la vez la virtud y el límite de SMOTE: interpola, nunca extrapola. Si la clase minoritaria forma dos grupos separados por una región de la clase mayoritaria y $k$ es grande, algunos segmentos cruzan esa región y los puntos sintéticos caen donde no hay minoritarios reales, creando ruido de etiqueta. La parte (ii) indica que los puntos sintéticos están más concentrados que los reales: en ese modelo simplificado su dispersión es $2/3$ de la original (con vecinos cercanos la contracción es local, pero del mismo tipo). SMOTE solo tiene sentido en características numéricas continuas (interpolar entre dos categorías codificadas en one-hot da vectores que no corresponden a ninguna categoría) y, como cualquier remuestreo, se aplica solo a $\Dtr$ después de dividir: si se aplica antes, un punto sintético construido con un ejemplo de prueba termina en entrenamiento, que es fuga en el sentido de la Definición~\ref{def:fuga}.

\subsection{Aumento de datos}\label{sec:aumento}

\noindent \emph{Motivación.} Cuando recolectar más ejemplos de la clase minoritaria es caro, se pueden fabricar a partir de los existentes, siempre que se sepa que ciertas modificaciones no cambian la clase. Para que ``no cambiar la clase'' sea una condición verificable hay que formularla.

\begin{definicion}[Transformación que preserva la etiqueta]\label{def:preserva}
\noindent Supóngase que la etiqueta de cada entrada $x$ está determinada por una función $h^*$ (la etiqueta verdadera) definida en un conjunto $\mathcal{X}$ de entradas posibles. Una biyección $g:\mathcal X\to\mathcal X$ \emph{preserva la etiqueta} si $h^*(gx)=h^*(x)$ para todo $x\in\mathcal X$. El \emph{aumento de datos} con un conjunto finito $G$ de tales transformaciones (que contiene la identidad) reemplaza $\Dtr$ por los pares $(gx_i,y_i)$, $g\in G$.
\end{definicion}

\noindent Para imágenes de la Definición~\ref{def:imagen}, dos transformaciones típicas son el \emph{volteo horizontal} $F$ y la \emph{rotación de $180^\circ$} $R$:
\[
(FI)(u,v,c)=I(u,W-1-v,c),\qquad (RI)(u,v,c)=I(H-1-u,W-1-v,c).
\]
Ambas son permutaciones de los píxeles, y cumplen $F\circ F=R\circ R=\mathrm{id}$ porque $W-1-(W-1-v)=v$ y lo mismo en filas. Los \emph{ajustes de color} son transformaciones afines por canal, $I(\cdot,\cdot,c)\mapsto a_cI(\cdot,\cdot,c)+b_c$, del tipo de la Proposición~\ref{prop:norm-pix}(iii). Que una de ellas preserve la etiqueta depende del problema.

\begin{ejemplo}[Una rotación que no preserva la etiqueta]\label{ej:seis-nueve}
\noindent En reconocimiento de dígitos manuscritos, la rotación $R$ lleva la imagen de un ``$6$'' a una imagen que se lee como ``$9$'': $h^*(RI)=9\ne6=h^*(I)$. Aumentar con $R$ enseñaría al modelo etiquetas falsas. El volteo horizontal no preserva la lateralidad de una radiografía, y sí la clase ``gato'' en fotografías.
\end{ejemplo}

\begin{proposicion}[Riesgo aumentado]\label{prop:aumento}
\noindent El riesgo empírico sobre los datos aumentados es $\hat R_G(f)=\frac1{n\abs G}\sum_{i=1}^n\sum_{g\in G}\ell(f(gx_i),y_i)$. Si $f$ es invariante ($f(gx)=f(x)$ para todo $g\in G$ y todo $x$), entonces $\hat R_G(f)=\hat R(f)$. Si todo $g\in G$ preserva la etiqueta, cada par aumentado $(gx_i,y_i)$ cumple $y_i=h^*(gx_i)$ siempre que $y_i=h^*(x_i)$.
\end{proposicion}
\begin{proof}
\noindent La primera fórmula es la definición del riesgo empírico sobre los $n\abs G$ pares aumentados. Si $f$ es invariante, cada sumando interior es $\ell(f(x_i),y_i)$, y el promedio sobre $g$ de $\abs G$ términos iguales es ese término. La última afirmación es la Definición~\ref{def:preserva}: $h^*(gx_i)=h^*(x_i)=y_i$.
\end{proof}

\noindent La proposición muestra qué hace el aumento: no cambia el riesgo de los modelos ya invariantes y penaliza a los que no lo son, de modo que empuja el ajuste hacia la invariancia. Es la vía por los datos hacia la misma propiedad que el Teorema~\ref{teo:invariancia-cnn} obtiene por la arquitectura. Todo descansa en la hipótesis de preservación, que el Ejemplo~\ref{ej:seis-nueve} muestra que puede fallar. El aumento de texto (reemplazo por sinónimos, inserciones) es una función sobre sucesiones de tokens de la Definición~\ref{def:token} con la hipótesis análoga, más frágil: el Ejemplo~\ref{ej:negacion} muestra que cambiar o quitar un solo token puede invertir la etiqueta. Como indica la guía, el aumento parte de los datos existentes, mientras que los datos sintéticos se generan sin usar el conjunto original.

\subsection{Métricas de sesgo previas al entrenamiento}\label{sec:clarify}

\noindent SageMaker Clarify cuantifica, antes de entrenar, cuán representado está cada grupo. Sea una \emph{faceta} una característica que define dos grupos: $a$ (el que se sospecha favorecido) y $d$ (el desfavorecido), con $n_a$ y $n_d$ registros, de los cuales $n_a^{(1)}$ y $n_d^{(1)}$ tienen la etiqueta positiva.

\begin{definicion}[CI y DPL]\label{def:ci}
\noindent El \emph{desbalance de clases} de la faceta es $\mathrm{CI}=\dfrac{n_a-n_d}{n_a+n_d}$ y la \emph{diferencia en proporciones de etiquetas} es $\mathrm{DPL}=q_a-q_d$, con $q_a=n_a^{(1)}/n_a$ y $q_d=n_d^{(1)}/n_d$.
\end{definicion}

\begin{proposicion}[Rangos]\label{prop:ci}
\noindent $\mathrm{CI}\in[-1,1]$, con $\mathrm{CI}=0$ si y solo si $n_a=n_d$ y $\mathrm{CI}=\pm1$ si y solo si uno de los grupos está vacío. Para etiquetas binarias, $\mathrm{DPL}\in[-1,1]$, con $\mathrm{DPL}=0$ si y solo si ambos grupos tienen la misma proporción de positivos.
\end{proposicion}
\begin{proof}
\noindent Como $n_a,n_d\ge0$ y $n_a+n_d>0$, $\abs{n_a-n_d}\le n_a+n_d$, con igualdad si y solo si $n_a=0$ o $n_d=0$; el cociente vale $0$ si y solo si $n_a=n_d$. Para DPL, $q_a,q_d\in[0,1]$ porque $0\le n^{(1)}\le n$, y la diferencia de dos números de $[0,1]$ está en $[-1,1]$.
\end{proof}

\noindent En el ejemplo de la guía, con $10^6$ transacciones, $n_a=999\,000$ legítimas y $n_d=1000$ fraudulentas, $\mathrm{CI}=998\,000/10^6=0.998$. Como advierte el documento original, usar la etiqueta misma como faceta hace que CI mida solo el desbalance de la etiqueta; el uso previsto es una faceta demográfica, y entonces DPL es la diferencia de las tasas de resultado favorable entre grupos, el numerador del estadístico de la prueba de dos proporciones. Ambas métricas se calculan con conteos, sin necesidad del servicio.

\section{División de datos y fuga}\label{sec:division}

\noindent La Definición~\ref{def:fuga} fija el principio: toda transformación ajustada se ajusta con $\Dtr$ y se aplica sin cambios a validación y prueba. Esta sección estudia cómo construir la partición.

\medskip

\noindent \emph{Tamaño del conjunto de prueba.} Si la exactitud verdadera es $a$ y el conjunto de prueba tiene $n_{\mathrm{te}}$ registros independientes, la exactitud medida es una proporción binomial con desviación estándar $\sqrt{a(1-a)/n_{\mathrm{te}}}$ (la varianza de una media de $n_{\mathrm{te}}$ variables Bernoulli$(a)$). Con $a=0.9$ y $n_{\mathrm{te}}=1000$ es $0.0095$: diferencias de un punto porcentual entre modelos están en el nivel del ruido. Las proporciones de la guía ($60$--$80\%$ para entrenamiento, $10$--$20\%$ para validación y prueba) son un compromiso entre esta varianza y la cantidad de datos para aprender.

\begin{proposicion}[Clases raras en una división aleatoria]\label{prop:hiper}
\noindent Si un conjunto de $N$ registros con $K$ positivos se divide eligiendo al azar $n$ registros para prueba, el número de positivos en prueba es hipergeométrico, con media $nK/N$ y
\[
\Prob(\text{ningún positivo en prueba})=\frac{\binom{N-K}{n}}{\binom Nn}=\prod_{i=0}^{K-1}\frac{N-n-i}{N-i}\le\Big(1-\frac nN\Big)^K.
\]
\end{proposicion}
\begin{proof}
\noindent Todos los subconjuntos de tamaño $n$ son equiprobables, y los que no contienen positivos son los $\binom{N-K}n$ subconjuntos de los $N-K$ negativos. Escribiendo los binomiales con factoriales,
\[
\frac{\binom{N-K}n}{\binom Nn}=\frac{(N-K)!\,(N-n)!}{(N-K-n)!\,N!}=\frac{(N-n)(N-n-1)\cdots(N-n-K+1)}{N(N-1)\cdots(N-K+1)},
\]
ya que $(N-n)!/(N-n-K)!$ y $N!/(N-K)!$ son productos de $K$ factores consecutivos descendentes. Cada factor cumple $\frac{N-n-i}{N-i}=1-\frac n{N-i}\le1-\frac nN$. La media es $K$ veces la probabilidad $n/N$ de que un positivo dado caiga en prueba.
\end{proof}

\noindent Con $N=10\,000$, $K=10$ fraudes ($0.1\%$) y $n=1000$ ($10\%$ para prueba), el número esperado de fraudes en prueba es $1$ y la probabilidad de no tener ninguno es $\approx0.349$ (la cota da $0.9^{10}\approx0.349$). Una de cada tres divisiones aleatorias dejaría el conjunto de prueba sin la clase que se quiere detectar. La nota de precisión del documento original (``la división aleatoria no asegura proporciones'') queda así cuantificada.

\medskip

\noindent \emph{División estratificada.} Se divide cada clase por separado con la misma fracción $f$: la clase $c$ aporta $\lfloor fn_c\rfloor$ o $\lceil fn_c\rceil$ registros a prueba, de modo que la proporción de cada clase en prueba difiere de la global solo por redondeo (menos de un registro por clase). Elimina la variabilidad de la Proposición~\ref{prop:hiper}.

\medskip

\noindent \emph{División por clave.} Si varios registros comparten una clave (un cliente, un paciente), están correlacionados. Repartirlos entre entrenamiento y prueba produce el efecto de la Proposición~\ref{prop:duplicado} de forma atenuada: el modelo se evalúa en casi duplicados de lo que vio. Asignar claves completas a una sola partición mide el desempeño en claves nuevas, que es lo que ocurrirá en producción.

\medskip

\noindent \emph{División ordenada.} En series de tiempo, entrenar con el pasado y evaluar con el futuro es la única partición compatible con la Definición~\ref{def:admisible}; una partición aleatoria permite que el modelo interpole entre observaciones vecinas en el tiempo, lo que no podrá hacer al pronosticar.

\medskip

\noindent \emph{Orden de las operaciones.} Con lo anterior, el orden correcto es: deduplicar $\to$ dividir (estratificada, por clave u ordenada según el caso) $\to$ ajustar con $\Dtr$ imputación, límites de atípicos, transformaciones, escaladores, codificadores, selección y PCA $\to$ remuestrear o ponderar solo $\Dtr$ $\to$ aplicar las transformaciones ya ajustadas a validación y prueba. Un \texttt{Pipeline} de scikit-learn evaluado con validación cruzada ejecuta este orden dentro de cada pliegue.

\section{Servicios de AWS donde se aplican las técnicas}\label{sec:aws}

\noindent Esta sección sigue el reparto de servicios del documento original. Advertencias vigentes según ese documento (verificadas por él el 23 de septiembre de 2026): Amazon SageMaker se llama ahora Amazon SageMaker AI; SageMaker Data Wrangler está integrado en SageMaker Canvas; SageMaker Ground Truth y SageMaker Clarify ya no admiten clientes nuevos (los existentes pueden seguir usándolos); Amazon Mechanical Turk cierra el 30 de septiembre de 2026. Los nombres de transformaciones son los de la documentación de AWS y conviene confirmarlos en la versión vigente.

\medskip

\noindent \emph{Herramientas transversales.} \textbf{SageMaker Data Wrangler} (dentro de Canvas) es la herramienta visual de bajo código para casi todas las técnicas de este texto y exporta el flujo a S3, a un pipeline de SageMaker, a Feature Store o a un script. \textbf{AWS Glue DataBrew} es la alternativa visual sin código de la familia Glue, con recetas reutilizables que se ejecutan como trabajos serverless. Los \textbf{notebooks de SageMaker} permiten usar scikit-learn, pandas y \texttt{category\_encoders} con el código de la guía (\texttt{SimpleImputer}, \texttt{StandardScaler}, \texttt{PowerTransformer}, \texttt{OneHotEncoder}, \texttt{FeatureHasher}), y los \textbf{trabajos de AWS Glue} (PySpark) ejecutan la misma lógica de forma repetible sobre el volumen completo.

\medskip

\begin{center}
\small
\begin{tabularx}{\linewidth}{@{}>{\raggedright\arraybackslash}p{3.4cm}Y@{}}
\toprule
Técnica (sección) & Servicios de AWS\\
\midrule
Valores faltantes (\ref{sec:faltantes}) & Data Wrangler: \emph{Handle missing} (\emph{Impute missing} con media, mediana o moda; \emph{Fill missing}; \emph{Add indicator for missing}; \emph{Drop missing}). Glue DataBrew: transformaciones integradas de relleno. Notebooks: \texttt{SimpleImputer}.\\
Atípicos (\ref{sec:atipicos}) & Data Wrangler: \emph{Handle outliers} (desviación estándar, desviación estándar robusta, cuantiles, mínimo--máximo; recortar, eliminar o invalidar). Glue DataBrew: marcar, eliminar, reemplazar o reescalar atípicos por puntuación Z o IQR. Notebooks: pandas y SciPy.\\
Deduplicación (\ref{sec:dedup}) & Glue DataBrew (opción que la guía recomienda por simplicidad).\\
Selección y PCA (\ref{sec:seleccion}) & Data Wrangler: \emph{Dimensionality Reduction} (PCA por número de componentes o umbral de varianza, con opción de centrar y escalar). Algoritmo integrado PCA de SageMaker AI. Notebooks: información mutua y $\chi^2$ de scikit-learn.\\
Escalado (\ref{sec:escalado}) & Data Wrangler: \emph{Process numeric} (\emph{Standard Scaler}, \emph{Robust Scaler}, \emph{Min Max Scaler}, \emph{Max Absolute Scaler}). Glue DataBrew: \texttt{SCALE} (\texttt{MIN\_MAX}, \texttt{SCALE\_BETWEEN}, \texttt{MEAN\_NORMALIZATION}, \texttt{Z\_SCORE}). Notebooks y Glue jobs.\\
Transformaciones no lineales (\ref{sec:nolineales}) & Notebooks o trabajos de procesamiento de SageMaker: \texttt{PowerTransformer} (Box--Cox, Yeo--Johnson), \texttt{np.log1p}. Data Wrangler: \emph{Custom formula} con funciones de Spark SQL. Glue DataBrew: transformaciones de sesgo.\\
Discretización (\ref{sec:binning}) & Glue DataBrew: \texttt{BUCKETIZATION}. Notebooks: \texttt{KBinsDiscretizer}.\\
Categóricas (\ref{sec:categoricas}) & Data Wrangler: \emph{Encode categorical} (\emph{Ordinal encode}, \emph{One-hot encode}, \emph{Similarity encode}). Glue DataBrew: \texttt{ONE\_HOT\_ENCODING}, \texttt{CATEGORICAL\_MAPPING}. Notebooks: \texttt{OrdinalEncoder}, \texttt{OneHotEncoder}, \texttt{BinaryEncoder}, \texttt{FeatureHasher}.\\
Series de tiempo (\ref{sec:series}) & Data Wrangler, transformaciones de series de tiempo: \emph{Featurize datetime} (modo ordinal o cíclico), \emph{One-hot encode}, \emph{Lag features}, \emph{Rolling window features}, \emph{Extract features}, además de remuestreo e imputación temporal. Feature Store: consultas \emph{point-in-time}.\\
Imágenes (\ref{sec:imagenes}) & SageMaker JumpStart: modelos de visión preentrenados (CNN) para extraer embeddings. Amazon Rekognition: etiquetas, objetos, texto, rostros y cuadros delimitadores por API. Data Wrangler: normalización y PCA. Feature Store: almacenamiento y reutilización.\\
Texto (\ref{sec:texto}) & Amazon Comprehend: entidades, frases clave, sentimiento, idioma. Amazon Textract: OCR, formularios y tablas. Algoritmos integrados de SageMaker AI: BlazingText (Word2Vec) y LDA (tópicos). Data Wrangler: \emph{Featurize text} (\emph{Character statistics}, \emph{Vectorize} con TF-IDF). Glue DataBrew: \texttt{TOKENIZATION}. Amazon Bedrock: modelos fundacionales con tokenización por subpalabras y cobro por tokens.\\
Etiquetado & SageMaker Ground Truth (sin clientes nuevos).\\
Desbalance (\ref{sec:desbalance}) & Data Wrangler: \emph{Balance data} (\emph{Random oversampling}, \emph{Random undersampling}, SMOTE). Ponderación dentro de los algoritmos (\texttt{scale\_pos\_weight} en XGBoost). SageMaker Clarify: métricas CI y DPL previas al entrenamiento (sin clientes nuevos; se calculan igual con pandas).\\
División (\ref{sec:division}) & Data Wrangler: \emph{Split data} (\emph{Randomized}, \emph{Ordered}, \emph{Stratified}, \emph{Split by key}), con semilla reproducible.\\
Almacenar y servir características & SageMaker Feature Store: almacén \emph{online} para inferencia en tiempo real y \emph{offline} en S3 para entrenamiento, con consultas \emph{point-in-time} que evitan la fuga (Definición~\ref{def:admisible}) y el \emph{training-serving skew}.\\
Datos de origen & Amazon S3 como lago de datos y fuente única de verdad; AWS Glue Data Catalog y \emph{crawlers}; Amazon Athena para consultas; AWS Lake Formation para permisos por columna, fila y celda.\\
\bottomrule
\end{tabularx}
\end{center}

\noindent El hilo común con el desarrollo matemático es la Definición~\ref{def:ajustada}: cada herramienta separa el ajuste de la aplicación de las transformaciones. Data Wrangler y DataBrew lo hacen al guardar el flujo o la receta y reaplicarla a datos nuevos; Feature Store lo hace al calcular cada característica una sola vez y servir el mismo valor al entrenamiento y a la inferencia. Es esa separación la que impide la fuga de datos y el \emph{training-serving skew}.

\begin{thebibliography}{99}
\bibitem{boxcox} G.~E.~P.~Box y D.~R.~Cox, An analysis of transformations, \emph{J. Roy. Statist. Soc. B} \textbf{26} (1964), 211--252.
\bibitem{chawla} N.~V.~Chawla, K.~W.~Bowyer, L.~O.~Hall y W.~P.~Kegelmeyer, SMOTE: Synthetic minority over-sampling technique, \emph{J. Artificial Intelligence Res.} \textbf{16} (2002), 321--357. arXiv:1106.1813.
\bibitem{kuhn} M.~Kuhn y K.~Johnson, \emph{Feature Engineering and Selection: A Practical Approach for Predictive Models}, CRC Press, 2019. Versión en línea: \url{https://feat.engineering/}.
\bibitem{levy} O.~Levy e Y.~Goldberg, Neural word embedding as implicit matrix factorization, \emph{Advances in Neural Information Processing Systems} \textbf{27} (2014).
\bibitem{mikolov} T.~Mikolov, I.~Sutskever, K.~Chen, G.~Corrado y J.~Dean, Distributed representations of words and phrases and their compositionality, \emph{Advances in Neural Information Processing Systems} \textbf{26} (2013).
\bibitem{glove} J.~Pennington, R.~Socher y C.~D.~Manning, GloVe: Global vectors for word representation, \emph{Proc. EMNLP} (2014).
\bibitem{rubin} D.~B.~Rubin, Inference and missing data, \emph{Biometrika} \textbf{63} (1976), 581--592.
\bibitem{sennrich} R.~Sennrich, B.~Haddow y A.~Birch, Neural machine translation of rare words with subword units, \emph{Proc. ACL} (2016).
\bibitem{shiffler} R.~E.~Shiffler, Maximum Z scores and outliers, \emph{Amer. Statist.} \textbf{42} (1988), 79--80.
\bibitem{tukey} J.~W.~Tukey, \emph{Exploratory Data Analysis}, Addison-Wesley, 1977.
\bibitem{weinberger} K.~Weinberger, A.~Dasgupta, J.~Langford, A.~Smola y J.~Attenberg, Feature hashing for large scale multitask learning, \emph{Proc. ICML} (2009). arXiv:0902.2206.
\bibitem{yeo} I.-K.~Yeo y R.~A.~Johnson, A new family of power transformations to improve normality or symmetry, \emph{Biometrika} \textbf{87} (2000), 954--959.
\bibitem{aws} Amazon Web Services, \emph{Amazon SageMaker AI Developer Guide} (transformaciones de Data Wrangler en Canvas, BlazingText) y \emph{AWS Glue DataBrew Developer Guide} (acciones de receta), consultadas en septiembre de 2026.
\end{thebibliography}

\end{document}
